%~Mouliné par MaN_auto v.0.40.4 (bbb487d5) 2026-04-16 14:48:05
\documentclass[AHL,Unicode,longabstracts,published]{cedram}

\usepackage{dsfont}

\DeclareMathOperator{\diag}{diag}
\DeclareMathOperator{\dist}{dist}
\DeclareMathOperator{\re}{Re}
\DeclareMathOperator{\im}{Im}
\DeclareMathOperator{\voll}{vol}
\DeclareMathOperator{\rank}{rank}
\DeclareMathOperator{\Neigh}{Neigh}
\DeclareMathOperator{\divv}{div}
\DeclareMathOperator{\Graph}{Graph}
\DeclareMathOperator{\supp}{supp}
\DeclareMathOperator{\itt}{it}
\DeclareMathOperator{\SPS}{SPS}%statuer en relecture pour l'instant $\mathit{SPS}
\DeclareMathOperator{\SQS}{SQS}%statuer en relecture pour l'instant $\mathit{SQS}
\DeclareMathOperator{\idd}{id}
\DeclareMathOperator{\irr}{ir}

\newcommand{\Opp}{\mathrm{Op}}

\newcommand{\iTO}{\mathit{TO}}

\newcommand{\clK}{\mathcal{K}}
\newcommand{\clD}{\mathcal{D}}
\newcommand{\clS}{\mathcal{S}}
\newcommand{\clJ}{\mathcal{J}}
\newcommand{\clE}{\mathcal{E}}
\newcommand{\clO}{\mathcal{O}}
\newcommand{\clF}{\mathcal{F}}
\newcommand{\clH}{\mathcal{H}}

\newcommand\1{\mathds{1}}


\newcommand{\R}{\mathbb{R}}
\newcommand{\C}{\mathbb{C}}
\newcommand{\Z}{\mathbb{Z}}
\newcommand{\N}{\mathbb{N}}

\newcommand{\dd}{\mathrm{d}}
\renewcommand{\S}{\mathbb{S}}
%%%%%%%%%%%%%%%%%%%%%%%%%%%%%%%%%%%%%%%%%%%%%%%%%%%%%%%%%%%%%%%%%%%%%%%%%%
%Better widebar

\makeatletter
\let\save@mathaccent\mathaccent
\newcommand*\if@single[3]{%
  \setbox0\hbox{${\mathaccent"0362{#1}}^H$}%
  \setbox2\hbox{${\mathaccent"0362{\kern0pt#1}}^H$}%
  \ifdim\ht0=\ht2 #3\else #2\fi
  }
%The bar will be moved to the right by a half of \macc@kerna, which is computed by amsmath:
\newcommand*\rel@kern[1]{\kern#1\dimexpr\macc@kerna}
%If there's a superscript following the bar, then no negative kern may follow the bar;
%an additional {} makes sure that the superscript is high enough in this case:
\newcommand*\widebar[1]{\@ifnextchar^{{\wide@bar{#1}{0}}}{\wide@bar{#1}{1}}}
%Use a separate algorithm for single symbols:
\newcommand*\wide@bar[2]{\if@single{#1}{\wide@bar@{#1}{#2}{1}}{\wide@bar@{#1}{#2}{2}}}
\newcommand*\wide@bar@[3]{%
  \begingroup
  \def\mathaccent##1##2{%
%Enable nesting of accents:
    \let\mathaccent\save@mathaccent
%If there's more than a single symbol, use the first character instead (see below):
    \if#32 \let\macc@nucleus\first@char \fi
%Determine the italic correction:
    \setbox\z@\hbox{$\macc@style{\macc@nucleus}_{}$}%
    \setbox\tw@\hbox{$\macc@style{\macc@nucleus}{}_{}$}%
    \dimen@\wd\tw@
    \advance\dimen@-\wd\z@
%Now \dimen@ is the italic correction of the symbol.
    \divide\dimen@ 3
    \@tempdima\wd\tw@
    \advance\@tempdima-\scriptspace
%Now \@tempdima is the width of the symbol.
    \divide\@tempdima 10
    \advance\dimen@-\@tempdima
%Now \dimen@ = (italic correction / 3) - (Breite / 10)
    \ifdim\dimen@>\z@ \dimen@0pt\fi
%The bar will be shortened in the case \dimen@<0 !
    \rel@kern{0.6}\kern-\dimen@
    \if#31
      \overline{\rel@kern{-0.6}\kern\dimen@\macc@nucleus\rel@kern{0.4}\kern\dimen@}%
      \advance\dimen@0.4\dimexpr\macc@kerna
%Place the combined final kern (-\dimen@) if it is >0 or if a superscript follows:
      \let\final@kern#2%
      \ifdim\dimen@<\z@ \let\final@kern1\fi
      \if\final@kern1 \kern-\dimen@\fi
    \else
      \overline{\rel@kern{-0.6}\kern\dimen@#1}%
    \fi
  }%
  \macc@depth\@ne
  \let\math@bgroup\@empty \let\math@egroup\macc@set@skewchar
  \mathsurround\z@ \frozen@everymath{\mathgroup\macc@group\relax}%
  \macc@set@skewchar\relax
  \let\mathaccentV\macc@nested@a
%The following initialises \macc@kerna and calls \mathaccent:
  \if#31
    \macc@nested@a\relax111{#1}%
  \else
%If the argument consists of more than one symbol, and if the first token is
%a letter, use that letter for the computations:
    \def\gobble@till@marker##1\endmarker{}%
    \futurelet\first@char\gobble@till@marker#1\endmarker
    \ifcat\noexpand\first@char A\else
      \def\first@char{}%
    \fi
    \macc@nested@a\relax111{\first@char}%
  \fi
  \endgroup
}
\makeatother

\let\oldbar\bar
\renewcommand*{\bar}[1]{{\mathchoice{\widebar{#1}}{\widebar{#1}}{\widebar{#1}}{\oldbar{#1}}}}
%%%%%%%%%%%%%%%%%%%%%%%%%%%%%%%%%%%%%%%%%%%%%%%%%%%%%%%%%%%%%%%%%%%%%%%%%%%%%%%%%%%%%%%%%%%%%%%%%%%%%%%%%%%
\graphicspath{{./figures/}}

\newcommand*{\mk}{\mkern -1mu}
\newcommand*{\Mk}{\mkern -2mu}
\newcommand*{\mK}{\mkern 1mu}
\newcommand*{\MK}{\mkern 2mu}

\hypersetup{urlcolor=purple, linkcolor=blue, citecolor=red}

\newcommand*{\relabel}{\renewcommand{\labelenumi}{(\theenumi)}}
\newcommand*{\romanenumi}{\renewcommand*{\theenumi}{\roman{enumi}}\relabel}
\newcommand*{\Romanenumi}{\renewcommand*{\theenumi}{\Roman{enumi}}\relabel}
\newcommand*{\alphenumi}{\renewcommand*{\theenumi}{\alph{enumi}}\relabel}
\newcommand*{\Alphenumi}{\renewcommand*{\theenumi}{\Alph{enumi}}\relabel}
\let\oldtilde\tilde
\renewcommand*{\tilde}[1]{\mathchoice{\widetilde{#1}}{\widetilde{#1}}{\oldtilde{#1}}{\oldtilde{#1}}}
\let\oldhat\hat
\renewcommand*{\hat}[1]{\mathchoice{\widehat{#1}}{\widehat{#1}}{\oldhat{#1}}{\oldhat{#1}}}
\let\oldexists\exists
\renewcommand*{\exists}{\mathrel{\oldexists}}
\let\oldforall\forall
\renewcommand*{\forall}{\mathrel{\oldforall}}

\newcommand{\llbracket}{\mathopen{[\mkern -2.7mu[}}
\newcommand{\rrbracket}{\mathclose{]\mkern -2.7mu]}}
\newcommand{\llb}{\llbracket}
\newcommand{\rrb}{\rrbracket}
%%%%%%%%%%%%%%%%%%%%%%%%%%%%%%%%%%%%%%%%%%%%%%%%%%%%%%%%%%%%%%%%%%%%%%%%%%%%%%%%%%%%%%%%%%
\renewcommand*{\to}{\mathchoice{\longrightarrow}{\rightarrow}{\rightarrow}{\rightarrow}}


\let\oldmapsto\mapsto
\renewcommand*{\mapsto}{\mathchoice{\longmapsto}{\oldmapsto}{\oldmapsto}{\oldmapsto}}

\renewcommand*{\implies}{\mathchoice{\Longrightarrow}{\Rightarrow}{\Rightarrow}{\Rightarrow}}

%~ This fixes spacing issues around \left( and \right) and other similar symbols
\let\originalleft\left 
\let\originalright\right
\renewcommand{\left}{\mathopen{}\mathclose\bgroup\originalleft}
\renewcommand{\right}{\aftergroup\egroup\originalright}
%%%%%%%%%%%%%%%%%%%%%%%%%%%%%%%%%%%%%%%%%%%%%%%%%%%%%%%%%%%%%%%%%%%%%%%%%%%%%%%%%%%%%%%%%%%

\title[Analytic Szeg\H{o} kernel and Fourier integral operators]{The Szeg\H{o} kernel in analytic regularity and analytic Fourier integral operators}

\alttitle{Noyau de Szeg\H{o} en régularité analytique et opérateurs de Fourier intégraux}

\subjclass{32A25, 32W10, 32T15, 35A18, 35A20, 35A27, 35S30, 58J15, 58J40}


\keywords{Szeg\H{o} kernel, Toeplitz operators, analytic microlocal analysis}



\author[\initial{A.} \lastname{Deleporte}]{\firstname{Alix} \lastname{Deleporte}}
\address{Université Paris-Saclay, CNRS\\
Laboratoire de Mathématiques d'Orsay\\
91405 Orsay (France)}
\email{alix.deleporte@universite-paris-saclay.fr}
\thanks{The author acknowledges support by the  ANR-24-CE40-5905-01 ``STENTOR'' project.}
\CDRGrant[STENTOR]{ANR-24-CE40-5905-01}

\begin{abstract}
We build a general theory of microlocal (homogeneous) Fourier integral operators in real-analytic
regularity, following the general construction in the smooth case by H\"ormander and Duistermaat. In particular, we prove that the Boutet--Sj\"ostrand parametrix for the Szeg\H{o} projector at the boundary of a strongly pseudo-convex real-analytic domain can be realised by an analytic Fourier integral operator. We then study some applications, such as FBI-type transforms on compact, real-analytic Riemannian manifolds and propagators of one-homogeneous (pseudo)differential operators.
\end{abstract}

\begin{altabstract}
Nous construisons une théorie générale des
opérateurs de Fouriers intégraux microlocaux (homogènes) en régularité analytique, en suivant la construction générale due, dans le cas lisse, à H\"ormander et Duistermaat. En particulier, nous montrons que la paramétrice de Boutet--Sj\"ostrand pour le projecteur de Szeg\H{o} sur le bord d'un domaine analytique fortement pseudoconvexe peut être réalisée comme un opérateur de Fourier intégral analytique. Nous étudions ensuite quelques applications, comme les transformées de type FBI sur les variétés riemanniennes analytiques compactes, et les propagateurs d'opérateurs (pseudo)différentiels d'ordre un.
\end{altabstract}

\datereceived{2023-06-28}
\daterevised{2026-02-10}
\dateaccepted{2026-02-19}

\editors{S.~V{\={u}}~Ng{\d{o}}c and P.~Gérard}
%%%%%%%%%%%%%%%%%%%%%%%%%%%%%%%%%%%%%%%%%%%%%%%%%%%%%%%%%%%%%%%%%%%%%%%%%%%%%%%%%%%%%%%%%%%%%%%%%%%%%%%%%%%%%%%%%%%%%%%%%%%%%%%%%%%%%%%%%%%%%%%%%%%%%%%%%%%%%%%%%%%%%%%%%%%%%%%%%%%%%
\dateposted{2026-06-12}
\begin{document}
\maketitle

\section{Introduction}
\label{sec:introduction}
This project began as an investigation of the following problem: let $\Omega\subset \C^n$ be an open set with real-analytic boundary; can we characterise the analytic singularities of the functions on the boundary $\partial \Omega$ which are boundary values of holomorphic functions inside $\Omega$?

The problem is well-posed when $\Omega$ is \emph{strongly pseudoconvex} (see Definition~\ref{def:pseudoconvex}) and a precise description modulo smooth functions of the Szeg\H{o} projector $S$ (from $L^2(\partial \Omega)$ onto the subspace of boundary values of holomorphic functions) is well-known in the case where $\partial \Omega$ is smooth~\cite{boutet_de_monvel_sur_1975,fefferman_bergman_1976}. It is then natural to expect that, if $\partial \Omega$ is analytic, then an even more precise description of $S$ can be reached. This microlocal problem admits a semiclassical (small parameter) counterpart, about the description of spaces of holomorphic sections of high powers of ample line bundles over K\"ahler manifolds. In the latter case, a description of the Szeg\H{o} kernel in analytic regularity was recently obtained~\cite{charles_analytic_2021,deleporte_toeplitz_2018,deleporte_direct_2024,hezari_off-diagonal_2017,hezari_property_2021,rouby_analytic_2018}.

It quickly became obvious to the author that a natural course of action to study the analytic singularities of $S$ was to adapt the method of~\cite{boutet_de_monvel_sur_1975}, where $S$ is described as a {Fourier integral operator} with complex-valued phase function, and where the proof uses a microlocal normal form procedure (and in particular, a conjugation by a {Fourier integral operator}). Building a theory of analytic {Fourier integral operators} that lies as close as possible to the smooth constructions~\cite{hormander_analysis_1985} became our next goal and constitutes the bulk of this article.

In the microlocal as well as semiclassical case, the study of the Szeg\H{o} projector has important applications in geometric quantization and ``general purpose'' microlocal/semiclassical analysis. One particular instance of strongly pseudoconvex open set with real-analytic boundary is the Grauert tube around a real-analytic Riemannian manifold; the range of $S$ is then the target space of a natural FBI-type transformation. Thus the analytic counterpart of the results of~\cite{boutet_de_monvel_sur_1975} serves as a tool for other problems related to the spectrum and dynamics of differential and pseudo-differential operators. These techniques could be particularly relevant to the study of non-self-adjoint analytic (pseudo)differential operators.

\subsection{What this article contains}\label{sec:what-this-article}

Section~\ref{sec:bound-cauchy-riem} reviews some basic properties of the boundary Cauchy--Riemann problem for pseudoconvex open sets. The impatient reader only interested in ``off-the-shelf'' properties of analytic Fourier integral operators may skip it, however the pseudoconvex open set $\R^n\times S^{n-1}$, on which the Szeg\H{o} projector admits an exact formula, is useful in the construction of the general theory.

Section~\ref{sec:analytic-symbols} contains the ``formal analytic'' content of our construction. We describe spaces of analytic amplitudes (formal or not) in Section~\ref{sec:ampl-analyt-micr}, then we study the geometric properties of homogeneous phase functions and their associated conical Lagrangians in Section~\ref{sec:posit-nond-phas}. We also prove a stationary phase type theorem in Section~\ref{sec:stationary-phase}. Most of the contents of this section will look familiar to experts, but many details (in particular, the sense in which we perform Borel summations of formal analytic amplitudes in the microlocal setting) turn out to be non-trivial and crucial.

The core of the construction of analytic Fourier integral operators is Section~\ref{sec:four-integr-oper}. We properly define analytic {Fourier integral operators} in Section~\ref{sec:first-properties}, in an ascending level of geometric generality, and study the composition rule and the action of stationary phase; our results are gathered in Theorem~\ref{prop:FIOs}. The algebra of analytic pseudodifferential operators, which of course is of utmost importance, is described in Section~\ref{sec:pseud-oper-1}; in our framework, to microlocalize the action of these operators\linebreak we introduce a FBI transform, which conjugates the problem to operators on $\R^n\times S^{n-1}$ whose microlocal structure resembles that of the Szeg\H{o} projector. As in~\cite{boutet_de_monvel_spectral_1981}, we call them ``Toeplitz operators'', even though they are not best described as $SQS$ for~$Q$ a pseudodifferential operators, but rather by a {Fourier integral operator} formula, akin to the ``covariant Toeplitz operators'' of~\cite{charles_berezin-toeplitz_2003}. These Toeplitz operators turn out to be a very efficient microlocal replacement of pseudodifferential operators, allowing us to describe refined properties of Fourier integral operators (notably inversion) microlocally in Section~\ref{sec:advanced-properties}. Besides pseudodifferential operators and Toeplitz operators, important examples of {Fourier integral operators} are ``quantized contact transformations'': one can quantize the action of a one-homogeneous symplectic change of variables into a unitary operator, and we do so in Proposition~\ref{prop:quantized-contact}.

We then turn our attention to the specific problem of the Szeg\H{o} projector for a general pseudoconvex open set in Section~\ref{sec:szegho-kern-toepl}. We prove a normal form theorem for the~$\bar{\partial}_b$ operator in analytic regularity, and use it to describe the Szeg\H{o} projector as an analytic Fourier integral operator in Theorem~\ref{prop:analytic-Szego}. We then generalise the properties of Toeplitz operators obtained in Section~\ref{sec:pseud-oper-1} to this more general geometric setting.

Some illustrations of our techniques are presented in Section~\ref{sec:few-applications}. We discuss in particular the case of Grauert tubes, and how a natural FBI-type transformation on compact analytic Riemannian manifolds helps to translate ``usual'' questions concerning differential or pseudodifferential operators to the Toeplitz framework, where they are more easily studied. We prove in Proposition~\ref{prop:quantum_propag} that propagators of one-homogeneous, self-adjoint analytic pseudodifferential operators are analytic {Fourier integral operators}. We also discuss briefly how to obtain semiclassical results from our microlocal toolbox. To our knowledge, there is no general theory of analytic {Fourier integral operators} available in the semiclassical case (and certainly not one which preserves analytic function spaces), and we hope that our results can be of use in this direction.

The Appendix~\ref{sec:algebra-symb-pseud}--\ref{sec:funct-spac-analyt} presents some well-known facts about analytic semiclassical analysis, which we gather for the comfort of the reader. We discuss in particular the properties of spaces of analytic functions, their dual spaces (called \emph{analytic functi\-onals}), and a convenient quotient space of analytic functionals called hyperfunctions.

\subsection{Comparison with earlier work}\label{sec:comp-with-earl}

Microlocal analysis in real-analytic regularity has a long history, but its early achievements have been shadowed by the success and popularity of the $C^{\infty}$ techniques.

Building on an observation that in some
\emph{formal} sense, there are classes of pseudodifferential operators adapted to real-analytic regularity~\cite{boutet_de_monvel_pseudo-differential_1967}, a general framework was proposed~\cite{sato_microfunctions_1972} but ultimately, this theory never reached the same foundational status as the $C^{\infty}$ counterpart. In particular, a method of proof of Theorem~\ref{prop:analytic-Szego} using the tools of~\cite{sato_microfunctions_1972} appears in~\cite{kashiwara_analyse_1976}, but the status of this proof and of the statement itself is disputed; we mention~\cite{kaneko_introduction_1989} for a more detailed, but incomplete, discussion of Kashiwara's approach. One of the initial goals of this work was the hope to settle once and for all the status of Theorem~\ref{prop:analytic-Szego}.

In contrast, the theory of \emph{semiclassical} analysis in real-analytic regularity is now well-developed, after the seminal work~\cite{sjostrand_singularites_1982}. In this setting, one is interested in improving remainders from $O(\hbar^{\infty})$ to $O(e^{-c\hbar^{-1}})$ for some $c>0$, rather than improving remainders from smooth to real-analytic. The semiclassical theory makes full use of complex-valued phase functions, whose manipulation is perhaps even easier than in the $C^{\infty}$ case because of the naturalness of extensions into the complex.

It is often said that microlocal analysis is a particular case of semiclassical analysis. We found out, however, that the usual semiclassical constructions of pseudodifferential operators, other Fourier integral operators, or even Borel summations of analytic symbols, was not applicable here. This is chiefly due to the fact that, when presented with an oscillating integral with small parameter, of the form $\int e^{i\frac{\phi(x)}{\hbar}}a(x)\dd x$, we may remove the areas where $\im(\phi)>0$ by a simple (smooth) cut-off argument, leading to $O(e^{-c\hbar^{-1}})$ errors. As a consequence, at its present stage of development, analytic semiclassical analysis deals with errors of size $O(e^{-c\hbar^{-1}})$ in $C^{k}$ topology for fixed~$k$. In the microlocal setting, we cannot afford to introduce cut-offs even in areas of phase space which we presume are away from the analytic wave front set. Even apart from stationary phase, cut-and-paste arguments, used to pass from local definitions of objects in coordinate charts to global descriptions, must be replaced by slightly more refined arguments of cohomological nature: broadly speaking, we are able to define how objects act locally and then we can use somewhat abstract gluing tools. This is a standard approach in the study of holomorphic or real-analytic functions, but it turns out (see the Appendix~\ref{sec:algebra-symb-pseud}--\ref{sec:funct-spac-analyt}) that these arguments already work in the $C^{\infty}$ setting.

Conversely, at least the elementary results of analytic semiclassical analysis can be obtained from our microlocal setting by considering Fourier modes in an auxiliary variable; we describe this briefly in Section~\ref{sec:from-micr-semicl}.

The book~\cite{treves_analytic_2022} presents an overview of the existing literature on general-purpose analytic microlocal analysis and some of its applications. Sections~\ref{sec:analytic-symbols}--\ref{sec:four-integr-oper} have been somewhat motivated by this book, as we found out that adapting the homogeneous FIO framework to the real-analytic case is not as straightforward as it seems and many statements and constructions of this book deserved a more thorough approach.

\subsection*{Acknowledgments}

We are most grateful to Michael Hitrik and Johannes Sj\"ostrand for their encouragement and useful remarks throughout the realisation of this project.
This work is dedicated to the memory of Steve Zelditch, who motivated us to give a complete proof of the analytic parametrix of the Szeg\H{o} kernel in the first place.

\section{The boundary Cauchy--Riemann problem}\label{sec:bound-cauchy-riem}

This section is a primer on the theory of boundary Cauchy--Riemann problems; we present some facts which, besides being the basis for Theorem~\ref{prop:analytic-Szego}, will be useful for the general construction of Fourier integral operators in real-analytic regularity.

\subsection{Strongly pseudoconvex domains and CR geometry}\label{sec:strongly-pseud-doma}

\begin{defi}\label{def:pseudoconvex}
Let $Y$ be a (paracompact, boundaryless) complex manifold of
dimension $n\geq 2$.
\begin{itemize}
\item Let $U\subset Y$ open. A function $\rho\in C^2(U,\R)$ is \emph{strongly plurisubharmonic} when at every point of $U$, in local coordinates, the Hermitian matrix $[\frac{\partial^2\rho}{\partial z_j\partial
\bar{z_k}}]_{j,k}$ is positive definite. (This does not depend on the coordinates).
\item A compact real hypersurface $X\subset Y$ is \emph{strongly pseudoconvex} when there\linebreak exists an open neighbourhood $U$ of $X$ and a strongly plurisubharmonic function $\rho:U\to R$, such that $\dd \rho$ never vanishes and $X=\{\rho=0\}$.
\end{itemize}
\end{defi}

The situation of interest is the case where $X$ separates $Y$ into two connected components; the ``inside'' component $\Omega$ (the one which contains $\{\rho<0\}$) is then called a \emph{strongly pseudoconvex domain}, independently on the behaviour of $\rho$ far inside~$\Omega$.

Given a general open set $\Omega\subset Y$ with regular boundary, one may ask how holomorphic functions on $\Omega$ behave near $\partial
\Omega$: whether they can be extended (using, for instance, the Hartogs theorem) or whether their restrictions to $\partial \Omega$ admit a simple description. As it turns out, these two questions are related and strongly pseudoconvex domains are exactly those for which restrictions of holomorphic functions are solutions of a hypoelliptic linear PDE, the boundary Cauchy--Riemann operator $\bar{\partial}_b$~\cite{kohn_harmonic_1963}. The Szeg\H{o} projector $S:L^2(X)\to \ker_{L^2}(\bar{\partial}_b)$ allows to describe the space of solutions, as well as (after application of the Poisson kernel) the Bergman space of holomorphic functions on $\Omega$. The description of $S$ modulo smooth functions~\cite{boutet_de_monvel_sur_1975,fefferman_bergman_1976} when $X$ is smooth has numerous crucial applications in complex and K\"ahler geometry as well as geometric quantization. For example, the local definition of a strongly pseudoconvex hypersurface above, or even an intrinsic definition mimicking the properties of $X$ gathered in the next Proposition~\ref{prop:CR-structure}, are equivalent to the realisation of $X$ as the boundary of a strongly pseudoconvex open set in $\C^n$, with a globally defined and strongly plurisubharmonic $\rho$, as soon as $n\geq 3$~\cite{boutet_de_monvel_integration_1974}.

\pagebreak{}
\begin{prop}\label{prop:CR-structure}
A strongly pseudoconvex hypersurface $X\subset Y$ inherits a CR-structure as follows: $TX\otimes \C$ decomposes as a direct sum of three involutive components $T^{(0,1)}X\oplus T^{(1,0)}X\oplus V$, where:
\begin{itemize}
\item $T^{(1,0)}X=T^{(1,0)}Y\cap (TX\otimes \C)$ is the $n-1$-dimensional bundle of holomorphic tangent vectors.
\item $T^{(0,1)}X=T^{(0,1)}Y\cap (TX\otimes \C)$ is the $n-1$-dimensional bundle of anti-holomorphic tangent vectors.
\item $V$ is a line bundle over $X$.
\end{itemize}
The line bundle $\Sigma=(T^{(0,1)}X\oplus T^{(1,0)}X)^*\subset T^*X\otimes
\C$ is spanned by the real one-form $\alpha:(x,\xi)\mapsto \dd
\rho(x)(J\xi)$ whenever $\rho$ is a pseudoconvex defining function for $X$. Moreover $\alpha$ is a contact form on $X$, and $\Sigma$ is a symplectic submanifold of $T^*X$.
\end{prop}

\begin{proof}
Let us prove that $\alpha$ is a contact form; this follows from an explicit computation in local coordinates
\[
\alpha\wedge (\dd\alpha)^{n-1}=\|\nabla
\rho\|
\det([\partial_{z_j}\partial_{\bar{z_k}}\rho]_{j,k})\dd \voll,
\]
where $\dd \voll$ is the standard volume form in coordinates and the Hessian of $\rho$ is restricted to $T^{(0,1)}X\otimes T^{(1,0)}X$, where it is still non-degenerate.

The other claims follow directly from there.
\end{proof}

Given $u\in C^1(X,\C)$, one naturally forms $\bar{\partial}_bu$ as the section of $(T^{(0,1)}X)^*$ obtained by restricting $\dd u$ to $T^{(0,1)}X$. Restrictions to $X$ of holomorphic functions naturally satisfy $\bar{\partial}_bu=0$, and reciprocally elements of $\ker_{L^2}
\bar{\partial}_b$ extend into $H^{\frac 12}$ holomorphic functions on $\Omega$ in the strongly pseudoconvex case~\cite{kohn_harmonic_1963}.

In this article we will only be interested in the situation where $X$ (that is, $\rho$) is real-analytic. The symbol of the differential operator $\bar{\partial}_b$ can then be extended in the complex, which allows to define several manifolds of interest.

\begin{prop}\label{prop:Z_and_Sigma}
Let $Z=\{(x,\xi)\in \widetilde{T^*X},
\widetilde{\sigma(\bar{\partial}_b)}(x,\xi)=0\}$. Then $Z$ is a complex submanifold of $\widetilde{T^*X}$ of complex codimension $n-1$.

$Z$ is coisotropic, and defining $\bar{Z}=\{(x,\xi)\in \widetilde{T^*X},
\widetilde{\sigma(\partial_b)}(x,\xi)=0\}$, then $Z$ and $\bar{Z}$ intersect cleanly and
\[
Z\cap \bar{Z}=\widetilde{\Sigma}.
\]

Moreover $T\Sigma$ is transverse to the null-distribution $(TZ)^{\perp}=\{v\in T(\widetilde{T^*X}), \forall v'\in TZ, \omega(v,v')=0\}$.
\end{prop}

\begin{proof}
Let us consider holomorphic coordinates mapping a neighbourhood in $Y$ of a point of $X$ to a neighbourhood of $0$ in $\C^n$, with $\partial_{\bar{z_n}}\rho\neq 0$. Then as a subset of $\widetilde{T^*\C^n}=\{(x,y,\xi,\eta)\in \C^{4n}\}$, $Z$ is the joint zero set of the functions
\begin{align*}
p_j:(x,y,\xi,\eta)&\mapsto
\xi_j+i\eta_j-\frac{(\partial_{x_j}+i\partial_{y_j})\widetilde{\rho}(x,y)}{(\partial_{x_n}+i\partial_{y_n})\widetilde{\rho}(x,y)}(\xi_n+i\eta_n)&
&1\leq j \leq n-1,\\
\widetilde{\rho}:(x,y,\xi,\eta)&\mapsto \widetilde{\rho}(x,y),\\
\widetilde{\dd \rho}:(x,y,\xi,\eta)&\mapsto \sum_{j=1}^n\xi_j\partial_{x_j}\widetilde{\rho}(x,y)+\eta_j\partial_{y_j}\widetilde{\rho}(x,y).
\end{align*}
The differentials of these functions are linearly independent, and one can check that the functions $p_j$ Poisson-commute with each other. Therefore $Z$ is a coisotropic submanifold of $\widetilde{T^*X}$ whose codimension is the number of functions $p_j$ (i.e. $n-1$).

To define $\bar{Z}$, we replace $p_j$ with the functions
\begin{align*}
p_j^*:(x,y,\xi,\eta)\mapsto
\xi_j-i\eta_j-\frac{(\partial_{x_j}-i\partial_{y_j})\widetilde{\rho}(x,y)}{(\partial_{x_n}-i\partial_{y_n})\widetilde{\rho}(x,y)}(\xi_n-i\eta_n),&
 &1\leq j \leq n-1.
\end{align*}
These functions again Poisson-commute with each other.

Now, clearly $\Sigma=Z_{\R}=(Z\cap \bar{Z})_{\R}$. The symplectic gradients of the functions $p_1,\dots,p_{n-1},p_1^*,\dots,p_{n-1}^*$ are linearly independent on $\Sigma$, and in fact the matrix of Poisson brackets $\frac{1}{i}[\{p_j,p_k^*\}]_{j,k}$ is everywhere either positive definite (on the half-line bundle $\Sigma_+$ of positive multiples of $\alpha$) or negative definite (on the half-line bundle $\Sigma_-$ of negative multiples of $\alpha$), see the computations in~\cite{boutet_de_monvel_sur_1975}. This concludes the proof; in particular, the Hamiltonian vector fields of $p_1,\dots,p_{n-1}$, which span $(TZ)^{\perp}$, are transverse to $\Sigma$.
\end{proof}

From now on we decompose $\Sigma=\Sigma_+\cup\Sigma_-\cup\{0\}$ where $\Sigma_\pm$ is the half-line bundle of positive multiples of $\pm
\alpha$, on which $\pm\frac{1}{i}[\{p_j,p_k^*\}]_{j,k}$ is positive definite.

An important consequence of the geometry of $\Sigma$ and $Z$ is the following CR-extension result for $\Sigma$-Lagrangians.

\begin{prop}\label{prop:completing_Lagrangians}
Let $\lambda\subset \Sigma_+\times \Sigma_+$ be a real-analytic conical Lagrangian for the twisted symplectic form $\omega_1-\omega_2$.

There exists a unique $\Lambda\subset Z\times \bar{Z}$, Lagrangian in $T^*X\times T^*X$ for the twisted symplectic form, which contains $\lambda$. Moreover $\Lambda$ is positive (in the sense of~\cite[Definition~3.3]{melin_fourier_1975}).
\end{prop}

\begin{proof}
We proceed as in~\cite{boutet_de_monvel_sur_1975}. Let $\Lambda\subset Z\times \bar{Z}$ be a Lagrangian which contains~$\widetilde{\lambda}$. Then the null-distribution $(TZ^{\perp}\times
\{0\})\oplus(\{0\}\times T\bar{Z}^{\perp})$ belongs to $T\Lambda$. By Proposition~\ref{prop:Z_and_Sigma}, this distribution is involutive (because $Z$ is coisotropic); it is transverse to $T\tilde{\lambda}$; its complex dimension is $2(n-1)=\dim(\Lambda)-\dim(\widetilde{\lambda})$. Therefore $\Lambda$ is exactly the union of the leaves of this distribution which pass through $\widetilde{\lambda}$. The positivity of $\Lambda$ is then a direct consequence of the positivity of the matrix $\frac 1i[\{p_j,p_k^*\}]_{j,k}$ on $\Sigma_+$.
\end{proof}

\subsection{Classical normal form}\label{sec:class-norm-form}
At the level of the operator $\bar{\partial}_b$, the strong pseudoconvexity of $X$ is reflected in the following facts:
\begin{itemize}
\item the real characteristic $\Sigma=\{(x,\xi)\in T^*X,\sigma(\bar{\partial}_b)=0\}$ is symplectic;
\item denoting $(p_1,\dots,p_{n-1})$ the components of $\sigma(\bar{\partial}_b)$, one has $\{p_j,p_k\}=0$ for all $j,k$, while the matrix $\frac{1}{i}[\{p_j,\bar{p_k}\}]_{j,k}$ is everywhere either positive definite or negative definite.
\end{itemize}
There is a universal local model (depending only on the dimension) for the principal symbols of such operators, whose construction is presented in~\cite[Appendix~II]{boutet_de_monvel_hypoelliptic_1974} in the smooth case. We review this proof in the analytic setting. This local model is an important tool in the description of the Szeg\H{o} projector in the smooth case, and we will also use it in our setting. The local model we will use is the vector-valued differential operator $D_0$ acting on $\R^{n-1}_x\times \R^n_y$:
\begin{align*}
&(D_0)_j=-i\frac{\partial}{\partial x_j}-x_j\frac{\partial}{\partial y_1},& &1\leq j \leq n-1.
\end{align*}
This differs from the local model considered in~\cite{boutet_de_monvel_sur_1975}. Our choice of $D_0$, with a polynomial total symbol, will facilitate handling subprincipal terms in the quantum normal form in Section~\ref{sec:norm-form-overl}.



\begin{lemm}[{See also
\cite[Proposition~10.10]{boutet_de_monvel_hypoelliptic_1974} for the $C^{\infty}$ statement}]
\label{prop:extend_zero_bracket}
Let $1\leq r\leq m$. Let $\Sigma\subset T^*\R^m$ be an open cone. Let $p_1,\dots,p_r$ be $\frac 12$-homogeneous, real-analytic, complex-valued functions near $\Sigma$, such that:
\begin{itemize}
\item all the Poisson brackets $\{p_j,p_k\}$ and $\{\bar{p_j},\bar{p_k}\}$, for $1\leq j,k\leq r$, vanish on $\Sigma$;
\item all the Poisson brackets $\frac{1}{i}\{p_j,\bar{p_k}\}$, for $1\leq j,k\leq r$, are equal to $\delta_{jk}$ on $\Sigma$.
\end{itemize}
Then, in a conical neighbourhood of $\Sigma$, there exist $p_1',\dots,p_r'$, homogeneous, real-analytic, equal respectively to $p_1,\dots,p_r$ along with their first derivatives on $\Sigma$, and such that $\{p'_j,p'_k\}=\{p'_j,\bar{p'_k}\}=\frac{1}{i}\{p'_j,\bar{p'_k}\}-\delta_{jk}$ is equal to $0$ everywhere, and the ideal generated by $p_1,\dots,p_r$ coincides with the ideal generated by $p_1',\dots,p_r'$.
\end{lemm}
\begin{proof}
We proceed by induction on $r$. The claim follows immediately from the following two properties:
\begin{enumerate}
\item\label{prflemm2.5prop1} Let $V\subset T^*\R^n$ be a conical open set. Let $p:V\to \C$ be $\frac 12$-homogeneous, real-analytic; suppose that $\{p,\bar{p}\}$ bounded away from $0$ on $V$. (In particular, $\dd p$ is bounded away from $0$ on $V$.)

Then there exists a small neighbourhood $W$ of $\{p=0\}$, and $p':W\to \C$, with the same properties, such that $\{\widetilde{p}=0\}=\{\widetilde{p'}=0\}$, and $\frac{1}{i}\{p',\bar{p'}\}=1$. Moreover, for every function $q:V\to \C$ such that $\{p,q\}=\{\bar{p},q\}=0$, one has also $\{p,q\}=\{p',\bar{q}\}=0.$
\item\label{prflemm2.5prop2} Let $V\subset T^*\R^n$ be a conical open set. Let $p_1,\dots,p_r,q:V\to \C$ be $\frac 12$-homogeneous and real-analytic; suppose that $\{\widetilde{p}_1=\dots=\widetilde{p}_r=0\}$ is a regular energy level near $V$, and that $\{p_j,p_k\}=\{\bar{p_j},\bar{p_k}\}=\frac{1}{i}\{p_j,\bar{p_k}\}-\delta_{jk}=0$. Suppose also that $\{p_j,q\},\{\bar{p_j},q\}$ vanish on $\{p_j=\bar{p_j}=0\}$, on which $\{q,\bar{q}\}$ is bounded away from $0$.

Then there exists a small neighbourhood $W$ of $\{p_j\mk=\bar{p_j}\mk=0\}$ and ${q':W\mk\to\mk \C}$ a $\frac{1}{2}$-homogeneous, real-analytic function, which is equal to $q$ along with its first total differential on $\{p_j=\bar{p_j}=0\}$ (in particular $\{q',\bar{q'}\}$ is bounded away from $0$ and the ideal generated by $p_1,\dots,p_r,q$ coincides with the ideal generated by $p_1,\dots,p_r,q'$), and such that $\{p_j,q\}=\{\bar{p_j},q\}=0$ everywhere.
\end{enumerate}
Let us prove property~\eqref{prflemm2.5prop1} first. We want to correct $p$ into $p'=ap$, where $a$ is real-valued. After this modification, the symplectic bracket with the complex conjugate reads
\[
\{p',\bar{p'}\}=|a|^2\{p,\bar{p}\}+a\bar{p}\{p,a\}+ap\{a,\bar{p}\}.
\]
Letting $A=a^2$, this simplifies into
\[
\{p',\bar{p'}\}=A\{p,\bar{p}\}+\frac 12 p\{A,\bar{p}\}+\frac 12 \bar{p}\{p,A\}.
\]
Therefore we want to solve the transport equation
\[
Y\cdot A + \frac{1}{i}\{p,\bar{p}\}A=1
\]
where $Y$ is the (real-valued) vector field $\frac{1}{i}(-p\Xi_{\bar{p}}+\bar{p}\Xi_p)$ (where $\Xi_f$ denotes the symplectic gradient of $f$).

This vector field is singular along the codimension~2 set $\{Y=0\}=\{p=\bar{p}=0\}$. Moreover it enjoys an absorbing property: all trajectories of $\pm Y$ converge exponentially fast to the singularity in positive time. Indeed, since $\frac{1}{i}\{p,\bar{p}\}$ has constant sign, $|p|^2$ is a Lyapunov function for this vector field.

Therefore, by the method of characteristics, there exists a unique $A$ in a neighbourhood of $\{Y=0\}$, real-analytic, and such that $Y\cdot A + \frac{1}{i}\{p,\bar{p}\}A=1$, and moreover it is real-valued.

If $q$ is such that $\{p,q\}=\{\bar{p},q\}=0$, then both $Y$ and $\{p,\bar{p}\}$ commute with $q$, so that $A$ does as well.

To prove Property~\eqref{prflemm2.5prop2}, we simply observe that the complex submanifold $\{\widetilde{p_j}=\widetilde{\bar{p_j}}=0\}$ is transverse to the flows of $\widetilde{p_j}$ and $\widetilde{\bar{p_j}}$. We can therefore set $q'=\widetilde{q}$ on this set, then extend it to be constant along the flows of $\widetilde{p_j}$ and $\widetilde{\bar{p_j}}$.
\end{proof}


\begin{prop}\label{prop:canon_transf_to_local_model}
Let $P_0\in \Sigma_{\pm}$. There exists a real-analytic contact transformation $\kappa$ from a conical neighbourhood of $P_0$ in $T^*X\setminus \{0\}$ to a conical neighbourhood of $(0,0, 0,(\pm 1,0))$ in $T^*(\R^{n-1}\times \R^{n})$, mapping $P_0$ to $(0,0,0,(\pm 1, 0))$, and a real-analytic, $0$-homogenous map $C$ from a conical neighbourhood of $x_0$ in $T^*X\setminus \{0\}$ to $GL(\C^{n-1},\Omega^{(1,0)}(X))$, such that
\[
\sigma(\bar{\partial}_b)(z,\zeta)=C\sigma(D_0)(\kappa(z,\zeta)).
\]
\end{prop}

\begin{proof}[{Proof of Proposition~\ref{prop:canon_transf_to_local_model}}]
Let us apply Lemma~\ref{prop:extend_zero_bracket} to the $\bar{\partial}_b$ operator: the components $Z_1,\dots,Z_{n-1}$ of its symbol satisfy, near any point of the characteristic set $\Sigma^{\pm}$,
\begin{align*}
&\{Z_j,Z_k\}=0 \text{ on }\Sigma, & &\pm\frac{1}{i}\{Z_j,\bar{Z_k}\}\gg 0.
\end{align*}
Let $M:\Sigma^+\to GL(n-1)$ be the positive square root of the matrix $\pm\frac{1}{i}\{Z_j,\bar{Z_k}\}$; extend $M$ arbitrarily into a $\frac 12$-homogeneous, real-analytic function from a neighbourhood of $\Sigma^+$ into $GL(n-1)$. Let $z=M^{-1}Z$. The ideal generated by $z$ coincides with that generated by $Z$, and on particular $z$ vanishes on $\Sigma^+$, where
\begin{align*}
\{z_j,z_k\} &=\sum_{l,m}(M^{-1})_{jl}(M^{-1})_{mk}\{Z_l,Z_m\}=0,
\\
\intertext{whereas}
\frac{1}{i}\{z_j,\bar{z_k}\} &=\sum_{l,m}(M^{-1})_{jl}\bar{(M^{-1})_{km}}\frac{1}{i}\{Z_l,\bar{Z_m}\}=\pm
\delta_{jk}.
\end{align*}
Then, we let $u_j=\re(z_j)$ and $v_j=\pm \im(z_j)$ and directly obtain the situation of Lemma~\ref{prop:extend_zero_bracket}:
\begin{align*}
&\{u_j,u_k\}=\{v_j,v_k\}=0,& &\{v_j,v_k\}=\delta_{jk}.
\end{align*}

Letting $z'_j=u'_j+iv'_j$, then the ideal generated by $z'$ coincides with that generated by $Z$.

We are almost ready for the normal form. We let $w$ be any non-vanishing,\linebreak $\frac 12$-homogeneous, real-valued function on a neighbourhood of $\Sigma$ which commutes with $z'_1,\dots,z'_{d}$, and we consider the following functions:
\[
x_j=u'_jw^{-1}, \qquad \xi_j=v_j'w, \qquad \eta_1=\pm w^2.
\]
By construction, the $x_j$ are $0$-homogeneous while $\eta_1$ and the $\xi_j$ are $1$-homogeneous, and $\{x_j,x_k\}=\{x_j,\eta_1\}=\{\xi_j,\eta_1\}=\{x_j,\xi_k\}-\delta_{jk}=0$. The ideal of real-analytic functions generated by the one-homogeneous family $(d_0)_j=x_j\eta_1+i\xi_j=wz'_j$ coincides with the ideal of real-analytic functions generated by $z'$, which is the same as the ideal generated by $Z$. All in all, we obtain that, for some invertible matrix $C$, one has $Z=Cd_0$.

Completing the coordinates $(x_1,\dots,x_{d},\xi_1,\dots,\xi_{d},\eta_1)$ into a canonical transformation of a neighbourhood of $\Sigma^+$ into a neighbourhood of $x=\xi=0,y=0,\linebreak\pm \eta_1>0,(\eta_2,\dots,\eta_n)=0$ in $T^*\R^{n-1}\times T^*\R^n$, we can exhibit a canonical transform $\kappa$ as requested.
\end{proof}


\subsection{Results from the smooth theory}\label{sec:results-from-smooth}

Propositions~\ref{prop:Z_and_Sigma}, \ref{prop:completing_Lagrangians}, \ref{prop:extend_zero_bracket} and~\ref{prop:canon_transf_to_local_model} are also valid in the $C^{\infty}$ category, after a suitable replacement of holomorphic extensions by almost holomorphic extensions (in fact, our method of proof for the classical normal form applies to this situation as well, and is much simpler than the one presented in~\cite[Appendix~2]{boutet_de_monvel_hypoelliptic_1974}). The main result of~\cite{boutet_de_monvel_sur_1975} is then that, with $\lambda=\diag(\Sigma_+)$, and $\Lambda$ as in Proposition~\ref{prop:completing_Lagrangians}, then the Szeg\H{o} projector $S$ on $X$ is a {Fourier integral} operator associated with $\Lambda$, modulo an operator which sends $\clD'(X)$ into $\clD(X)$. Since $\Lambda$ is a half-line bundle over $X\times X$, one can represent this {Fourier integral} operator, globally, using a phase function with one phase variable, as follows:
\[
S(x,y)=\int_0^{+\infty}e^{it\psi(x,y)}s(x,y;t)\dd t.
\]
One can construct $\psi$ using the almost holomorphic extension of a defining function~$\rho$.

The {Lagrangian} $\Lambda$ is idempotent (and indeed $S$ is a projector), and a study of {Fourier integral operators} with {Lagrangian} $\Lambda$ is performed in~\cite{boutet_de_monvel_spectral_1981}. This book avoids the notion of {Fourier integral} operators with complex phases and only considers operators of the form $\mathit{SPS}$ where $P$ is a pseudodifferential operator on $X$, but it turns out that general {Fourier integral operators} with Lagrangian $\Lambda$, in the smooth category, are of this form. We will use the following precise version of this result.

\begin{prop}\label{prop:repr_BG_smooth}
Let $A$ be a (smooth) {Fourier integral} operator with {Lagrangian} $\Lambda$. Then there exists a pseudodifferential operator $Q$ on $X$ such that $[Q,S]$ and $A-SQD$ are continuous from $\clD'$ to $\clD$.
\end{prop}

\begin{proof}
By induction we only need to prove the result at principal order, and by~\cite[Proposition~2.13]{boutet_de_monvel_spectral_1981} we may drop the condition that $Q$ and $S$ commute. Therefore we are only searching for the principal symbol $q_0$ of a pseudodifferential operator $Q$, such that the principal symbol of $\mathit{SQS}$ is that of $A$. An application of stationary phase shows that $\mathit{SQS}$ has an distributional kernel of the form
\[
\mathit{SQS}(x,y)=\int_0^{+\infty}e^{it\psi(x,y)}a(x,y;t)\dd t,
\]
with principal symbol $a:\Sigma_+\to \R$ equal to $q_0$. Hence, given $A$, we can extend arbitrarily its principal symbol to the whole of $T^*X$ and obtain an approximation at first order.
\end{proof}

\subsection{The case of the cylinder}\label{sec:an-exampl-eucl}

A particular case of a pseudoconvex hypersurface (albeit not compact) is $\{z\in \C^n,\linebreak \lvert\im(z)|=1\}$; indeed the function $\rho:z\mapsto \lvert\im(z)|^2$ is strongly plurisubharmonic everywhere.

There is an explicit description of the Szeg\H{o} kernel in this case, and one can directly establish many properties which we will generalise only in Section~\ref{sec:szegho-kern-param}. This particular case will be useful in our construction and characterisations of general {Fourier integral operators} in real-analytic regularity in Section~\ref{sec:pseud-oper-1}.

\begin{prop}\label{prop:flat_Grauert}
The Szeg\H{o} kernel on $\R^n\times S^{n-1}$ is
\[
K(z,w)=\frac{1}{(2\pi)^n}\int_{\R^n}\frac{e^{i(z-\bar{w})\cdot
\xi}}{\int_{S^{n-1}}e^{2y\cdot \xi}\dd y}\dd \xi,
\]
which defines a holomorphic function of $(z,\bar{w})$ on $(\R^n\times \overline{B_{\R^n}(0,1)})^2\setminus \diag(\R^n\times S^{n-1})$.

In particular, the Szeg\H{o} projector $S$ has a distributional kernel analytic away from the diagonal.
\end{prop}

A related formula concerning the Bergman kernel appears on p.~100 of the manuscript~\cite{epstein_shrinking_1990}. %Statuer sur la forme du cite en relecture.
\begin{proof}
The first step is the study of the function
\[
m_n:\xi\mapsto \int_{S^{n-1}}e^{y\cdot \xi}\dd y.
\]
This defines an holomorphic function on $\C^n$. When restricted to $\R^n$, it is a radial function and, by the Laplace method, its asymptotics for large argument are
\begin{equation}\label{eq:mn_first_term}
m_n(r\omega)=2^{\frac{n-1}{2}}|S^{n-2}|\Gamma\left(\frac{n+1}{2}\right)r^{-\frac{n+1}{2}}e^r+O_{r\to +\infty}\left(e^rr^{-\frac{n+3}{2}}\right),
\end{equation}
where $|S^{n-2}|$ is the area of the unit $n-2$-sphere. Let $(z,w)\in (\R^n\times \bar{B_{\R^n}(0,1)})^2\setminus \diag(\R^n\times S^{n-1})$. Then $\lvert\im(z-\bar{w})|\leq 2$ and furthermore $\lvert\im(z-\bar{w})|=2\implies
\re(z-\bar{w})\neq 0$. Thus, by formula~\eqref{eq:mn_first_term}, the integral
\[
\int_{\R^n}e^{i(z-\bar{w})\cdot \xi}(m_n(2\xi))^{-1}\dd \xi
\]
is well-defined, real-analytic on its domain, and holomorphic with respect to $(z,\bar{w})$.


One can directly check that $K(z,w)$ is self-adjoint.

It remains to prove that, for every $f\in L^2(\R^n\times S^{n-1})$ which extends into a holomorphic function inside $\R^n\times B(0,1)$, one has $f(z)=\int K(z,w)f(w)\dd w$. By~density, it suffices to prove this fact when $f\in
\clS(\R^n\times S^{n-1})$. Given such an $f$, its Fourier transform with respect to the first variable reads
\begin{align*}
\hat{f}(\xi;y) &=\frac{1}{(2\pi)^{\frac n2}}\int_{\R^n} e^{-i\xi\cdot x'}f(x'+iy)\dd x'
\\
\intertext{in the sense that}
f(x+iy) &=\frac{1}{(2\pi)^{\frac n2}}\int_{\R^n} e^{i\xi\cdot x}\hat{f}(\xi;y)\dd \xi.
\end{align*}
We now apply a contour deformation to the integral defining $\hat{f}(\xi;y)$, replacing $x'$ with $x'+i(w-y)$ for $w\in
\S^{n-1}$. We are at liberty to do so: in the homotopy of contours $x'\mapsto x'+it(w-y)$ for $t\in [0,1]$, on $\{t\in\mathopen{]}0,1\mathclose{[}\}$ the function is holomorphic, and it extends smoothly up to the boundary. The fact that $f\in \clS$ ensures sufficient decay as $|x'|\to +\infty$ to guarantee that one can apply the contour deformation.

Thus, for every $w\in S^{n-1}$,
\[
\hat{f}(\xi;y)=\frac{1}{(2\pi)^{\frac n2}}\int_{\R^n} e^{-i\xi\cdot (iy-(x'+iw))}f(x'+iw)\dd x'.
\]
We now average this over $w\in S^{n-1}$ with respect to the probability density
\[
w\mapsto \frac{e^{-2w\cdot \xi}}{m_n(2\xi)},
\]
which yields
\[
\hat{f}(\xi;y)=\frac{1}{(2\pi)^{\frac n2}}\int_{\R^n\times S^{n-1}}e^{i\xi\cdot(iy-(x'-iw))}(m_n(2\xi))^{-1}f(x'+iw)\dd x' \dd w.
\]
Applying now the Fourier inversion formula, we obtain finally
\[
f(x+iy)=\frac{1}{(2\pi)^n}\int_{\R^n\times \R^n\times S^{n-1}}e^{i\xi\cdot(x+iy-(x'-iw))}(m_n(2\xi))^{-1}f(x'+iw)\dd \xi\dd x' \dd w.
\]
This concludes the proof of Proposition~\ref{prop:flat_Grauert}.
\end{proof}

\begin{prop}\label{prop:loc-anal-hypoell-Grauert}
$\bar{\partial}_b$ is locally analytic hypoelliptic on $\R^n\times S^{n-1}$, by which we mean the following property: if $u\in L^2(\R^n\times S^{n-1})$ and $\bar{\partial}_bu$ is real-analytic on an open set $V$, then $(1-S)u$ is real-analytic on $V$.
\end{prop}
\begin{proof}
It suffices to prove the claim in the case where $V$ is bounded. Let $V_1\Subset V$ be a smaller open set and let $(\chi_N)_{N\in \N}$ be a family of smooth functions supported on $V$, such that $\chi_N=1$ on $V_1$, and
\[
\exists \rho, \;\forall j\leq N,\qquad |\nabla^j\chi_N|\leq (\rho N)^j.
\]
Such cutoffs will be used systematically in Section~\ref{sec:ampl-analyt-micr}, and we postpone until Proposition~\ref{prop:Ehrenpreis} the proof of their existence. Let us recall from the Leibniz formula that
\[
u\in C^{\omega}(V)\implies \; \exists C,R,\; \forall n\in \N,\qquad \|\nabla^n(\chi_nu)\|_{L^2}\leq C(Rn)^n\implies u\in C^{\omega}(V_1).
\]
Since $V$ is relatively compact, we have more generally the following characterisation: let $(Z_1,\dots,Z_k)$ be a finite family of real-analytic vector fields on a neighbourhood of $V$, which span $TV$ at each point. Then
\begin{multline*}
u\in C^{\omega}(V)\implies \;\exists C,R,\;\forall n\in \N, \;\forall \mu\in
\llbracket 1,k\rrbracket^n,\\
\left\|Z_{\mu_1}Z_{\mu_2}\dots Z_{\mu_n}(\chi_nu)\right\|_{L^2}\leq C(Rn)^n\implies u\in C^{\omega}(V_1).
\end{multline*}
Here, by convention, when $Z$ is a vector field on $\R^n\times S^{n-1}$ and $u$ is a function on $\R^n\times S^{n-1}$, we call $Zu$ the Lie derivative of $u$ with respect to $Z$.

The following vector fields are particularly suited to our problem:
\begin{align*}
X_j&=\frac{\partial}{\partial x_j} & &1\leq j\leq n,\\
Y_{j,k}&=\omega_k\frac{\partial}{\partial
\omega_j}-\omega_j\frac{\partial}{\partial
\omega_k}-x_k\frac{\partial}{\partial x_j}+x_j\frac{\partial}{\partial x_k} & &1\leq j,k\leq n.
\end{align*}
Indeed, these vector fields span $T(\R^n\times S^{n-1})$, and they preserve all of the geometry: $X_j$ acts by translation on the $x_j$ variable, and $Y_{j,k}$ acts by rotation in the plane of the variables $(y_j,y_k)$. Thus, the Lie derivatives along these vector fields all commute with $\bar{\partial}_b$ and its adjoint $\bar{\partial}_b^*$, hence they commute with the Szeg\H{o} projector $S$. Now, given $u\in L^2$ such that $\bar{\partial}_bu$ is real-analytic on $V$, one has, for every $n\in \N$ and every multi-index $\mu\in \llbracket 1,k\rrbracket^n$,
\begin{align*}
\|Z_{\mu_1}\dots Z_{\mu_n}(1-S)(\chi_nu)\|_{L^2}&=\|(1-S)Z_{\mu_1}\dots Z_{\mu_n}(\chi_nu)\|_{L^2}\\
&\leq \left\|\bar{\partial}_bZ_{\mu_1}\dots Z_{\mu_n}(\chi_nu)\right\|_{L^2}\\
&=\left\|Z_{\mu_1}\dots Z_{\mu_n}\bar{\partial}_b(\chi_nu)\right\|_{L^2};
\end{align*}
where we used Kohn's hypoellipticity result $\|(1-S)u\|_{L^2}\leq
\|\bar{\partial}_bu\|_{L^2}$~\cite{kohn_harmonic_1963}.

To conclude, since $\chi_n=1$ on $V_1$ and the kernel of $S$ is analytic away from the diagonal by Proposition~\ref{prop:flat_Grauert}, then for every $V_2\Subset V_1$, the sequence $((1-S)(1-\chi_n)u)_{n\in \N}$ is bounded in some fixed space of analytic functions on $V_2$. Therefore we can conclude that
\[
\exists C,R, \forall \mu\in \llbracket 1,k\rrbracket^n,\qquad\|Z_{\mu_1}\dots Z_{\mu_n}\bar{\partial}_b(1-S)u\|_{L^2(V_2)}\leq C(Rn)^n,
\]
and the proof of Proposition~\ref{prop:loc-anal-hypoell-Grauert} is complete: $(1-S)u$ is real-analytic on $V_2$.
\end{proof}

\section{Amplitudes and phases}\label{sec:analytic-symbols}

This section serves as a technical preliminary for the definition and study of {Fourier integral operators} in Section~\ref{sec:four-integr-oper}. We study in detail spaces of amplitudes adapted to the real-analytic regularity, and we state and prove a stationary phase result for these spaces of amplitudes.


\subsection{Amplitudes in analytic microlocal analysis}
\label{sec:ampl-analyt-micr}

This section takes inspiration from~\cite{deleporte_toeplitz_2018,sjostrand_singularites_1982,treves_analytic_2022}.



\begin{enonce*}{Convention}
Given $U\subset \R^n$ open, $x\in U$, and $u\in C^j(U,\C)$, we let
\[
\left|\nabla^ju(x)\right|=\left(\sum_{\alpha\in
\N^n,|\alpha|=j}|\partial^{\alpha}u(x)|^2\right)^{\frac 12}.
\]
\end{enonce*}

\begin{defi}\label{def:formal-amp}
Let $\Omega\subset \R^m_x\times \R^n_{\theta}$ be an open cone in the second variable, and let $d\in \R$. A \emph{(homogeneous classical) formal analytic amplitude} of degree $d$ is a sequence $(a_k)$ of real-analytic functions on $\Omega$ such that $a_k$ is homogeneous of degree $d-k$, and such that there exists $C>0, \rho>0,R>0,m>0$ satisfying the following: for every $j,k,\ell\in \N_0$, for every $(x,\theta)\in \Omega$,
\[
\left|\nabla^{j}_x\nabla^{k}_{\theta}a_\ell(x,\theta)\right|\leq C\frac{\rho^{j+k}R^\ell(j+k+\ell)!}{(1+j+k+\ell)^m}|\theta|^{d-\ell-k}.
\]
For every $m,\rho,R$, the best constant $C$ in this inequality is written $\|a\|_{S^{\rho,R}_m(\Omega)}$.
\end{defi}

If $(a_{\ell})_{\ell\in \N}$ is a formal analytic amplitude, then for every $\ell$, $a_{\ell}$ is a real-analytic function of $(x,\theta)$, which extends to a fixed size neighbourhood of the real locus, where it satisfies $\|a_{\ell}(1+|\theta|)^{-d+\ell}\|_{L^{\infty}}\leq C(R')^{\ell}\ell!$. This turns out to be an equivalent definition of a formal analytic amplitude. If $\Omega$ is compact with respect to the first variable, one-homogeneous analytic changes of variables in $x,\theta$ preserve the space of all formal analytic amplitudes, even though they may not individually preserve the Banach spaces $S^{\rho,R}_m(\Omega)$.

Given a paracompact conical analytic manifold, we may then define formal analytic amplitudes on $X$ as sequences of real-analytic functions on $X$ with decreasing homogeneity, which satisfy the controls of Definition~\ref{def:formal-amp} locally in charts.



\begin{defi}\label{def:analytic_amp}
Let $\Omega\subset \R^m_x\times \R^n_{\theta}$ be an open cone in the second variable, and let $d\in \R$. An \emph{analytic amplitude} of degree $d$ on $\Omega$ is a smooth function $a$ on $\Omega$ such that there exists $R>0,\rho>0,C>0$ satisfying the following:
\[
\forall (x,\theta)\in \Omega,\;\forall k\leq |\theta|/R,\;\forall j,\qquad \left|\nabla_{x}^j\nabla_{\theta}^ka(x,\theta)\right|\leq C\rho^{j+k}(j+k)!|\theta|^{d-k}.
\]
\end{defi}

An analytic amplitude is real-analytic with respect to $x$, and extends into a holomorphic function of $x$; an alternative definition of an analytic amplitude is the estimate
\[
\forall k \leq |\theta|/R,\qquad \left|\nabla_{\theta}^ka(x,\theta)\right|\leq C\rho^kk!|\theta|^{d-k}
\]
on a whole neighbourhood of the real locus in $x$. Thus, again, we may speak of analytic amplitudes on $\Omega\subset X\times \R^n_{\theta}$ when $X$ is a paracompact analytic manifold.

\begin{rema}\label{rem:name_analytic_ampl}
Despite their names, the analytic amplitudes as in Definition~\ref{def:analytic_amp} are \emph{not} real-analytic functions. The main reason for this definition is that the Borel summation of a formal analytic amplitude into an analytic amplitude (Proposition~\ref{prop:real_formal_symb}) will involve a lowest term summation, which requires cut-off functions in the variable $|\theta|$. Trèves calls such functions ``pseudo-analytic''~\cite{treves_analytic_2022} but unfortunately this term, like many other variations upon the word ``analytic'', already refer to well-established mathematical objects.
\end{rema}

Analytic amplitudes are stable by ``phasification'' of the variables: denoting $x=(x_1,x_2)$, where $x_1\in \R^{m_1}$ and $x_2\in \R^{m_2}$, and similarly $\R^{m_2+n}\ni \theta=(\theta_1,\theta_2)\ni \R^{m_2}\times \R^n$, if some amplitude $a$ is analytic on $\Omega\subset \R^{m_1+m_2}\times (\R^{n}\setminus \{0\})$, then so is
\[
(\R^{m_1})\times (\R^{m_2}\times (\R^n\setminus \{0\}))\supset
\Omega \ni (x_1,\theta)\mapsto a((x_1,\theta_1/|\theta_2|),\theta_2).
\]
The converse, however, is not true. We will say that an amplitude $a$ is \emph{maximally analytic} if, writing $\theta=\lambda\omega$ in spherical coordinates, the function $((x,\omega),\lambda)\mapsto a(x,\lambda\omega)$ is an analytic amplitude.

A last immediate consequence of Definition~\ref{def:analytic_amp} is that analytic amplitudes belong to the usual H\"ormander symbol classes, even as they are extended holomorphically with respect to the first variable.


In order to build analytic amplitudes starting from formal analytic amplitudes, the following construction of cut-off functions will be useful.

\begin{prop}[{See
\cite[Lemma~2.2]{hormander_uniqueness_1971} and~\cite[Section~3.2 (Ehrenpreis cutoffs)]{treves_analytic_2022}}]
\label{prop:Ehrenpreis}
Let $K,L$ be closed sets of $\R^d$ at positive distance from each other. There exists a sequence of compactly supported functions $(\chi_N)_{N\geq 1}\in C^{\infty}(\R^d,[0,1])$ such that $\chi_N=0$ on $L$ and $\chi_N=1$ on $K$, and a constant $\rho>0$ such that
\begin{equation}\label{eq:Ehrenpreis}
\forall N\in \N,\; \forall j\in \N,\qquad j\leq N\implies |\nabla^j\chi_N|\leq (\rho N)^j.
\end{equation}
\end{prop}
\begin{proof}
Let $K\subset K'$ such that $K'\cap L=\partial K'\cap K=\emptyset $, and let
\[
c=\min\bigl(\dist(K',L),\;\dist(K,\partial K')\bigr)>0.
\]

Let $\phi\in C^{\infty}(\R^d,[0,+\infty))$ be supported on $B(0,1)$ and such that $\int_{\R^d}\phi=1$. Given $\alpha>0$, let
\[
\phi_{\alpha}:x\mapsto \alpha^{-d}\phi(\alpha^{-1}x).
\]
The claim follows from the following choice of cut-off functions:
\[
\chi_N=\underbrace{\phi_{cN^{-1}}*\dots *\phi_{cN^{-1}}}_{N\text{ times }}*\1_{K'}.
\]
Indeed, letting $\rho=\|\nabla \phi_c\|_{L^1}$, one has $\|\nabla
\phi_{cN^{-1}}\|_{L^1}=\rho N$, while $\|\phi_{cN^{-1}}\|_{L^1}=1$ and $\|\1_{K'}\|_{L^{\infty}}=1$; altogether we find the desired bound.
\end{proof}


\begin{prop}\label{prop:real_formal_symb}
Let $\Omega\subset \R^m_x\times \R^n_{\theta}$ be an open cone with respect to the second variable. Let $(a_k)_{k\in \N}$ be a degree $d$ formal analytic amplitude on $\Omega$. Then there exists a degree $d$ maximally analytic amplitude $a$ on $\Omega$, and constants $C,\rho$ such that, uniformly on a complex neighbourhood of $\Omega$ in the spherical variable, for every $N\in \N$,
\[
\left|a(x,\theta)-\sum_{k=0}^{N-1}a_k(x,\theta)\right|\leq C\rho^NN!|\theta|^{d-N}.
\]

Any analytic amplitude whose extension satisfies the bound above will be called a realisation of the formal amplitude $(a_k)_{k\in
\N}$. Two extensions of analytic amplitudes realising the same formal amplitude have a difference bounded by $e^{-c|\theta|}$, for some $c>0$, on a whole conical complex neighbourhood of $\Omega$ with respect to the first variable.
\end{prop}

\begin{proof}
Let us first prove the last fact: letting $\widetilde{\Omega}$ be a complex neighbourhood with respect to the first variable, suppose that
\[
\forall (x,\theta)\in \widetilde{\Omega}, \;\forall N\in \N,\qquad |a(x,\theta)|\leq C\rho^NN!|\theta|^{d-N}.
\]
Writing $N=\alpha|\theta|$ and applying the Stirling formula gives
\[
|(x,\theta)|\leq C|\theta|^d(\alpha\rho/e)^{\alpha|\theta|} ;
\]
For $\theta$ large enough there is an integer $N$ between $\frac{e|\theta|}{4\rho}$ and $\frac{e|\theta|}{2\rho}$, and we obtain
\[
|a(x,\theta)|\leq C|\theta|^d(1/2)^{\frac{e|\theta|}{4\rho}};
\]
hence the claim.

We now give a direct construction of $a$. Let $(\chi_N)_{N\in \N}:[0,+\infty)\to [0,1]$ be a sequence of Ehrenpreis cutoffs as in Proposition~\ref{prop:Ehrenpreis}, equal to $1$ on $[0,1]$ and to $0$ on $\R\setminus [0,2]$.

Given $(a_\ell)_\ell\in S^{\rho,R}_0(\Omega)$, the claim follows from the choice
\[
a:(x,\theta)\mapsto
\sum_{\ell=1}^{\infty}a_{\ell}(x,\theta)\left(1-\chi_{\ell+1}\left(\frac{c|\theta|}{\ell+1}\right)\right),
\]
where $c$ is chosen small enough depending on $R$.

From now on we pass to spherical coordinates and write $\theta=\lambda\omega,X=(x,\omega)$. Let us prove that $a$ is an analytic amplitude. Given $(X,\lambda)\in \Omega$, $j\in \N$ and $k\leq \lambda/R$, one has
\[
\nabla^j_{X}\partial^k_{\lambda}a(X,\lambda)=\sum_{\ell\leq c\lambda}\nabla^j_{X}\partial^k_{\lambda}(a_{\ell}(X,\lambda))\left(1-\chi_{\ell+1}\left(\frac{c\lambda}{\ell+1}\right)\right).
\]

Let us first control the terms where no derivatives hit $\chi_{\ell+1}$. If $c\leq \frac{1}{4R}$, then
\begin{align*}
\left|\sum_{\ell\leq c\lambda}\left(1-\chi_{\ell+1}\left(\frac{c\lambda}{\ell+1}\right)\right)\nabla^j_{X}\partial^k_{\lambda}a_\ell(X,\lambda)\right| &
\leq \|a\|_{S^{\rho,R}_0}\sum_{\ell\leq c\lambda}\lambda^{d-\ell-k}(j+k+\ell)!\rho^{j+k}R^{\ell}\\
&\leq
\|a\|_{S^{\rho,R}_0}(j+k)!(2\rho)^{j+k}\lambda^{d-k}\sum_{\ell\leq c\lambda}(2R)^{\ell}\ell!\lambda^{-\ell}\\
&\leq 2\|a\|_{S^{\rho,R}_0}(j+k)!(2\rho)^{j+k}\lambda^{d-k},
\end{align*}
where we used the inequality $(j+k+\ell)!\leq 2^{j+k+\ell}(j+k)!\ell!$ and the fact that, on $0\leq \ell\leq c\lambda$, the family $(2R)^{\ell}\ell!\lambda^{-\ell}$ decreases with $\ell$ faster than a geometric sequence of ratio $\frac 12$.

Now we treat the case where one of the derivatives has hit $\chi_{\ell+1}$. The support of $\lambda\mapsto
\chi_{\ell+1}'(\frac{c\lambda}{\ell+1})$ is included in $\frac{c\lambda}{2}\leq \ell+1\leq c\lambda$. If moreover $k\leq
\epsilon \lambda$ with $\epsilon \leq \frac{c}{2}$, then for every $m\leq k$, one has $m\leq \frac{c\lambda}{2}\leq \ell+1$. Therefore
\begin{multline*}
\left|\sum_{\ell\leq c\lambda}\sum_{m=1}^{k}\binom{k}{m}\partial_{\lambda}^m\chi_{\ell+1}(\frac{c\lambda}{\ell+1})\nabla^j_X\partial^{k-m}_{\lambda}a_{\ell}(X,\lambda)\right|\\
\begin{aligned}
&\leq\|a\|_{S^{\rho,R}_0}\sum_{\frac{c\lambda}{2}-1\leq\ell\leq c\lambda-1}\sum_{m=1}^k\binom{k}{m}c^m\rho^{j+k-m}R^\ell(j+k-m+\ell)!\lambda^{d-k-\ell+m}\\
&\leq \|a\|_{S^{\rho,R}_0}(4\rho)^{j+k}(j+k)!\lambda^{d-k}\sum_{\frac{c\lambda}{2}-1\leq\ell\leq c\lambda-1}\sum_{m=1}^kc^m(2\rho)^{-m}(2R)^{\ell}(\ell-m)!\lambda^{m-\ell}.
\end{aligned}
\end{multline*}
Since $m\leq \frac{\ell}{2}$, we perform a change of variables and obtain, provided $c\leq\frac{1}{4R}$,
\begin{multline*}
\left|\sum_{\ell\leq c\lambda}\sum_{m=1}^{k}\binom{k}{m}\partial_{\lambda}^m\chi_{\ell+1}\left(\frac{c\lambda}{\ell+1}\right)\nabla^j_X\partial^{k-m}_{\lambda}a_{\ell}(X,\lambda)\right|\\
\begin{aligned}
&\leq
\|a\|_{S^{\rho,R}_0}k(8\rho^2R)^{j+k}(j+k)!\lambda^{d-k}\sum_{0\leq p\leq c\lambda-1}(2R)^p\lambda^{-p}(p+1)!\\
&\leq 2\|a\|_{S^{\rho,R}_0}(16\rho^2R)^{j+k}(j+k)!\lambda^{d-k}
\end{aligned}
\end{multline*}
It remains to show that $a=\sum_{\ell=0}^Na_\ell+O_N(\lambda^{d-N-1})$ with the required growth in $N\geq 0$. First
\[
(X,\lambda)\mapsto\sum_{\ell=0}^Na_\ell(X,\lambda)\chi_{\ell+1}\left(\frac{c\lambda}{\ell+1}\right)
\]
has compact support, and is uniformly bounded as follows:
\[
\left|\sum_{\ell=0}^Na_\ell(X,\lambda)\chi_{\ell+1}\left(\frac{c\lambda}{\ell+1}\right)\right|
\leq
\|a\|_{S^{\rho,R}_0}\sum_{\ell=0}^NR^{\ell}\ell!\lambda^{d-\ell}\chi_{\ell+1}\left(\frac{c\lambda}{\ell+1}\right);
\]
for this sum to be nonzero, we need $c\lambda\leq N+1$, moreover $\ell\mapsto R^{\ell}\ell!\lambda^{d-\ell}$ is log-convex so that
\[
\left|\sum_{\ell=0}^Na_\ell(X,\lambda)\chi_{\ell+1}\left(\frac{c\lambda}{\ell+1}\right)\right|
\leq
\|a\|_{S^{\rho,R}_0}\lambda^d(N+1)\max\left(1,R^NN!\lambda^{-N}\right)\1_{\lambda\leq
\frac{N+1}{c}}.
\]
Moreover if $\lambda\leq \frac{N+1}{c}$ then $N!\lambda^{-N}\geq
\left(\frac{c}{e}\right)^N$, so that, all in all, for some $R_1$ large enough,
\[
\left|\sum_{\ell=0}^Na_\ell(X,\lambda)\chi_{\ell+1}\left(\frac{c\lambda}{\ell+1}\right)\right|
\leq
\|a\|_{S^{\rho,R}_0}R_1^{N+1}(N+1)!\lambda^{d-N-1}
\]

It remains to bound
\[
\left|\sum_{\ell=N+1}^{+\infty}a_\ell(X,\lambda)(1-\chi_{\ell+1}\left(\frac{c\lambda}{\ell+1}\right)\right|\leq \|a\|_{S^{\rho,R}_0}
\sum_{\ell=N+1}^{c\lambda}R^{\ell}\ell!\lambda^{d-\ell}.
\]
If $c\leq \frac{1}{2R}$, the first term in this sum is $R^{N+1}(N+1)!\lambda^{d-N-1}$, and the ratio between two consecutive terms is smaller than $\frac{1}{2}$; thus
\[
\left|\sum_{\ell=N+1}^{+\infty}a_\ell(X,\lambda)\chi_{\ell+1}\left(\frac{c\lambda}{\ell+1}\right)\right|\leq 2\|a\|_{S^{\rho,R}_0}R^{N+1}(N+1)!\lambda^{d-N-1}.
\]
This concludes the proof of Proposition~\ref{prop:real_formal_symb}.
\end{proof}

\begin{rema}\label{rem:Resurgence}
It is not clear to us which natural conditions one can impose on the formal amplitude $a$ to be the expansion at large $|\theta|$ of an analytic function; many estimates and considerations in this section are much easier in this case. This seems to be an instance of
\emph{resurgence}. It would be interesting to study if the space of such formal amplitudes is stable under the natural Fourier integral operator manipulations.
\end{rema}

\begin{prop}\label{prop:Cauchy_product}
Let $(a_j)_{j\in \N}$, $(b_j)_{j\in \N}$ be two formal analytic amplitudes on an open cone $\Omega$. Let $a,b$ be respective analytic realisations of $(a_j)_{j\in \N},(b_j)_{j\in \N}$. Then $ab$ is an analytic amplitude, which realises the Cauchy product $((a*b)_j)_{j\in \N}=(\sum_{\ell=0}^ja_{\ell}b_{j-\ell})_{j\in \N}$.
\end{prop}
\begin{proof}
One can directly check that $ab$ is an analytic amplitude and that\linebreak $((a*b)_j)_{j\in \N}$ is a formal analytic amplitude. Now (counting some terms twice)
\begin{multline*}
\left|a(x,\theta)b(x,\theta)-\sum_{k=0}^{N-1}(a*b)_k(x,\theta)\right|\\ \leq
\left|a(x,\theta)b(x,\theta)-\sum_{j=0}^{\frac{N-1}{2}}\sum_{\ell=0}^{\frac{N-1}{2}}a_j(x,\theta)b_\ell(x,\theta)\right|+\sum_{j+\ell=\frac{N-1}{2}}^{N-1}
\left|a_j(x,\theta)b_{\ell}(x,\theta)\right|.
\end{multline*}
The first term of the right-hand side is equal to
\[
a(x,\theta)\left(b(x,\theta)-\sum_{\ell=0}^{\frac{N-1}{2}}b_\ell(x,\theta)\right)+\sum_{\ell=0}^{\frac{N-1}{2}}b_{\ell}(x,\theta)\left(a(x,\theta)-\sum_{j=0}^{\frac{N-1}{2}}a_j(x,\theta)\right),
\]
so that
\[
\left|a(x,\theta)b(x,\theta)-\sum_{j=0}^{\frac{N-1}{2}}\sum_{\ell=0}^{\frac{N-1}{2}}a_j(x,\theta)b_\ell(x,\theta)\right|\leq C|\theta|^{d_a+d_b-\frac{N-1}{2}}\rho^{\frac{N-1}{2}}\left(\frac{N-1}{2}\right)!.
\]

To treat the second term, observe that, given $j+\ell\in [(N-1)/2,N-1]$,
\begin{multline}\label{eq:ajbl}
|a_j(x,\theta)b_{\ell}(x,\theta)|\leq C_{a,b}|\theta|^{d_a+d_b-j-\ell}\rho^{j+\ell}(j+\ell)!\\ \leq C_{a,b}|\theta|^{d_a+d_b}\max\left(|\theta|^{-N+1}\rho^{N-1}(N-1)!,\;|\theta|^{-\frac{N-1}{2}}\rho^{\frac{N-1}{2}}\left(\frac{N-1}{2}\right)!\right),
\end{multline}
where in the last inequality we used the fact that the middle quantity is a log-convex function of $j+\ell$. Altogether if $|\theta|^{-\frac{N-1}{2}}\rho^{\frac{N-1}{2}}(\frac{N-1}{2})!\leq |\theta|^{-N+1}\rho^{N-1}(N-1)!$, the proof is complete.

Similarly, one can write
\begin{multline*}
\left|a(x,\theta)b(x,\theta)-\sum_{k=0}^{N-1}(a*b)_k(x,\theta)\right|\\ \leq
\left|a(x,\theta)b(x,\theta)-\sum_{j=0}^{N-1}\sum_{\ell=0}^{N-1}a_j(x,\theta)b_\ell(x,\theta)\right|+\sum_{j+\ell=N-1}^{2N-2}|a_j(x,\theta)b_{\ell}(x,\theta)|,
\end{multline*}
and if $|\theta|^{-2N+2}\rho^{2N-2}(2N-2)!\leq |\theta|^{-N+1}\rho^{N-1}(N-1)!$, we can conclude using~\eqref{eq:ajbl} again.

The sequence $(v_k)_{k\in \N}=(|\theta|^{-k}\rho^kk!)_{k\in \N}$ is decreasing on $\{k\leq |\theta|/\rho\}$ and increasing on $\{k\geq |\theta|/\rho\}$. Therefore if both $v_{\frac{N-1}{2}}$ and $v_{2N-2}$ are greater than $v_{N-1}$, then $N-1\in [\frac{|\theta|}{2\rho},2\frac{|\theta|}{\rho}]$. Under these circumstances, by the Stirling formula, all of $v_{\frac{N-1}{2}},v_{N-1},v_{2N-2}$ can be bounded by $e^{-c|\theta|}$ for some $c>0$, and for $\rho'$ large enough, $e^{-c|\theta|}\leq (\rho')^{N-1}|\theta|^{-N+1}(N-1)!$; this concludes the proof of Proposition~\ref{prop:Cauchy_product}.
\end{proof}

\begin{prop}\label{prop:product_preparation}
Let $\Omega_1\subset \R^{n_x}\times \R^{n_{\theta}}\times
\R^{n_y}$ and $\Omega_2\subset \R^{n_y}\times \R^{n_{\tau}}\times
\R^{n_z}$ be open cones with respect to their middle variables. Let $c>0$ and let
\begin{equation*}
\Omega=\left\{
\begin{aligned}
&(x,\theta,y,\tau,z)\in \R^{n_x}\times \R^{n_{\theta}}\times
\R^{n_y}\times \R^{n_{\tau}}\times
\R^{n_z},\\ &(x,\theta,y)\in \Omega_1,(y,\tau,z)\in \Omega_2,\quad c|\theta|<|\tau|<c^{-1}|\theta|
\end{aligned}
\right\}.
\end{equation*}
Let $a_1:\Omega_1\to \R$ and $a_2:\Omega_2\to \R$ be analytic amplitudes (real-analytic with respect to their first and third variables). Then
\[
A:\Omega\ni (x,\theta,y,\tau,z)\mapsto a_1(x,\theta,y)a_2(y,\theta,z)
\]
is an analytic amplitude (real-analytic with respect to its first and last variable). Moreover, if $a_1$ and $a_2$ are respectively realisations of formal analytic amplitudes $(a_{1;j})_{j\in
\N},(a_{2;j})_{j\in \N}$, then $A$ is a realisation of the Cauchy product
\[
\left(\sum_{\ell=0}^ja_{1;\ell}(x,\theta,y)a_{2;j-\ell}(y,\tau,z)\right)_{j\in
\N},
\]
the latter being a formal analytic amplitude on $\Omega$.
\end{prop}
\begin{proof}
The key property is that, on $\Omega$, the quantities $|\theta|,|\tau|$ and $(|\tau|^2+|\theta|^2)^{\frac 12}$ are comparable. Therefore $(x,\theta,y,\tau,z)\mapsto a_1(x,\theta,y)$ is an analytic amplitude on $\Omega$, which realises the formal analytic amplitude $((x,\theta,y,\tau,z)\mapsto a_{1;j}(x,\theta,y))_{j\in \N}$, and a similar property holds for $a_2$. By Proposition~\ref{prop:Cauchy_product}, the proof is complete.
\end{proof}

We conclude this section with a practical example, related to Proposition~\ref{prop:flat_Grauert}, and which will allow us to prove that the Szeg\H{o} kernel takes the form of an analytic Fourier integral operator, in the sense of Section~\ref{sec:four-integr-oper}.

\begin{prop}\label{prop:m_n_amplitude}
Define the following open cone in $\C^d$:
\[
\Omega=\left\{z\in \C^d,\; \lvert\im(z)|^2<\lvert\re(z)|^2\right \};
\]
recall that the function $r:\Omega\ni z\mapsto
\sqrt{\sum_{j=1}^dz_j^2}$ is well-defined and holomorphic.

Then
\[
a:z\mapsto e^{-r(z)}\int_{S^{n-1}}e^{z\cdot y}\dd y
\]
is an analytic amplitude which realises the formal analytic amplitude
\[
\left(2^{\frac{d-1}{2}}\left|S^{d-2}\right|\frac{(d-1)(d-3)\dots (d-2j+1)}{2^jj!}\Gamma\left(\frac{d+1}{2}+j\right)r(z)^{-\frac{d+1}{2}-j}\right)_{j\in \N}.
\]
\end{prop}
\begin{proof}
Let us first study the elementary properties of $a$ on $\Omega$.

$S^{d-1}$ is a real contour in the complex manifold $\widetilde{S^{d-1}}:=\{\omega\in \C^d,
\sum_j \omega_j^2=1\}$, on which $\omega\mapsto e^{z\cdot \omega}$ is holomorphic. This manifold is preserved by the linear action of
\[
\widetilde{O(d)}:=\left\{M\in M_d(\C),\; M^TM=I\right\},
\]
whose action is transitive; more generally, for every $\lambda\in
\C\setminus \R^-$, $\widetilde{O(d)}$ acts transitively on $\{z\in \Omega_1, r(z)^2=\lambda\}$. The function $r$ is itself $\widetilde{O(d)}$-invariant.


To prove that $a$ itself is $\widetilde{O(d)}$-invariant, let $z\in \Omega$. By the previous considerations, there exists $M\in \widetilde{O(d)}$ such that $Mz$ is of the form $\alpha x$ where $x\in \R^{d}$ and $\re(\alpha)>0$. Since $\widetilde{O(d)}$ is connected, one can deform the contour of integration into $\{M^{-1}\omega,\linebreak \omega\in S^{d-1}\}$, and then $z\cdot M^{-1}\omega=\alpha z\cdot \omega$. Thus $a(z)=a(Mz)$, so $a$ is $\widetilde{O(d)}$-invariant.

Using this $\widetilde{O(d)}$ invariance, to prove that $a$ is the realisation of a formal analytic amplitude, it suffices to prove the required estimate for $a(R\alpha e_1)$, where $R\to +\infty$, $|\alpha|=1$, $\re(\alpha)>\lvert\im(\alpha)|$, and $e_1$ is the first vector of the base.

One has
\[
a(R\alpha e_1)=\left|S^{d-2}\right|\int_{-1}^1 e^{-R\alpha (x-1)}\left(1-x^2\right)^{\frac{d-1}{2}}\dd x,
\]
where $|S^{d-2}|$ is the area of the unit sphere in $\R^{d-1}$. Letting $x=1-y$, we rewrite this integral into
\[
a(R\alpha e_1)=2^{\frac{d-1}{2}}|S_{d-2}|\int_{0}^{2} e^{\alpha R y}y^{\frac{d-1}{2}}\left(1-\frac{y}{2}\right)^{\frac{d-1}{2}}\dd y.
\]
Let us decompose this integral into two parts. If $y\in [1,2]$, then $|e^{-\alpha Ry}|=e^{-\re(\alpha)Ry}\geq e^{-\frac{R}{\sqrt{2}}}$ is exponentially small. On $[0,1]$, we will replace $(1-\frac{y}{2})^{\frac{d-1}{2}}$ by its Taylor series at $0$. First
\[
\frac{\partial^j}{\partial y^j}\left(1-\frac{y}{2}\right)^{\frac{d-1}{2}}=\frac{(d-1)(d-3)\dots(d-2j+1)}{2^j}\left(1-\frac{y}{2}\right)^{\frac{d-1}{2}-j}.
\]
Let $c$ be a small constant (to be determined later on), and let $J_R=\lfloor cR\rfloor$. We write, for $y\in [0,1]$,
\[
\left(1-\frac{y}{2}\right)^{\frac{d-1}{2}}=\sum_{j=0}^{J_R}\frac{(d-1)(d-3)\dots(d-2j+1)}{2^jj!}y^j+y^{J_R+1}F_{J_R}(y),
\]
where
\[
\sup_{y\in [0,1]}|F_{J_R}(y)|\leq
\frac{(d-1)(d-3)\dots(d-2J_R+1)}{4^{J_R}J_R!}.
\]
Notice that the sequence
\[
\left(\frac{(d-1)(d-3)\dots(d-2j+1)}{2^jj!}\right)_{j\in \N}
\]
is bounded; for every $j\geq d/2$, from $j$ to $j+1$ the general term is multiplied by $(d-2j-1)(j+1)\in [-1,0]$.

For every $k>0$, $y>0$ and $|\alpha|=1$ with $\re(\alpha)>\lvert\im(\alpha)|,$
\[
\left|e^{-\alpha y}y^k\right|\leq e^{-\re(\alpha)y}y^k\leq e^{-\frac{y}{\sqrt{2}}}y^k\leq k^ke^{k\left(\log(\sqrt{2})-1\right)}\leq k^k.
\]
We apply this inequality at $k=J_R+1+\frac{d-1}{2}$, and obtain
\[
\left|\int_0^1e^{\alpha R y}y^{\frac{d-1}{2}+J_R}F_{J_R}(y)\dd y\right|\leq (cR)^{cR}R^{-cR}\leq Ce^{-c'R}
\]
for some $c'>0$, if $c$ is small enough. It remains
\[
2^{\frac{d-1}{2}}\left|S^{d-2}\right|\sum_{j=0}^{J_R}\frac{(d-1)(d-3)\dots(d-2j+1)}{2^jj!}\int_0^1e^{\alpha R y}y^{\frac{d-1}{2}+j}\dd y.
\]
And to this end we use that, for every $k\leq \frac{R}{\sqrt{2}}$,
\[
\sup_{y\in [1,+\infty)}\left|e^{-\alpha Ry}y^k\right|\leq e^{-R/\sqrt{2}};
\]
applied at $k=j+\frac{d-1}{2}+2$, we obtain
\[
\left|\sum_0^{J_R}\frac{(d-1)(d-3)\dots(d-2j+1)}{2^jj!}\int_1^{+\infty} e^{\alpha R y}y^{\frac{d-1}{2}+j}\dd y\right|\leq Ce^{-c'R}.
\]
Now we can conclude:
\begin{align*}
\left|
\begin{multlined}
a(R\alpha e_1)-\sum_{j=0}^{J_R}2^{\frac{d-1}{2}}\left|S^{d-2}\right|\frac{(d-1)(d-3)\dots (d-2j+1)}{2^jj!}\\
\Gamma\left(\frac{d+1}{2}+j\right)(R\alpha)^{-\frac{d+1}{2}-j}
\end{multlined}
\right|
\leq Ce^{-c'R}
\end{align*}
and by $\widetilde{O(d)}$-invariance, for every $z\in \Omega$,
\begin{equation*}
\left|
\begin{multlined}
a(z)-\sum_{j=0}^{J_R}2^{\frac{d-1}{2}}\left|S^{d-2}\right|\frac{(d-1)(d-3)\dots(d-2j+1)}{2^jj!}\\
\Gamma\left(\frac{d+1}{2}+j\right)r(z)^{-\frac{d+1}{2}-j}
\end{multlined}
\right|
\leq Ce^{-c'|z|}.
\end{equation*}
Using the Cauchy formula, we deduce a similar formula for all derivatives of $a$, which concludes the proof that $a$ satisfies the derivative bounds of an analytic amplitude.
\end{proof}

\subsection{Stationary phase}\label{sec:stationary-phase}
The goal of this section is to obtain a stationary phase theorem adapted to our notion of analytic amplitudes. We will have to deal with complex-valued phase functions, and our first step is a study of the contour deformations necessary to reach the critical points.

\begin{prop}\label{prop:contour}
Let $\clK\Subset \R^n_x\times \R^m_y$. Let $\phi$ be a real-analytic function from an open neighbourhood of $\clK$ to $\{z\in \C,\im(z)\geq 0\}$. Suppose that $\nabla^2_y\phi$, the matrix of second derivatives with respect to the second set of variables, is invertible everywhere on $\clK$ and that $\im(\phi)>0$ on $\partial K$.

Then there exists a constant $c>0$ and a covering of $\clK$ by two open sets $U_1$ and $U_2$ of $\R^{n+m}$ such that the following is true.
\begin{enumerate}
\item\label{prop3.10.1} There exists a real-analytic map $\kappa_1:U_1\to \R^n_x\times \C^m_y$ with the following properties:
\begin{itemize}
\item $\kappa_1$ acts as the identity on the first variable.
\item $\kappa_1$ is a diffeomorphism between $U_1$ and its image $\Gamma_1\subset \R^n_x\times \{\lvert\im(y)|<c\}$.
\item $\Gamma_1$ is a totally real submanifold of $\C^n\times \C^m$.
\item $\im(\phi)\geq c$ on $\partial \Gamma_1$.
\item $\kappa_1$ is isotopic to the identity among the maps satisfying the four previous properties. (We will henceforth call this isotopy a \emph{contour deformation}, and call $\Gamma_1$ a \emph{contour}.)
\item For every $x\in \R^n$, in each connected component of $\Gamma_1(y):=\{y\in \C^m,\linebreak (x,y)\in \Gamma_1\}$, the holomorphic extension $\widetilde{\phi}$ of $\phi$ admits exactly one critical point $y^*$ with respect to $x$ and the real part of $\widetilde{\phi}$ is constant.
\item $\nabla^2_y(\im\widetilde{\phi}\circ \kappa_1)\geq c$, and the phase, at the critical points, has non-negative imaginary part.
\end{itemize}
\item\label{prop3.10.2} There exists a contour $\Gamma_2$ deforming $U_2$, on which $\im(\widetilde{\phi})\geq c$.
\end{enumerate}
\end{prop}
\begin{proof}
Let $V$ be a neighbourhood of $\clK$ on which $\phi$ is defined and $\det\nabla^2_y\phi$ is bounded away from $0$. Let $\epsilon>0$ be small enough. We define
\begin{align*}
&U_1=\left\{(x,y)\in V, |\nabla_y\phi(x,y)|< \epsilon \right\},& &U_2= \left\{(x,y)\in V, |\nabla_y\phi(x,y)|> \epsilon/2\right\}.
\end{align*}
Since $\det\nabla^2_y\phi$ is bounded away from $0$, for every $(x,y)\in U_1$, there exists a unique critical point $(x,y^*(c))$ of $\phi$ closeby (at distance of order $\epsilon$). Moreover, again because of the nondegeneracy, a small neighbourhood of $(x,y^*(x))$ contains the connected component $V$ of $U_1\cap (\{x\}\times \R^m)$ to which $(x,y)$ belongs.

On this neighbourhood of the critical point, we apply the holomorphic Morse lemma: there exists a biholomorphism $\sigma$ (with real-analytic dependence on $x$) such that $\widetilde{\phi}(x,y)=\sigma(x,y)^2+\widetilde{\phi}(x,y^*(x))$.

We are able to deform $V$ into a neighbourhood of $0$ in $\{\sigma\in i\R^n\}$, first by straightening $\sigma(V)$ into its tangent space at some point $y_0$, then by translation and rotation, since the set of the totally real linear subspaces is open and connected.

On $U_2$, where $|\nabla_y\phi(x,y)|>\epsilon/2$, we deform using the ascending gradient flow of $\im(\widetilde{\phi})$ on $\C^m$. After a time $\epsilon$, the imaginary part of the phase is incremented by a quantity bounded below everywhere except in the vicinity of real points where $\nabla_y\re(\widetilde{\phi})$ is small. This can only happen if $\nabla_y\im(\widetilde{\phi})$ is large, which only happens if $\im(\phi)$ is large to begin with. Thus, after this gradient flow, $\im(\widetilde{\phi})$ is bounded from below everywhere.

To conclude the proof, it remains to show that, among the critical points sufficiently close to $U_1$, the imaginary part of the phase is never negative. We proceed by contradiction and suppose that there is such a critical point close to a point $P\in U_1$. Up to a linear transformation by an element of $GL_m(\R)$, near $P$, the second derivative of the phase with respect to the second variable $y_1$ is nonzero. We can therefore apply the analytic Morse lemma in the variable $y_1$, with all other variables as parameters. This transforms $U_1$ into a piece of real curve in $\C$, parametrised by the other variables. By hypothesis, the imaginary part of the phase is strictly smaller on the image of the real locus by the Morse biholomorphism than at the critical point. Hence, this piece of curve belongs to either one of the two quarter spaces in $\C$ where $\re(y_1^2)>0$. Since these quarter spaces have a corner at $0$, the image of $U_1$ cannot get close to $0$ without the second derivative of the phase along $U_1$ getting large. This is not possible because $\clK$ is compact.
\end{proof}

The contour deformation above allows us to prove a stationary phase result. Mind that there is no assumption of uniqueness of the critical point, so that the result of the stationary phase is a sum of different contributions coming from the different critical points.

\begin{prop}\label{prop:stat_phase}
Let $\Omega_1$ be an open subset of $\R^{n_{x}}\times
\R^{n_{\theta}}$, conic with respect to the second variable and relatively compact, and let $\Omega_2$ be an open, relatively compact subset of $\R^{n_y}$. Let $\phi:\Omega_1\times \Omega_2\to \C$ be real-analytic and one-homogeneous with respect to $\theta$, which extends to a small neighbourhood of $\Omega_1\times \Omega_2$. Suppose the following:
\begin{itemize}
\item $\phi$ takes values in $\{{\rm Im(z)}\geq 0\}$.
\item $\nabla^2_y\phi$ is nonsingular.
\item $\im(\phi)>0$ on $\partial(\Omega_1\times \Omega_2)$.
\end{itemize}

Let $a:\Omega_1\times \Omega_2\to \C$ be such that $((x,y),\theta)\mapsto a(x,\theta,y)$ is an analytic amplitude.

Then there exists a finite collection of open conic subsets $V_1,\dots,V_N$ of $\Omega$, real-analytic zero-homogeneous functions $y^*_j:V_j\to
\widetilde{\Omega_2}$, analytic amplitudes $b_j$ on $V_j$, and $\epsilon>0$, such that, for every $(x,\theta)$ in a small complex neighbourhood of $\Omega$ with respect of the first variable,
\[
\int_{\Omega_2}e^{i\phi(x,\theta,y)}a(x,\theta,y)\dd y -
\sum_{j=1}^N\1_{(x,\theta)\in V_i}e^{i\widetilde{\phi}\left(x,\theta,y^*_j(x,\theta)\right)}b_j(x,\theta)=O(e^{-\epsilon|\theta|}).
\]
Moreover $\im(\phi)(x,\theta,y^*_j(x,\theta))\geq 0$ on $V_j$ and $\im(\phi)(x,\theta,y^*_j(x,\theta))>\epsilon|\theta|$ on $\partial V_j$.


If $a$ realises a formal analytic amplitude, then the $b_j$ realise the formal amplitudes obtained by formal stationary phase (and in particular, said formal amplitudes are analytic).
\end{prop}
\begin{proof}
We first apply the contour deformations of Proposition~\ref{prop:contour}, which allows us to identify the open sets $V_j$ and the functions $y^*_j$. Recalling the open sets $U_1$ and $U_2$ of $\Omega_1\times \Omega_2$, for each $(x,\theta)\in \Omega_1$, each connected component of $\{y\in \Omega_2,(x,\theta)\in \Omega_1\}$ will correspond to a single critical point of $\phi$ with respect to $y$. As $x,\theta$ vary, these connected components never merge or split (as $\lvert\det(\nabla^2_y\phi)|$ is bounded from below, critical points are far away from each other) but simply appear or disappear, as the associated critical point moves too far away from the real locus.

Therefore, if we define the sets $V_j$'s as the projections onto the $(x,\theta)$ variables of the connected components of $U_1$, and $y^*_j$ as the corresponding critical points, one has indeed $\im(\phi)(x,\theta,y_j^*(x,\theta))\geq 0$, and positivity on the boundary.

Let $(x,\theta)\in \Omega_1$; decompose $\Omega_2=U_2(x,\theta)\cup
\bigcup_{j\in \clJ(x,\theta)}U_{1;j}(x,\theta)$, where $U_2(x,\theta)=\{y\in \Omega_2,(x,\theta,y)\in U_2\}$, $\clJ(x,\theta)=\{1\leq j\leq n, (x,\theta)\in V_j\}$ and $U_{1;j}(x,\theta)$ is the connected component of $\{y\in
\Omega_2,(x,\theta,y)\in U_2\}$ whose complex neighbourhood contains $y_j^*(x,\theta)$. Since there exists $\epsilon>0$ such that $\im(\phi)>\epsilon|\theta|$ on $U_2(x,\theta)$, we are left with
\[
\sum_{j\in
\clJ(x,\theta)}\int_{U_{1;j}(x,\theta)}e^{i\phi(x,\theta,y)}a(x,\theta,y)\dd y.
\]
By Proposition~\ref{prop:contour}, there exists a homotopy $\Gamma_t$ of real-analytic contours deforming $U_{1;j}(x,\theta):=\Gamma_{0;j}(x,\theta)$ into some $\Gamma_{1;j}(x,\theta)$ containing $y^*_j(x,\theta)$ on which $\re(\phi)$ is constant; this homotopy satisfies $\im(\phi)>\epsilon|\theta|$ on $\partial
\Gamma_{t;j}(x,\theta)$.

Extending $a$ and $\phi$ holomorphically on $y$, Stokes' theorem gives
\[
\int_{\Gamma_{0;j}(x,\theta)}e^{i\widetilde{\phi}(x,\theta,y)}a(x,\theta,y)\dd y -
\int_{\Gamma_{1;j}(x,\theta)}e^{i\widetilde{\phi}(x,\theta,y)}\widetilde{a}(x,\theta,y)\dd y=\int_{W_1}e^{i\widetilde{\phi}(x,\theta,y)}\widetilde{a}(x,\theta,y)\dd y;
\]
here, on $W_1=\bigcup_{t\in [0,1]}\partial \Gamma_{t;j}(x,\theta)$, one has $\im(\widetilde{\phi})>\epsilon |\theta|$. Hence, the difference between the integrals on the two contours is $O(e^{-c|\theta|})$.

Applying the Morse lemma on $\widetilde{\phi}$ on $\Gamma_{1;j}$, and an analytic change of variables on $y$, it remains to prove that, when $a$ is an analytic amplitude of $((x,y),\theta)$, then so is
\begin{equation}\label{eq:std_phase}
(x,\theta)\mapsto \int_{B(0,1)}e^{-|\theta|y^2}a(x,\theta,y)\dd y,
\end{equation}
and that this stationary phase commutes with the realisation of formal analytic amplitudes.

We first show a result ``at fixed order'', then re-sum the obtained amplitude in the spirit of Proposition~\ref{prop:real_formal_symb}.

Let $N\in \N$ and let $u\in C^N(B_{\R^d}(0,1),\R)$. The Taylor expansion of $u$ at $0$ is
\[
u(y)=\sum_{|\mu|<N}\frac{\partial^{\mu}u(0)y^{\mu}}{\mu!}+R(y),
\]
where
\[
|R(y)|\leq \frac{\rho_d^N|y|^N}{N!}\left\|\nabla^Nu\right\|_{L^{\infty}};
\]
here $\rho_d$ depends only on the dimension $d$.

Since $\frac{|y|^N}{N!}\leq e^{|y|}$, we obtain
\[
\left| \int_{B(0,1)}e^{-\lambda y^2}R(y)\dd y\right|\leq C_d\rho_d^N\lambda^{-\frac{N}{2}} \left(\frac{N}{2}\right)!\frac{\left\|\nabla^Nu\right\|_{L^{\infty}}}{N!}.
\]
Noting that, on $\R^+$, the function $y\mapsto y^Ne^{-y}$ reaches its maximum at $y=N$, then as long as $\frac{N+d+1}{2}<\lambda$,
\begin{align*}
y\notin B(0,1)\to |y^{\mu}e^{-\lambda y^2}| &\leq e^{-\lambda}|y|^{-d-1};
\\
\intertext{hence}
\left|\int_{\R^d\setminus B(0,1)}e^{-\lambda y^2}\sum_{|\mu|<N}\frac{\partial^\mu u(0)y^{\mu}}{\mu!}\dd y\right| &\leq C_de^{-\lambda}\sum_{j=0}^{N-1}\frac{\rho_d^j\left\|\nabla^ju\right\|_{L^{\infty}}}{j!}.
\end{align*}
It remains to compute the explicit integral
\[
\int_{\R^d}e^{-\lambda y^2}\sum_{|\mu|<N}\frac{\partial^\mu u(0)}{\mu!}y^{\mu}\dd y=\pi^{\frac d2}\lambda^{-\frac d2}\sum_{j<\frac{N}{2}}\frac{\Delta^ju(0)}{(4\lambda)^jj!}.
\]
Altogether we have proved, for some $\rho_d>0$, for every $N<2\lambda-d-1$, for every $u\in C^N$,
\begin{multline*}
\left|\int e^{-\lambda y^2}u(y)\dd y-\pi^{\frac d2}\lambda^{-\frac d2}\sum_{j<
\frac{N}{2}}\frac{\Delta^ju(0)}{(4\lambda)^jj!}\right|\\
\leq C_de^{-\lambda}\sum_{j=0}^{N-1}\frac{\rho_d^j\left\|\nabla^ju\right\|_{L^{\infty}}}{j!}+C_d\rho_d^N\lambda^{-\frac N2}\left(\frac N2\right)!\frac{\left\|\nabla^Nu\right\|_{L^{\infty}}}{N!}.
\end{multline*}
If, on the contrary, $N\geq 2\lambda-d-1$, then
\[
\left|\int e^{-\lambda y^2}u(y)\dd y\right|\leq C_0\|u\|_{L^{\infty}}\leq C_0\lambda^{-\frac N2}\left(\frac N2\right)!\rho_d^{\frac N2}\|u\|_{L^{\infty}}
\]
and, for $j<\frac N2$,
\[
\lambda^{-\frac d2}\frac{\Delta^ju(0)}{(4\lambda)^jj!}\leq
\lambda^{-\frac N2}\left(\frac N2\right)!\rho_d^{\frac N2}\frac{\left\|\nabla^{2j}u\right\|_{L^{\infty}}}{(2j)!}.
\]
Therefore, without condition on $N$ and $\lambda$, for every $u\in C^N$,
\begin{multline}\label{eq:stat_phase_fixed_order}
\left|\int e^{-\lambda y^2}u(y)\dd y-\pi^{\frac d2}\lambda^{-\frac d2}\sum_{j<
\frac{N}{2}}\frac{\Delta^ju(0)}{(4\lambda)^jj!}\right|\\
\leq C_d\left(e^{-\lambda}+\rho_d^N\lambda^{-\frac N2}\left(\frac N2\right)!\right)\sum_{j=0}^N\frac{\left\|\nabla^ju\right\|_{L^{\infty}}}{j!}.
\end{multline}

Plugging $u:y\mapsto a(x,\lambda \omega,y)$, we obtain the second part of the result.

It remains to prove that~\eqref{eq:std_phase} defines an analytic amplitude. The derivatives on $x$ and on the spherical variables of $\theta$ play no role, and we are left with proving that, for some $\epsilon>0$ and $\rho>0$,
\[
\forall k\leq \epsilon \lambda,\;\left|\frac{\partial^k}{\partial
\lambda^k}\int_{B(0,1)}e^{-\lambda y^2}a(x,y,\lambda)\dd y
\right|\leq \rho^{k}k!\lambda^{d-\frac{n_y}{2}-k}.
\]
One term of the binomial sum looks like
\[
\binom{k}{j}\int_{B(0,1)}e^{-\lambda y^2}y^{2j}\partial_\lambda^{k-j}a(x,y,\lambda)\dd y
\]
and, applying~\eqref{eq:stat_phase_fixed_order} at order $N=2j+n_y$, we obtain
\begin{multline*}
\left|\int_{B(0,1)}e^{-\lambda y^2}y^{2j}\partial_\lambda^{k-j}a(x,y,\lambda)\dd y\right|\\ \leq C_d\left(e^{-\lambda}\sum_{\ell=0}^{N-1}\frac{\rho_d^\ell\left\|\nabla^\ell_y\partial^{k-j}_{\lambda}(y^{2j}a)\right\|_{L^{\infty}}}{\ell!}+\lambda^{-\frac{n_y}{2}}\sum_{\ell=j}^{j+\frac{n_y}{2}}\frac{\Delta^\ell_y\left(y^{2j}\partial_{\lambda}^{k-j}a\right)(0)}{(4\lambda)^\ell \ell!}\right.\\
+\left.\rho_d^{2j+n_y}\lambda^{-\frac{n_y}{2}-j}\frac{\left\|\nabla^{2j+n_y}_yy^{2j}\partial^{k-j}_{\lambda}a\right\|_{L^{\infty}}}{(2j+n_y)!}\right).
\end{multline*}
As long as $\frac{k}{\lambda}$ is small enough, if $a$ is an analytic amplitude, then this satisfies the required claim.
\end{proof}


\subsection{Positive nondegenerate phases and their Lagrangians}\label{sec:posit-nond-phas}

In this short subsection, we avail ourselves of a few geometrical results concerning what will be the oscillating phases of the Fourier {integral operators} considered in this article. We will directly deal with \emph{complex-valued} phase functions; real-analytic regularity makes their treatment quite transparent, and every object will be defined in a complex-geometric setting. For instance, ``Lagrangians'' will always be complex submanifolds (of a larger complex symplectic space) on which the symplectic form vanishes.

We will only use the relatively standard notion of non-degenerate phase functions.
\begin{defi}\label{def:good_phase}
Let $U\subset \R^{n_x}$, $W\subset
\R^{n_y}$ be relatively compact open sets and $V\subset
\R^{n_{\theta}}$ an open cone.

A (strictly) \emph{positive non-degenerate} phase function $\phi:U_x\times V_{\theta}\times W_y\to \C$ is a one-homogeneous function of $\theta$, real-analytic o $\{\theta\neq 0\}$, such that, with $\Sigma_{\phi}=\{(x,\theta,y)\in \widetilde{U}\times
\widetilde{V}\times \widetilde{W},
\nabla_{\theta}\widetilde{\phi}=0\}$, the following conditions apply:
\begin{enumerate}
\item\label{defi3.12.1} The gradients $\nabla_x\widetilde{\phi}$ and $\nabla_y\widetilde{\phi}$, restricted to $\Sigma_{\phi}$, do not vanish.
\item\label{defi3.12.2} $\nabla_{\theta}\widetilde{\phi}$ is a defining function for $\Sigma_{\phi}$ (that is to say, $\nabla_{x,\theta,y}\nabla_{\theta}\widetilde{\phi}$ has rank $n_{\theta}$ and its columns span the conormal bundle of $\Sigma_{\phi}$).
\item\label{defi3.12.3} There exists $c>0$ such that, on the real set $U\times V\times W$, one has $\im(\phi)\geq c|\theta||\nabla_{\theta}\phi|^2$.
\end{enumerate}
\end{defi}

\begin{prop}\label{prop:lagr}
Let $\phi$ be a positive non-degenerate phase function. Then $\Sigma_{\phi}$ is a real-analytic conic submanifold diffeomorphic to the conical Lagrangian
\[
\Lambda_{\phi}=\left\{\left(x,\dd_x\widetilde\phi(x,\theta,y),y,-\dd_y\widetilde\phi(x,\theta,y)\right),\;(x,\theta,y)\in
\Sigma_{\phi}\right\}\subset (T^*\C^n\setminus \{0\})^2,
\]
called the canonical relation of $\phi$.
\end{prop}
\begin{proof}
This is an elementary and well-known fact, whose proof is the same whenever $\phi$ is a real-valued phase function, a $C^{\infty}$ complex-valued phase function with almost analytic extensions, or a real-analytic complex-valued phase function. See for instance~\cite[Proposition~25.4.4]{hormander_analysis_1985}.
\end{proof}

Definition~\ref{def:good_phase} is independent of a real-analytic change of variables in $\theta$, as long as it has a real-analytic dependence in $x$ and $y$. This motivates the next definition.
\begin{defi}\label{def:good_contour}
Given a real-analytic function $\phi:U\times V\times W\to \C$ satisfying the first two conditions of Definition~\ref{def:good_phase}, a totally real, conical submanifold of $\widetilde{V}$ with analytic dependence in $x$ and $y$ will be called a \emph{good contour} if $\widetilde{\phi}$, restricted to this contour, satisfies the third condition.
\end{defi}


\begin{defi}\label{def:composition}
Two Lagrangians $\Lambda_1\subset (T^*X\setminus \{0\})\times (T^*Y\setminus \{0\})$ and $\Lambda_2\subset (T^*Y\setminus \{0\})\times (T^*Z\setminus \{0\})$ are said to have \emph{transverse composition} at a point of $\Lambda_1\times \Lambda_2\cap [(T^*X\setminus \{0\})\times
\Delta(T^*Y\setminus \{0\})\times (T^*Z\setminus \{0\})]$ if the intersection of these two submanifolds is transverse at this point.
\end{defi}

\begin{prop}\label{prop:compo_phase}
Let $\phi_1:\Omega_1\to \C$ and $\phi_2:\Omega_2$ be positive nondegenerate phase functions whose Lagrangians have transverse composition. Then for every $\epsilon>0$, $\Phi: (x,z,\theta,y,\tau)\mapsto \phi_1(x,\theta,y)+\phi_2(y,\tau,z)$ is a positive nondegenerate phase function on
\[
\left\{(x,\theta,y)\in \Omega_1,\quad(y,\tau,z)\in
\Omega_2,\quad\epsilon|\theta|<|\tau|<\epsilon^{-1}|\theta|\right\},
\]
up to a small contour homotopy in $y$. Its Lagrangian is
\begin{equation*}
\Lambda_1\circ \Lambda_2:=\left\{
\begin{aligned}
&(x,\xi,z,\zeta)\in T^*X\setminus
\{0\}\times T^*Z\setminus \{0\},\\
&\exists (y,\eta)\in T^*Y\setminus
\{0\}, \quad(x,\xi,y,\eta)\in \Lambda_1,\quad(y,\eta,z,\zeta)\in
\Lambda_2
\end{aligned}
\right\},
\end{equation*}
and its real locus is exactly
\begin{equation*}
(\Lambda_1\circ \Lambda_2)_\R=\left\{
\begin{aligned}
&(x,\xi,z,\zeta)\in (T^*X)_{\R}\setminus
\{0\}\times (T^*Z)_{\R}\setminus \{0\},\\
&\exists (y,\eta)\in (T^*Y)_{\R}\setminus
\{0\}, \quad(x,\xi,y,\eta)\in (\Lambda_1)_{\R},\quad(y,\eta,z,\zeta)\in (\Lambda_2)_{\R}
\end{aligned}
\right\}.
\end{equation*}
\end{prop}
\begin{proof}
See~\cite[Proposition~25.5.4]{hormander_analysis_1985}. Note that the notion of positivity used there is more general than ours and we have to prove positivity by hand; a priori we only have
\[
\im(\Phi)\geq c|\theta||\nabla_{\theta}\Phi|^2+c|\tau||\nabla_{\tau}\Phi|^2.
\]
We now apply a contour deformation on $y$, consisting in the positive gradient flow of $(|\theta|^2+|\tau|^2)^{-\frac{1}{2}}\im(\Phi)$. Along this flow, $\im(\Phi)$ increases at a rate $(|\theta|^2+|\tau|^2)^{-\frac{1}{2}}|\nabla_y\Phi|^2$. Therefore, after a small time of evolution, we obtain
\[
\im(\Phi)\geq c|\theta||\nabla_{\theta}\Phi|^2+c|\tau||\nabla_{\tau}\Phi|^2+c(|\theta|^2+|\tau|^2)^{-\frac{1}{2}}|\nabla_y\Phi|^2.
\]
This is the desired result since $|\theta|, |\tau|,$ and $(|\theta|^2+|\tau|^2)^{\frac 12}$, are comparable.

The condition on the real locus is an easy consequence of the fact that the imaginary part of the phases are always nonpositive, and vanish exactly on the real locus of their Lagrangians.
\end{proof}

Phase functions are seldom more than continuous at the origin (unless they are linear), so that it is necessary to remove conical neighbourhoods of $\theta=0$ and $\tau=0$ above. When actually composing Fourier {Integral} operators in the next section, it will be necessary to deal with these neighbourhoods in a different way. This will be facilitated by the fact that $\im(\Phi)+(|\theta|^2+|\tau|^2)^{\frac 12}|\nabla_y\Phi|^2$ is bounded away from $0$ on this set.



In the $C^{\infty}$ theory, Fourier {integral operators} are completely described by the Lagrangians of their phase functions; the usual proof consists in a term-by-term identification or construction (see for instance~\cite[Proposition~25.1.5]{hormander_analysis_1985}). The analytic theory is a bit subtler, and we first describe how to remove ``spurious'' phase variables; we will prove in effect that Fourier {integral operators} associated with ``complicated'' phase functions can also be written using simpler phases, which will be enough for our needs.

\begin{prop}\label{prop:excess}
Let $\phi$ be a positive
nondegenerate phase. Then on every point of $\Sigma_{\phi}$,
\[
n_{\theta}-\rank\nabla^2_{\theta}\phi=n_x+n_y-\rank\left(\pi:\Lambda_{\phi}\to\C^{n_x+n_y}\right)
\]
\end{prop}
\begin{proof}
The proof is the same as in the real-valued case, see~\cite[Proposition~18.4.1]{treves_analytic_2022}.
\end{proof}

This proposition interacts very well with the condition of positivity in our definition of phase functions. Indeed, on $\Sigma_{\phi}^{\R}$, in the directions where $\nabla^2_{\theta}\phi$ is non-degenerate, it has positive imaginary part; we will therefore be able to apply the stationary phase without difficulty.


\section{Fourier integral operators}\label{sec:four-integr-oper}

\subsection{Basic properties}\label{sec:first-properties}

The main goal of this subsection is to give a suitable definition of Fourier integral operators in a real-analytic setting; they will be associated with positive phase functions in the sense of Definition~\ref{def:good_phase} and analytic amplitudes in the sense of Definition~\ref{def:analytic_amp}. Among other properties, we wish Fourier integral operators to preserve real-analytic functions and to behave well under composition.

The main obstacle to a direct and general definition is the fact that most of our objects are only defined locally, in small neighbourhoods of the subset where the imaginary part of the phase vanishes. The zones where the phase has positive imaginary part should not contribute to the integrals modulo real-analytic functions, but without the possibility of introducing cut-offs, it is difficult to prove a priori that it is the case.

We first describe a particular case of oscillatory integrals, in coordinates, then we proceed by successive generalisations and patching arguments.

\begin{defi}\label{def:osc_int}
Let $U\subset \R^{n_x}$ be an open set, and let $\Gamma\subset \R^{n_{\theta}}$ be a conical open set. Let $\phi$ be a one-homogeneous function on $U\times \Gamma$, real-analytic with respect to $x$, and such that $\im(\phi)\geq 0$ everywhere.

Then, given an analytic amplitude $a$ on $U\times \Gamma$, we define the following distribution on $U$:
\[
K_{\Gamma,\phi}(a):x\mapsto \int_{\Gamma} e^{i\phi(x,\theta)}a(x,\theta)\dd \theta.
\]
\end{defi}

\begin{prop}\label{prop:pos_phase_reg_int}
Let $U$ and $\Gamma$ be two open sets as above, let $\phi$ be a one-homogeneous function on a neighbourhood of $U\times \Gamma$, real-analytic with respect to $x$ and such that $\im(\phi)\geq 0$ everywhere. Let $a$ be an analytic amplitude on a neighbourhood of $U\times \Gamma$. Then $K_{\Gamma,\phi}(a)$ is real-analytic on
\[
\{x\in U, \:\forall \theta\in \Gamma,\;
\im(\phi)(x,\theta)>c|\theta|\}
\]
as well as on
\[
\left\{x_0\in U, \exists C,c>0,\:\forall \theta\in \Gamma,\: \forall x\in \Neigh(x_0,\widetilde{U}),\; |\widetilde{a}|(x,\theta)<Ce^{-c|\theta|}\right\}.
\]
\end{prop}
\begin{proof}
Under either of the conditions above, the integrand $(x,\theta)\mapsto e^{i\phi(x,\theta)}a(x,\theta)$ extends holomorphically into an integrable function of $\theta$ for $x$ close to the real locus.
\end{proof}

A particular case of oscillatory integral above consists in singular integral kernels of operators acting as follows:
\begin{equation}\label{eq:local_FIO}
x\mapsto \int_{\Gamma} e^{i\phi(x,\theta,y)}a(x,\theta,y)v(y)\dd \theta \dd y,
\end{equation}
that is, the singular integral kernel of a Fourier integral operator (with phase $\phi$ and amplitude $a$). If $\phi$ is a positive phase function and $a$ is also real-analytic with respect to $y$, then these integral operators preserve real-analytic functions, as long as $\im(\phi)>0$ on the boundary of $\Gamma$.

\begin{prop}\label{prop:reg_to_reg}
Let $U\subset \R^{n_x}$, $\Gamma\subset
\R^{n_{\theta}+n_y}$ be open sets, with $\Gamma$ conical in its first variable and relatively compact in its second variable. Let $\phi$ be a positive phase function on a neighbourhood of $U\times \Gamma$ such that
\[
\exists c>0, \:\forall (x,\theta,y)\in U\times \partial \Gamma,\quad \im(\phi)(x,\theta,y)\geq c|\theta|,
\]
and let $a$ be an analytic amplitude on a neighbourhood of $U\times \Gamma$. Then for every $v\in C^{\omega}(\R^{n_y})$, the function
\[
x\mapsto \int_{\Gamma} e^{i\phi(x,\theta,y)}a(x,\theta,y)v(y)\dd \theta \dd y
\]
is real-analytic on $\bar{U}$.
\end{prop}
\begin{proof}
Let $v\in C^{\omega}(\R^{n_y})$ and consider
\[
u:x\mapsto \int_{\Gamma}e^{i\phi(x,\theta,y)}a(x,\theta,y)v(y)\dd \theta \dd y.
\]
We will change the contour of integration in the variable $y$. Since $\phi$ is a positive phase function, one has
\[
\exists c>0,\:\forall (x,\theta,y)\in \bar{U}\times \Gamma,\quad
\im(\phi)+|\nabla_y\phi|^2\geq c|\theta|.
\]
We now deform contours by following the flow of $\nabla_y\im(\widetilde{\phi})$. After a small time, we obtain a contour $\Gamma_1$ on which $\im(\widetilde{\phi})>c|\theta|$. Now, by the Stokes formula,
\[
u(x)=\int_{\Gamma_1}e^{i\widetilde{\phi}(x,\theta,y)}\widetilde{a}(x,\theta,y)\widetilde{u}(y)\dd
\theta \dd y + \int_{\Gamma_2}e^{i\widetilde{\phi}(x,\theta,y)}\widetilde{a}(x,\theta,y)\widetilde{u}(y)\dd
\theta \dd y,
\]
where $\Gamma_2$ lies within a small neighbourhood of $\partial \Gamma$ in $\R^{n_{\theta}}\times \C^{n_y}$, so that $\im(\widetilde{\phi})>c|\theta|$ on $V_2$. Thus, by Proposition~\ref{prop:pos_phase_reg_int}, $u$ is real-analytic.
\end{proof}

Over the course of this article, the open sets on which phases and amplitudes are well-defined will seldom be of a product form $U\times
\Gamma$, unless restricted to small neighbourhoods of a point. We have to patch together expressions of the form~\eqref{eq:local_FIO}, which is possible thanks to Proposition~\ref{prop:patching_analytic_func_or_distros}. Before doing so we introduce a relevant function space.
\begin{defi}\label{def:spaces}
If $K$ is a compact set of a real-analytic paracompact manifold, we denote by $\clE(K)$ the space of smooth functions on $K$ in the sense of Whitney~\cite{whitney_analytic_1934}, and by $\clO(K)$ the space of functions which extend to real-analytic functions on a neighbourhood of $K$.

Let $U$ be a relatively compact open set of a real-analytic paracompact manifold. We denote
\[
\clF(U)=\clE'(\bar{U})/(\clE'(\partial U)\cup \clO(\bar{U})).
\]
\end{defi} We refer to Appendix~\ref{sec:funct-spac-analyt} for basic properties of these spaces.



\begin{defi}\label{def:FIO_fixed_fibre}
Let $X$ and $Y$ be two paracompact analytic manifolds. Let $U\subset X$ and $V\subset Y$ be open and relatively compact. Let $\Omega\subset X\times \R^{n_{\theta}}\times V$ be open and conical in its second variable.

Let $\phi$ be a positive phase function on a neighbourhood of $\Omega$, such that $\im(\phi)\geq 0$ everywhere and
\begin{multline}\label{eq:boundary_positivity}
\exists c>0, \forall x\in U, \forall (\theta,y)\in
\partial\bigl(\{(\theta',y')\in \R^{n_{\theta}}\times V,(x,\theta,y)\in \Omega\}\bigr),\\ \im(\phi)(x,\theta,y)\geq c|\theta|.
\end{multline}
Let $a$ be an analytic amplitude on a neighbourhood of $\Omega$, and let $v\in \clF(V)$. We define $I_{\phi,\Omega,U}(a)v$ as the element of $\clF(U)$ defined by the following recipe.
\begin{enumerate}
\item For every $x_0\in U$, for some small neighbourhood $\Gamma_0\times V_0$ of
\[
\bigl\{(\theta,y)\in \R^{n_{\theta}}\times V,\quad(x_0,\theta,y)\in
\Omega,\quad\im(\phi)(x,\theta,y)=0\bigr\},
\]
consider $\chi\in C^{\infty}(X,\R)$ be equal to $1$ on a small neighbourhood $U_0$ of $x_0$ and supported on a small neighbourhood of $U_0$. Then
\[
x\mapsto
\int_{\Gamma}e^{i\phi(x,\theta,y)}a(x,\theta,y)\chi(x)v(y)\dd \theta
\dd y
\]
is the action of a ``usual'' Fourier integral operator on $v$, and in particular it belongs to $\clE'(X)$ (in fact it is supported near $U_0$). This element of $\clE'(X)$ can then be restricted to an element of $\clE'(U_0)/\clE'(\partial U_0)$, see Proposition~\ref{prop:presheaf}. By Proposition~\ref{prop:reg_to_reg}, if $v\in C^{\omega}(\bar{V_0})$ then this produces an element of $C^{\omega}(\bar{U_0})$.

\item Perform the last item on a finite family of open sets covering $U$.
\item By Proposition~\ref{prop:pos_phase_reg_int}, any two such distributions agree on the intersection of their defining supports modulo a real-analytic function, and by Proposition~\ref{prop:reg_to_reg}. Therefore, by Proposition~\ref{prop:patching_analytic_func_or_distros}, we can patch these distributions together into a uniquely, well-defined element of $\clF(U)$, which is $0$ if $v\in C^{\omega}(\bar{V})$. In conclusion, we've built $I_{\phi,\Omega,U}(a):\clE'(\bar{V})/C^{\omega}(\bar{V})\to\clF(U)$.
\item One has then, by Proposition~\ref{prop:reg_to_reg},
\begin{multline*}
SS_a(I_{\Omega,U,\phi}(a)v)\\
\subset \bigl\{x\in U, \exists (\theta,y)\in
\R^{n_{\theta}}\times SS_a(v), \;(x,\theta,y)\in \Omega,
\im(\phi)(x,\theta,y)=0\bigr\},
\end{multline*}
where $SS_a$ denotes the analytic singular support (complement of the largest open set on which a function is real-analytic).

In particular, if $v\in \clE'(\partial V)$, then $SS_a(I_{\Omega,U,\phi}(a)v)=0$. Therefore, we can conclude:
\[
I_{\phi,\Omega,U}(a):\clF(V)\to \clF(U).
\]
\end{enumerate}
\end{defi}

\begin{rema}\label{rem:EtoO}
It is natural to ask whether this construction works when $v$ is a
\emph{singularity hyperfunction}, that is, an element of $\clO'(\bar{V})/(\clO'(\partial V)+\clO(V))$. The answer is yes, but to do so we need to study the analytic wave front set of the kernel $K_{\Gamma,\phi}$; we will do so in Section~\ref{sec:advanced-properties}, after having defined the analytic wave front set as the set of analytic singularities after conjugation by a Fourier integral operator; for the moment we only invoke the fact that the smooth wave front set of $K_{\Gamma,\phi}$ is known and therefore analytic Fourier {integral operators} act nicely on distributions. After making sure that Definition~\ref{def:FIO_fixed_fibre} makes sense as acting on singularity hyperfunctions, the rest of Section~\ref{sec:four-integr-oper} will extend to this case without any difficulty.
\end{rema}

As in the smooth case, these operators behave well under composition and stationary phase under natural geometric hypotheses, and particularly so when the analytic amplitudes in question are realisations of formal analytic amplitudes.

\begin{prop}\label{prop:compo_well_def}
Let $X,Y,Z$ be paracompact
real-analytic manifolds, and let $U\subset X,V\subset Y,W\subset Z$ be relatively compact open sets. Let $\Omega_1\subset X\times \R^{n_{\theta}}\times V$ and $\Omega_2\subset Y\times \R^{n_{\tau}}\times W$ be open and conic with respect to their middle variables; let $\phi_1$ and $\phi_2$ positive non-degenerate phase functions on open neighbourhoods of, respectively, $\Omega_1$ and $\Omega_2$. Suppose that $(\phi_1,\Omega_1,U)$ and $(\phi_2,\Omega_2,V)$ both satisfy~\eqref{eq:boundary_positivity}. Let $a_1:\Omega_1\to \C$ and $a_2:\Omega_2\to \C$ be analytic amplitudes.

Let
\begin{align*}
\Phi:(x,\theta,y,\tau,z)&\mapsto
\phi_1(x,\theta,y)+\phi_2(y,\tau,z)
\\
A:(x,\theta,y,\tau,z)&\mapsto a_1(x,\theta,y)a_2(y,\tau,z).
\end{align*}

There exists a contour deformation $\Gamma\in \widetilde{Y}$, with real-analytic dependence on $x,\theta,\tau,z$, and a neighbourhood $\Omega$ of
\[
\left\{(x,\theta,\tau,z)\in X\times \R^{n_{\theta}+n_{\tau}}\times Z,\; y\in
\Gamma(x,\theta,\tau,z),\nabla_{y,\theta,\tau}\Phi=0\right\},
\]
contained in $\{\epsilon|\theta|<|\tau|<\epsilon^{-1}|\theta|\}$ for some $\epsilon>0$, such that $(\phi_1+\phi_2,\Omega,U_1)$ satisfies~\eqref{eq:boundary_positivity}, and
\[
I_{\phi_1,\Omega_1,U_1}(a_1)\circ I_{\phi_2,\Omega_2,U_2}(a_2)=I_{\Phi,\Omega,U_1}(A)
\]
(as maps from $\clF(W)$ to $\clF(U)$). Moreover, if $a_1$ and $a_2$ respectively realise formal analytic amplitudes, then there exists a maximally analytic amplitude $A'$ on $\Omega$ such that
\[
I_{\Phi,\Omega,U_1}(A-A')=0.
\]
\end{prop}
\begin{proof}
The distributional kernel of $I_{\phi_1,\Omega_1,U_1}(a_1)\circ I_{\phi_2,\Omega_2,U_2}(a_2)$ is
\[
(x,z)\mapsto \int e^{i\Phi(x,\theta,y,\tau,z)}A(x,\theta,y,\tau,z)\dd y \dd \theta
\dd \tau.
\]
We begin with a contour deformation on $y$, which follows the positive gradient flow of $(|\theta|^2+|\tau|^2)^{-\frac 12}\nabla_y\im(\Phi)$. After a small time, along this contour $\Gamma_1$ one has
\[
\im(\Phi)\geq c|\theta||\nabla_{\theta}\Phi|^2+c|\tau||\nabla_{\tau}\Phi|^2+c\left(|\theta|^2+|\tau|^2\right)^{-\frac 12}|\nabla_y\Phi|^2.
\]
Compared to the integral on the original contour, the difference is, by Stokes' theorem, an integral of the form
\[
\int_{V} e^{i\widetilde{\Phi}(x,\theta,y,\tau,z)}\widetilde{A}(x,\theta,y,\tau,z)\dd y \dd \theta
\dd \tau
\]
where on $V= \bigcup_{t\in [0,1]}\partial
\Gamma_t(x,\theta,\tau,z)$, the phase $\widetilde{\Phi}$ has positive imaginary part, due to the fact that $(\phi_1,\Omega_1,U_1)$ satisfies~\eqref{eq:boundary_positivity}. Hence, by Definition~\ref{def:FIO_fixed_fibre}, the difference maps $\clE'(\bar{W})$ into $C^{\omega}(\bar{U})$ (note that, in this application of Definition~\ref{def:FIO_fixed_fibre}, $\Phi$ is not a positive phase function, because it is singular at $\theta=0$ and $\tau=0$; for the same reasons $A$ is not an analytic amplitude).

We now use Proposition~\ref{prop:pos_phase_reg_int} again to remove the complement set of a conical neighbourhood $\Omega$ of $\{\im(\Phi)=0\}$.

Let us prove that, on $\Omega$, $\theta$ and $\tau$ are bounded away from $0$. Indeed, if $\theta=0$, then $\Phi=\phi_2$, and $|\tau||\nabla_\tau\phi_2|^2$ is bounded away from $0$ except in the vicinity of $\Sigma_2$, where $|\tau|^{-1}|\nabla_y\phi|^2$ is bounded away from $0$. Therefore $\im(\Phi)$ is bounded away from $0$ on $\{\theta=0\}$, and on $\{\tau=0\}$ as well, by a symmetrical argument. This concludes\linebreak the proof.

To conclude, if $a_1$ and $a_2$ realise formal analytic amplitudes, then so does $A$, by Proposition~\ref{prop:product_preparation}; letting $A'$ be a maximally analytic realisation of the same analytic amplitude, one has $|A'-A|\leq Ce^{-c(|\theta|+|\tau|)}$ so that, by Proposition~\ref{prop:pos_phase_reg_int}, the difference between the associated Fourier {integral operators} maps $\clE'(\bar{W})$ into real-analytic functions.
\end{proof}

\begin{prop}\label{prop:stat_phase_in_FIO}
Let $X$ and $Y$ be paracompact real-analytic manifolds and let $U,V$ be relatively compact sets of $X$ and $Y$. Let $\Omega\subset X\times \R^{n_{\theta'}}\times
\R^{n_{\omega}}\times V$ be open and a relatively compact cone with respect to the second variable; let $\phi$ be a positive phase function on a neighbourhood of $\Omega$ such that $(\phi,\Omega,U)$ satisfies condition~\eqref{eq:boundary_positivity}. Suppose further that $\nabla_\omega^2\phi$ is nonsingular everywhere.

Let $a:\Omega\to \C$ be such that $((x,y,\omega),\xi)\mapsto a(x,\xi,\omega,y)$ is an analytic amplitude.

Then there exists $J\in \N$ and, for every $1\leq j\leq J$, a conic open set $\Omega_j\subset X\times
\R^{n_{\theta'}}\times V$, a non-degenerate phase function $\phi_j$ on a neighbourhood of $\Omega_j$ such that $(\phi_j,\Omega_j,U)$ satisfies condition~\eqref{eq:boundary_positivity}, and an analytic amplitude $b_j:\Omega_j\to \C$, such that
\[
I_{\phi,\Omega,U}(a)-\sum_{j=1}^JI_{\phi_j,\Omega_j,U}(b_j)
\]
maps $\clE'(\bar{V})$ into $C^{\omega}(\bar{U})$. Moreover, the union of the Lagrangians of $\phi_j$ is the Lagrangian of $\phi$, and if $a$ realises some formal analytic amplitude $(a_k)_{k\in \N}$, then each $b_j$ realises a formal analytic amplitude $(b_{j;k})_{k\in \N}$, which can be obtained from $(a_k)_{k\in \N}$ by the formal stationary phase.
\end{prop}
\begin{proof}
Choosing local expressions for $I_{\phi,\Omega,U}(a)$ as in Definition~\ref{def:FIO_fixed_fibre}, this is a restatement of Proposition~\ref{prop:stat_phase}.
\end{proof}

Many natural Fourier {integral operators} make sense of the fibre variable $\theta$ as belonging to the fibre of a bundle over the base~$X\times Y$ rather than a fixed space~$\R^{n_{\theta}}$. Similarly, in practical applications of Proposition~\ref{prop:stat_phase_in_FIO}, the decomposition of the phase variables into the directions where $\nabla^2\phi$ is non-degenerate and a complement set is seldom independent on $x,y$.

Both issues are formally taken care of by a one-homogeneous real-analytic change of variables $(x,\theta,y)\rightsquigarrow (x,\xi(x,\theta,y),y)$; however, analytic amplitudes do not have real-analytic dependence on $\theta$, and therefore, these change of variables destroy real-analyticity with respect to $x$ and $y$. We can remediate to this problem in the case of realisations of formal analytic amplitudes.

\begin{prop}\label{prop:change_phase_variables}
Let $X$ and $Y$ be
paracompact analytic manifolds, let $U\subset X$ and $V\subset Y$ be relatively compact. Let $\Omega\subset X\times \R^{n_{\theta}}_{\theta}\times V$ suppose that $\Omega$ is conical in its second variable.

Let $\xi_t:\Omega\to X\times \C^{n_{\theta}}\times V$ be a family of contours (they extend into complex-analytic diffeomorphisms, are one-homogeneous in the second variable, act as identity on the first and third variables, depending smoothly on the parameter $t$) with $\xi_0=\idd$. Let $\phi$ be a phase function on a neighbourhood of $\Omega$ and suppose that $(\phi_t\circ
\xi_t,\Omega,U)$ satisfies~\eqref{eq:boundary_positivity} for all $t\in [0,1]$. Let $a$ be the realisation of a formal analytic amplitude $(a_k)_{k\in \N}$ on a neighbourhood of $\Omega$.

Then, any realisation $b$ of $(a_k\circ \xi_1)_{k\in \N}$ is such that (as maps from $\clF(V)$ to $\clF(U)$)
\[
I_{\phi\circ \xi_1,\Omega,U}(b\det(D_{\theta}\xi_1))=I_{\phi,\Omega',U}(a).
\]
\end{prop}
\begin{proof}
Let $\xi_t$ be a homotopy of complex-analytic diffeomorphisms that are one-homogeneous and act as identity on the first and third variables, with $\xi_0=\idd$. Then $(a_k\circ \xi_t)_{k\in \N}$ is, uniformly for $t\in [0,1]$, a formal analytic amplitude; as in the proof of Proposition~\ref{prop:real_formal_symb}, one can realise it into
\[
b(t,x,y,\theta)=\sum_{k\in \N}a_k\circ
\xi_t(x,\theta,y)\left(1-\chi_k\left(\frac{c\theta}{k}\right)\right)
\]
for some $c>0$ small, where $\chi_k$ is a radial function, supported on $B(0,2)$ and equal to $1$ on $B(0,1)$, and such that $|\nabla^j\chi_k|\leq (Rk)^j$ whenever $j\leq k$.

Let now $x_0\in U$. Following Definition~\ref{def:FIO_fixed_fibre}, let $U_0$ be a small neighbourhood of $x_0$ in $U$, let $\Gamma_0,V_0$ an open neighbourhood of $\{\im(\phi)(x_0,y,\theta)=0\}$, and let $v\in \clE'(V_0)$, then for $t$ small, we are considering
\[
J(t):x\mapsto \int_{\Gamma_0\times V_0 } e^{i\phi(x,\xi_t(x,\theta,y),y)}b(t,x,y,\theta)\det(D_{\theta}\xi_t)v(y)\dd
\theta\dd y,
\]
For some $t_0>0$, one has still $\im(\phi)\circ \xi_t>0$ on $U_0\times \partial \Gamma_0\times V_0$ for all $t\in [0,t_0]$.

For every holomorphic function $f$,
\[
\frac{\dd}{\dd t}(f(x,\xi_t(x,\theta,y),y)\det(D_{\theta}\xi_t))=\divv_{\theta}\left(f(x,\xi_t(x,\theta,y))\partial_t\xi_t\right),
\]
and therefore, by Stokes' formula,
\begin{multline}\label{eq:Jprime}
\frac{\dd J(t)}{\dd t}(x)\\
=\int_{\Gamma_0\times V_0}e^{i\phi(x,\xi_t(x,\theta,y),y)}\det(D_{\theta}\xi_t)v(y)\sum_{k\in
\N}a_k\circ\xi_t(x,\theta,y) \partial_t\xi_t\cdot
\nabla_{\theta}\chi_k\left(\frac{c\theta}{k}\right)\dd
\theta \dd y\\
+\int_{\partial \Gamma_0\times V_0}e^{i\phi(x,\xi_t(x,\theta,y),y)}\det(D_{\theta}\xi_t)v(y)\sum_{k\in
\N}b(t,x,y,\theta)\dd
\theta \dd y.
\end{multline}
We can apply Proposition~\ref{prop:pos_phase_reg_int} to the second term of the right-hand side, since $b$ is an analytic amplitude and $\phi\circ\xi_t$ has positive imaginary part on $\partial \Gamma_0\times V_0$; therefore this term is an analytic function of $x$. As for the first term, let us prove that there exists $C>0, \epsilon>0$ such that, uniformly for $(x,y)$ in a complex neighbourhood,
\[
\left| \sum_{k\in
\N}a_k\circ\xi_t(x,\theta,y) \partial_t\xi_t\cdot
\nabla_{\theta}\chi_k\left(\frac{c\theta}{k}\right)\right|\leq Ce^{-\epsilon|\theta|}.
\]
Since $(a_k)_{k\in \N}$ is a formal analytic amplitude, one has
\[
|a_k\circ\xi_t(x,\theta,y)|\leq CR^kk!|\theta|^{-k}.
\]
Moreover $\nabla_{\theta}\chi_k(\frac{c\theta}{k})$ is only nonzero whenever $\frac{c|\theta|}{2}\leq k \leq c\theta$. In particular,
\[
\left|a_k\circ\xi_t(x,\theta,y)\nabla_{\theta}\left(\chi_k\left(\frac{c\theta}{k}\right)\right)\right|\leq C(cR)^{\frac{c|\theta|}{2}}.
\]
Moreover for fixed $\theta$ the number of nonzero terms is of order $|\theta|$, so that we finally obtain, whenever $c<\frac{1}{K}$, for some $\epsilon>0$, the required estimate.

In particular, we can apply Proposition~\ref{prop:pos_phase_reg_int} to the first term of the right-hand-side of~\eqref{eq:Jprime} as well; to conclude, for $t$ small, $\frac{\dd J(t)}{\dd t}$ continuously sends $\clE'(Y)$ into $\clO(U)$. Integrating this fact on $[0,t_0]$ for $t_0$ small, we obtain the desired identity with $\xi_1$ replaced with $\xi_{t_0}$. However the statement is clearly of a local nature, and we can cover $[0,1]$ with a finite number of open sets on which we can apply the argument above. This concludes the proof of Proposition~\ref{prop:change_phase_variables}.
\end{proof}

From now on we will only consider analytic Fourier integral operators whose amplitudes are realisations of formal analytic amplitudes; by Proposition~\ref{prop:pos_phase_reg_int}, they do not depend on the realisation, and we will simply pass formal amplitudes as arguments of $I_{\phi,\Omega,U}$.

Assembling the previous facts, we obtain the following general result.

\begin{theo}\label{prop:FIOs}
Let $X,Y$ be paracompact real-analytic manifolds. Let $E$ be a real-analytic cone bundle over $X\times Y$, and denote the base projection by $\pi=(\pi_X,\pi_Y)$. Let$V\subset Y$ be open and relatively compact. Let $\Omega\subset \pi^{-1}(X\times V)$ be open and conical, and let $\phi$ be a positive phase function on a neighbourhood of $\Omega$. Let $a$ be a formal analytic amplitude on a neighbourhood of $\Omega$.
\begin{enumerate}
\item Given $U\subset X$, open such that $(\phi,\Omega,U)$ satisfies~\eqref{eq:boundary_positivity}, any choice of local trivialisations of $E$ leads (via Definition~\ref{def:FIO_fixed_fibre}) to the same, well-defined operator $I_{\phi,\Omega,U}(a)$ mapping $\clF(V)$ to $\clF(U)$. When further quotiented by smooth functions, $I_{\phi,\Omega,U}(a)$ coincides with the usual formal construction of the Fourier {Integral} operator with phase $\phi$ and formal amplitude $a$. Furthermore, for all $v\in
\clE'(Y)$,
\begin{multline*}
SS_a(I_{\phi,\Omega,U}(a)v)\\
\subset \left\{x\in U, \exists \Theta \in \Omega,\;
\pi_X (\Theta)=x, \;\pi_Y(\Theta)\in SS_a(v),\; \im(\phi(\Theta))=0\right\}.
\end{multline*}
We call \emph{analytic Fourier integral operator} such an $I_{\phi,\Omega,U}(a)$.

\item If $U$ is relatively compact, $\Omega$ is a relatively compact open cone, and\linebreak ${\rank \nabla^2_{\theta}\phi\geq 1}$ on a neighbourhood of $\Omega$, then there exists a finite number of real-analytic vector bundles $(E_j)_{1\leq j\leq J}$, such that $\rank(E_j)=\rank(E)-1$, open conical sets $\Omega_j\subset E_j$, positive phase functions $\phi_j$ and formal analytic amplitudes $b_j$ on neighbourhoods of $\Omega_j$, such that
\[
I_{\phi,\Omega,U}(a)=\sum_{j=1}^JI_{\phi_j,\Omega_j,U}(b_j).
\]
As a formal amplitude, $b_j$ is obtained from $a$ by stationary phase.

\item Let $Z$ be a paracompact real-analytic manifold, $F$ a real-analytic cone bundle over $Y\times Z$, $W\subset Z$ open and conical, $\Omega'\subset \pi_F^{-1}(Y\times W)$ open and conical, $\phi'$ a positive phase function on a neighbourhood of $\Omega'$ such that $(\phi',F,V)$ satisfies~\eqref{eq:boundary_positivity}, $a'$ a formal analytic amplitude on a neighbourhood of $\Omega'$. Then the composition $I_{\phi,\Omega,U}(a)I_{\phi',\Omega',V}(a')$ is an analytic Fourier integral operator (constructed as in Proposition~\ref{prop:compo_well_def}).
\end{enumerate}
\end{theo}

To some extent, we have reproduced the $C^{\infty}$ theory in the analytic case. One thing still missing, and a crucial property in the $C^{\infty}$ case, is the fact that Fourier {integral} operators only depend on the positive Lagrangian associated with their phases: given two phases $\phi_1$ and $\phi_2$ with same Lagrangian, we expect any analytic Fourier integral operator associated with $\phi_1$ to be equal to some analytic Fourier integral operator associated with $\phi_2$. The proof in the $C^{\infty}$ case (see~\cite[Proposition~25.1.5]{hormander_analysis_1985}) relies on a construction of the amplitude order by order, which is difficult to translate to the analytic case (essentially, one would need to prove by hand that the constructed formal amplitude is analytic). We can, however, prove that some Lagrangians of interest possess a more or less canonical phase, so that all Fourier integral operators associated with said Lagrangians can be rewritten with this specific case. The first of these Lagrangians is the diagonal in $T^*\C^n$; we will prove in Proposition~\ref{prop:Kuranishi} the \emph{Kuranishi trick:} all such Fourier {Integral} operators are pseudodifferential operators in the sense that they can be rewritten using the usual phase and any good contour.

The second family of Lagrangians contains the Szeg\H{o} kernel parametrix on the boundary of any strictly pseudoconvex domain (the Lagrangian, of course, depends on the domain). These Lagrangians are idempotent, and Fourier {integral operators} with same Lagrangian are called ``covariant Toeplitz operators'' and form an algebra. To describe this algebra in practice, one writes all these Fourier integral operators with the same phase. We can do so, more generally, for Lagrangians whose base projection has corank 1.

\pagebreak{}
\begin{prop}\label{prop:FIO-rank1-phase}
Let $X,Y$ be paracompact analytic manifolds. Let $\Lambda\subset (T^*\widetilde{X}\setminus \{0\})\times (T^*\widetilde{Y}\setminus
\{0\})$ be a conical, real-analytic, positive Lagrangian and suppose that the projection onto the base $\pi:\Lambda\to \widetilde{X}\times \widetilde{Y}$ satisfies $\dim
\ker \dd\pi=1$ everywhere (so that $\Lambda$ is a half-line bundle over its projection $\pi(\Lambda)=Z$).

There exists a positive phase function on a half-line bundle over a neighbourhood of $Z$ such that $\Lambda_{\phi}=\Lambda$. For any such phase function, and any analytic Fourier integral operator with phase $\phi$, there exists a formal analytic amplitude $a$ on a neighbourhood of $Z$ such that $I_{\phi}(a)$ coincides with the analytic Fourier integral operator.
\end{prop}

\begin{proof}
By Theorem~\ref{prop:FIOs} and Propositions~\ref{prop:stat_phase_in_FIO} and~\ref{prop:excess}, possibly after application of stationary phase, an analytic Fourier integral operator with Lagrangian $\Lambda$ can be written with a number of phase variables equal to $1$, which means that the phase function is defined on a half-line bundle over a neighbourhood of $Z$. This proves the existence of such a phase.

Any such phase $\phi$ satisfies $Z=\{\phi=0\}$ and near $Z$ one has $\nabla_x\phi\neq 0,\nabla_y\phi\neq 0$. Since $Z$ has codimension~$1$, this means that $\phi$ is a defining function for $Z$. In particular, given another function $\phi_2$ satisfying the same properties, locally after identification of the half-line bundles over a neighbourhood of $Z$ on which the phases are defined, one has $\phi_2=f\phi$ for some non-vanishing function $f$. Thus, locally, Fourier {integral operators} with phase $\phi_2$ have integral kernels of the form
\[
(x,y)\mapsto \int_0^{+\infty}e^{iuf(x,y)\phi(x,y,1)}b(x,y,u)\dd u
\]
where we trivialised the half-line bundle and $b$ is a realisation of a formal analytic amplitude.

Now $u\mapsto f(x,y)u$ is a well-defined change of variables, so that by Proposition~\ref{prop:change_phase_variables}, there exists a formal analytic amplitude $a$ such that $I_{\phi}(a)$ is the same Fourier integral operator. This concludes the proof.
\end{proof}

\subsection{Pseudodifferential operators, FBI transforms, and Toeplitz operators}\label{sec:pseud-oper-1}

Pseudodifferential operators are a particular case of Fourier integral operators, with singular kernels of the form
\begin{equation}\label{eq:kernel_pseudos}
U\times V\ni (x,y)\mapsto \int_{\Gamma(x,y)} e^{i(x-y)\cdot (\xi+i\eta(x,\xi,y))}a(x,\xi,y)\dd \xi,
\end{equation}
where $U,V$ are open sets of $\R^n$, $\Gamma(x,y)$ is a conical contour, with real-analytic dependence in $(x,y)$, and which is \emph{positive}, meaning that
\begin{equation}\label{eq:good_ctr_pseudo}
\exists c, \:\forall (x,y)\in U\times V,\: \forall \xi\in \Gamma(x,y),\qquad\im((x-y)\cdot \xi)>c|\xi||x-y|^2;
\end{equation}
moreover $a$ realises a formal analytic amplitude on a neighbourhood of $\{(x,\xi,x),\linebreak x\in U\}$ in $U\times \R^n\times V$.

\begin{rema}\label{rem:pseudo_not_microlocal}
By the previous results, operators defined by~\eqref{eq:kernel_pseudos} do not depend on the realisation, nor do they depend on the choice of $\Gamma$. However, there is an important caveat at this stage: in order for condition~\eqref{eq:boundary_positivity} to hold, we need ${\Gamma(x,x)=\R^n}$. Indeed, otherwise the phase vanishes when $x=y\in U,\xi\in \partial \Gamma(x,x)$. This is, of course, a severe hindrance because we would like to study the action of real-analytic amplitudes defined near a point of phase space. This issue be addressed later on through the FBI transform. For the moment, we make note of this limitation and define pseudodifferential operators as follows.
\end{rema}

\begin{defi}\label{def:pseudos}
Let $U\Subset V\subset \R^n$, and let $a$ be a formal analytic amplitude on a neighbourhood of $\{(x,\xi,x),(x,\xi)\in U\times \R^n\}$ in $U\times \R^n\times V$. The pseudodifferential operator $\Opp(a)$ is defined as the operator with distributional kernel
\[
U\times V\ni (x,y)\mapsto
\int_{\Gamma(x,y)} e^{i(x-y)\cdot\xi}a^{\Gamma}(x,\xi,y)\dd \xi,
\]
where $\Gamma(x,y)$ is any contour satisfying~\eqref{eq:good_ctr_pseudo}, and where $a^{\Gamma}$ is such that the function $(x,\xi,y)\mapsto a^{\Gamma}(x,\theta(x,\xi,y),y)$ is a realisation of $(a_k(x,\theta(x,\xi,y),y))_{k\in \N}$, for a contour deformation $\theta$ which sends $\R^n$ to $\Gamma(x,y)$.

These operators preserve analytic functions, and they are uniquely defined as operators from $\clF(V)$ to $\clF(U)$.
\end{defi}

The Lagrangian associated with the phase of a pseudodifferential operator is always the diagonal of $(T^*U\setminus \{0\})^2$. Our first result is the \emph{Kuranishi trick}, which allows us to prove a reciprocal statement.

\begin{prop}\label{prop:Kuranishi}
Let $U\subset V\Subset \R^n$ and let $\pi:E\to \R^n\times \R^n$ be a cone bundle and let $\Omega\subset \pi^{-1}(E)$ be an open compact cone.. Let $\phi$ be a positive phase function on a neighbourhood of $\Omega$ such that $(\phi,\Omega,U)$ satisfies~\eqref{eq:boundary_positivity} and $\Lambda_{\phi}$ is the diagonal of $(T^*U\setminus \{0\})^2$. Let $a$ realise a formal analytic amplitude on a neighbourhood of $\Omega$.

There exists a formal analytic amplitude $b$ on a neighbourhood of $\{(x,\xi,x),(x,\xi)\in U\times \R^n\}$ in $U\times \R^n\times V$, such that
\[
I_{\phi,\Omega,U}(a)=\Opp(b)
\]
(as always, in the sense of operators from $\clF(V)$ to $\clF(U)$).
\end{prop}
\begin{proof}
We first apply Propositions~\ref{prop:excess} and~\ref{prop:stat_phase_in_FIO} to reduce ourselves to the case where $N=n$, since the right-hand side of the formula in Proposition~\ref{prop:excess} is equal to $n$. Note that when doing so, we obtain a finite sum of Fourier {integral operators}.

Since $\Lambda_{\phi}$ is the diagonal of $(T^*U\setminus \{0\})^2$, then $\phi(x,\theta,y)$ vanishes when $x=y$, and $\theta\mapsto
\partial_{\theta}\phi(x,\theta,x)\in \R^n$ is a surjective local diffeomorphism. Therefore it is a global diffeomorphism, and it extends to a neighbourhood of $x=y$ into a map $(x,\theta,y)\mapsto
\zeta(x,\theta,y)$ which is a global diffeomorphism between the real domain of its second variable and a totally real contour of $\C^n$, and such that
\[
\phi(x,\theta,y)=(x-y)\cdot \zeta(x,\theta,y).
\]
The contour is $\R^n$ when $x=y$, and it satisfies~\eqref{eq:good_ctr_pseudo}, because $\phi$ is positive. Removing a neighbourhood of $x=y$ (away from which the phase has positive imaginary part) and applying Proposition~\ref{prop:change_phase_variables}, we obtain a pseudodifferential operator.
\end{proof}

\pagebreak
Given a compact real-analytic manifold $M$, one can equivalently define analytic pseudodifferential operators on $M$ as:
\begin{itemize}
\item gluings of Definition~\ref{def:pseudos} in coordinate charts;
\item Fourier integral operators on $M$ (or on open subsets of $M$) with Lagrangian $\diag(T^*M)$.
\item Fourier integral operators defined with a suitable global phase function. An example is given by the function $\phi:T^*M\times M\to
\C$ defined as follows: $\iota:M\to \R^N$ is an embedding of $M$ in Euclidian space of sufficiently large dimension, $\pi_x:\R^n\to T_xM$ is the pull-back of the orthogonal projection from $\R^n$ to $\iota(T_xM)$, and
\[
\phi(x,\xi,y)=\xi(\pi_x(\iota(y)))+i|\xi|\dist(\iota(x),\iota(y))^2.
\]
\end{itemize}

Pseudodifferential operators are not uniquely determined by their amplitudes; there is ``one variable too many'' among $(x,\xi,y)$. The notion of \emph{symbol} is better suited to their analysis.

\begin{prop}\label{prop:ampl_to_symb_pseudos}
Let $U\subset \R^n$ open and relatively compact and let $a$ be a formal analytic symbol on a neighbourhood $\Omega$ of $\{(x,\xi,x)\in \R^{3n},x\in U\}$. Let $V$ be a small neighbourhood of $U$.

Given $j\in \N$ define
\[
b_j:(x,\xi)=\sum_{|\mu|\leq j}\frac{i^{|\mu|}}{\mu!}\left(\partial_y^{\mu}\cdot
\partial_{\xi}^{\mu}a_{j-|\mu|}\right)(x,\xi,x).
\]
Then $(b_j)$ is a formal analytic amplitude and $\Opp((a_j))=\Opp((b_j))$. $(b_j)$ is called the (left)
\emph{total symbol} of $(a_j)$.
\end{prop}

\begin{proof}
It follows from a direct computation that $b_j$ is a formal analytic
amplitude.

Let $V$ be a small neighbourhood of $U$ such that $(b_j)_{j\in \N}$ is well-defined on a complex neighbourhood of $U\times \R^n\times V$. Let $\theta:U\times \R^n\times V\to \C^n$ be such that $\kappa:(x,\xi,y)\mapsto (x,\theta(x,\xi,y),y)$ extends to a biholomorphism with $\kappa(x,\xi,x)=\xi$ and $\phi:(x,\theta,y)\mapsto \phi\circ \kappa^{-1}$ is a positive phase function.

Let $a$ and $b$ be respective realisations of $a_j\circ \kappa^{-1}$ and $b_j\circ \kappa^{-1}$ as built in Proposition~\ref{prop:real_formal_symb}. Namely, with $\chi_j\in C^{\infty}(\R^n,\R)$ radial and such that
\begin{gather*}
\1_{\R^n\setminus B(0,1)}\leq \chi_j\leq \1_{\R^n\setminus B\left(0,\frac 12\right)}\\
\exists \rho_0,\:\forall N\leq j,\qquad
\left\|\nabla^N\chi_j\right\|_{L^{\infty}}\leq (\rho_0j)^N,
\end{gather*}
we have defined, for $c>0$ small,
\begin{align*}
a(x,\theta,y)&=\sum_{j\in
\N}a_j(x,\xi(x,\theta,y),y)\left(1-\chi_{j+1}\left(\frac{c\theta}{j+1}\right)\right)\\
b(x,\theta,y)&=\sum_{j\in
\N}b_j(x,\xi(x,\theta,y),y)\left(1-\chi_{j+1}\left(\frac{c\theta}{j+1}\right)\right),
\end{align*}
and then
\begin{align*}
\Opp((a_j))(x,y)&=\int e^{i\phi(x,\theta,y)}a(x,\theta,y)\dd \theta\\
\Opp((b_j))(x,y)&=\int e^{i\phi(x,\theta,y)}b(x,\theta,y)\dd \theta,
\end{align*}
modulo a regularising operator. Define the contour ${\Gamma(x,y)\mk=\{\xi\mk\in\mk \C^n, \theta(x,\xi,y)\mk\in \mk\R^n\}}$, so that
\begin{align*}
\Opp((a_j))(x,y)&=\int_{\Gamma(x,y)} e^{i(x-y)\cdot \xi}a\circ \kappa(x,\xi,y)\dd \xi\\
\Opp((b_j))(x,y)&=\int_{\Gamma(x,y)} e^{i(x-y)\cdot \xi}b\circ \kappa (x,\xi,y)\dd \xi.
\end{align*}
Now we use the standard manipulation: for every fixed $N\in \N$ and smooth function $f$ on $V$
\[
\forall x\in U, \:\forall y \in B(x,\epsilon), f(y)=\sum_{|\mu|< N}\frac{\partial^\mu f(x)}{\mu!}(y-x)^{\mu}+\sum_{|\mu|=N}R_\mu(f)(x,y)(y-x)^{\mu},
\]
where
\[
|R_\mu(x,y)|\leq C\frac{\|\partial^{\mu}f\|_{L^{\infty}}}{N!}.
\]
Let $\psi_j=1-\chi_j$, and let $c_1\ll c$. To apply the above formula at an order $N$ which grows like $c'|\xi|$, we write
\begin{multline*}
a\circ \kappa(x,\xi,y)=a\circ \kappa(x,\xi,y)\chi_1(c'\theta(x,\xi,y))\\+\sum_{N\in
\N}\left(\sum_{|\mu|<N}\frac{\partial^{\mu}_y}{\mu!}a\circ
\kappa(x,\xi,y)|_{y=x}(y-x)^{\mu}+\sum_{|\mu|=N}R_\mu(a\circ
\kappa)(x,\xi,y)(y-x)^{\mu}\right)\\ 
\times \left(\psi_{N+1}\left(\frac{c'\theta(x,\xi,y)}{N+1}\right)-\psi_{N}\left(\frac{c'\theta(x,\xi,y)}{N}\right)\right).
\end{multline*}
This amounts to
\begin{multline*}
a\circ \kappa(x,\xi,y)
=\sum_{\mu\in
\N^d}\frac{\partial_y^{\mu}}{\mu!}(a\circ
\kappa(x,\xi,y))|_{y=x}(y-x)^{\mu}\chi_{|\mu|}\left(\frac{c'\theta(x,\xi,y)}{|\mu|}\right)\\
+\sum_{\mu\in
\N^d}R_{\mu}(a\circ \kappa)(x,\xi,y)(y-x)^{\mu}\left(\psi_{|\mu|}\left(\frac{c'\theta(x,\xi,y)}{|\mu|}\right)-\psi_{|\mu|-1}\left(\frac{c'\theta(x,\xi,y)}{|\mu|-1}\right)\right).
\end{multline*}
Before proceeding further, we integrate by parts in $\Opp(a)$ to trade the powers of $x-y$ for derivatives in $\xi$:
\begin{align*}
&\Opp(a)(x,y)=\int_{\Gamma(x,y)}\sum_{\mu\in
\N^d}\partial_{\xi}^{\mu}\left[\frac{i^{|\mu|}\partial_y^{\mu}}{\mu!}(a\circ
\kappa(x,\xi,y))|_{y=x}\chi_{|\mu|}\left(\frac{c'\theta(x,\xi,y)}{|\mu|}\right)\right]\dd
\xi\\
&+\Mk\int_{\Gamma(x,y)}\Mk\sum_{\mu\in
\N^d}|i|^{|\mu|}\partial_{\xi}^{\mu}\left[R_{\mu}(a\circ
\kappa)(x,\xi,y)\left(\psi_{|\mu|}\left(\frac{c'\theta(x,\xi,y)}{|\mu|}\right)\Mk-\psi_{|\mu|-1}\left(\frac{c'\theta(x,\xi,y)}{|\mu|-1}\right)\Mk\right)\Mk\right]\dd \xi.
\end{align*}

Each term in this expansion is real-analytic with respect to $(x,y)$ in the variables $(x,\theta,y)$. Indeed, given $1\leq k\leq d$,
\begin{align*}
\partial_{y_j}(a\circ \kappa)=&\partial_{y_j}\kappa\cdot (\nabla a)\circ \kappa=(\partial_{y_j}a)\circ \kappa+\partial_{y_j}\theta\cdot (\nabla_{\xi}a)\circ \kappa\\
\partial_{\xi_j}(a\circ \kappa)=&\partial_{\xi_j}\theta\cdot (\nabla_{\xi} a)\circ \kappa
\end{align*}
and both $(\nabla_{\xi,y}) \theta)\circ \kappa^{-1}$ and $a$ are real-analytic with respect to $(x,y)$. Therefore, by induction, for every $\mu\in \N^{d}$ $[\partial_y^{\mu}(a\circ \kappa)|_{y=x}]\circ \kappa^{-1}$ is real-analytic with respect to $(x,y)$. The term $R_{\mu}$ only depends on these kind of derivatives, as evidenced by the general remainder formula for Taylor series
\[
R_{\mu}(f)(x,y)=\frac{|\mu|}{\mu!}\int_0^1\partial^{\mu}f(tx+(1-t)y)(1-t)^{N-1}\dd t.
\]
Let us now prove that the remainders are small:
\begin{multline*}
\left|\sum_{\mu\in
\N^d}\partial_{\xi}^{\mu}\left[R_{\mu}(a\circ
\kappa)(x,\xi,y)\left(\psi_{|\mu|}\left(\frac{c'\theta(x,\xi,y)}{|\mu|}\right)-\psi_{|\mu|-1}\left(\frac{c'\theta(x,\xi,y)}{|\mu|-1}
\right)\right)\right]\right|\\
\leq Ce^{-c_1|\theta(x,\xi,y)|},
\end{multline*}
thereby obtaining (by Proposition~\ref{prop:reg_to_reg}) that this term in $a\circ \kappa$ contributes to $\Opp(a)$ by a regularising operator. Given $\mu\in \N^d$, one has
\[
\psi_{|\mu|}\left(\frac{c'\theta}{|\mu|}\right)-\psi_{|\mu|-1}\left(\frac{c'\theta}{|\mu|-1}\right)\neq 0 \implies \frac 12 (|\mu|-1)\leq c'|\theta|\leq |\mu|.
\]
Now, for some $\rho>0,R>0$,
\[
\forall j+k\leq \xi/R,\quad \left|\nabla^j_y\nabla^k_{\xi}(a\circ \kappa)(x,\xi,y)\right|\leq C\rho^{j+k}(j+k)!|\xi|^{d-k},
\]
so that, given $\mu\in \N^d$ and $\theta$ such that $ \frac 12 (|\mu|-1)\leq c'|\theta|\leq |\mu|$, given $\N^{d}\ni \nu\leq
\mu$,
\begin{multline*}
\frac{\mu!}{\nu!(\mu-\nu)!}\partial_{\xi}^\nu R_{\mu}(a\circ
\kappa)(x,\xi,y)\partial_{\xi}^{\mu-\nu}\left(\psi_{|\mu|}
\left(\frac{c'\theta(x,\xi,y)}{|\mu|}\right)-\psi_{|\mu|-1}
\left(\frac{c'\theta(x,\xi,y)}{|\mu|-1}\right)\right)\\
\leq C(a)|\mu|\frac{|\mu|!}{(|\mu|-|\nu|)!}\rho^{2|\mu|}(c')^{|\mu|-|\nu|}|\xi|^{d-|\nu|}.
\end{multline*}
With respect to $|\nu|$, the right-hand side is a log-convex function; if $c'$ is small enough, at both endpoints $|\nu|=0$ and $|\nu|=|\mu|$, it is smaller than $Ce^{c_1|\theta|}$. Therefore
\[
\left|\partial_{\xi}^{\mu}\left[R_{\mu}(a\circ
\kappa)(x,\xi,y)\left(\psi_{|\mu|}\left(\frac{c'\theta(x,\xi,y)}{|\mu|}\right)-\psi_{|\mu|-1}\left(\frac{c'\theta(x,\xi,y)}{|\mu|-1}\right)\right)\right]\right|\leq Ce^{-c_1|\theta|},
\]
and for $\theta$ fixed the number of non-zero terms in the sum over $\mu$ is polynomial in $|\theta|$.

Similarly, within
\[
\sum_{\mu\in
\N^d}\partial_{\xi}^{\mu}\left[\frac{i^{|\mu|}\partial_y^{\mu}}{\mu!}(a\circ
\kappa(x,\xi,y))|_{y=x}\chi_{|\mu|}\left(\frac{c'\theta(x,\xi,y)}{|\mu|}\right)\right],
\]
the terms where at least one derivative with respect to $\xi$ hits $\chi_{|\mu|}$ are exponentially small: indeed, given $\mu\in \N^d$, the support of $\nabla
\chi_{|\mu|}\left(\frac{c'\theta(x,\xi,y)}{|\mu|}\right)$ lies within $\{\frac 12 |\mu|\leq c'|\theta|\leq |\mu|\}$, and the estimates on the derivatives are the same as previously.

To conclude the proof, we inject the expression of $a\circ \kappa$ in
\[
\sum_{\mu\in
\N^d}\partial_{\xi}^{\mu}\frac{i^{|\mu|}\partial_y^{\mu}}{\mu!}(a\circ
\kappa(x,\xi,y))|_{y=x}\chi_{|\mu|}\left(\frac{c'\theta(x,\xi,y)}{|\mu|}\right)
\]
to prove that this is close to $b\circ \kappa$. We obtain
\[
\sum_{\mu\in \N^d}\sum_{j\in
\N}\frac{i^{|\mu|}\partial_{\xi}^{\mu}\partial_y^{\mu}}{\mu!}\left.\left[a_j(x,\xi,y)\chi_j\left(\frac{c\theta(x,\xi,y)}{j}\right)\right]\right|_{y=x}\chi_{|\mu|}\left(\frac{c'\theta(x,\xi,y)}{|\mu|}\right),
\]
and again one can prove that this is exponentially small unless no derivative hits~$\chi_j$: the support of $\nabla\chi_j\left(\frac{c\theta(x,\xi,y)}{j}\right)$ lies in $\{\frac 12 j\leq c|\theta|\leq j\}$, and one must also have $|\mu|<2c'|\theta|$; as long as $c'<\frac{c}{4}$, this ensures $2|\mu|\leq j$, hence for all $\nu_y\leq \mu$ and $\nu_{\xi}\leq \mu$,
\begin{multline*}
\left|\frac{\mu!}{\nu_y!(\mu-\nu_y)!\nu_{\xi}!(\mu-\nu_{\xi})!}\partial^{\nu_y}_y\partial^{\nu_{\xi}}_\xi a_j(x,\xi,y)\partial^{\mu-\nu_y}_y\partial^{\mu-\nu_{\xi}}_\xi\chi_j\left(\frac{c\theta(x,\xi,y)}{j}\right)\right|\\
\leq \rho^{j+2|\mu|}\frac{j!|\mu|!}{(|\mu|-|\nu_y|)!(|\mu|-|\nu_\xi|)!}c^{2|\mu|-|\nu_y|-|\nu_{\xi}|}|\xi|^{d-j-|\nu_{\xi}|};
\end{multline*}
for $c$ small this log-convex function of $|\nu_y|$ and $|\nu_{\xi}|$ is smaller than $e^{-c|\theta|}$ at all four endpoints $(0,0)$, $(0,|\mu|)$, $(|\mu|,0)$, $(|\mu|,|\mu|)$.

It remains
\[
\sum_{\mu\in \N^d}\sum_{j\in
\N}\frac{i^{|\mu|}\partial_{\xi}^{\mu}\partial_y^{\mu}a_j}{\mu!}(x,\xi,x)\chi_j\left(\frac{c\theta(x,\xi,x)}{j}\right)\chi_{|\mu|}\left(\frac{c'\theta(x,\xi,y)}{|\mu|}\right),
\]
where we recall that
\[
b(x,\xi,y)=\sum_{\mu\in \N^d}\sum_{j\in
\N}\frac{i^{|\mu|}\partial_{\xi}^{\mu}\partial_y^{\mu}a_j}{\mu!}(x,\xi,x)\chi_{j+|\mu|}\left(\frac{c\theta(x,\xi,x)}{j+|\mu|}\right);
\]
therefore we want to prove that the following is exponentially small
\[
\sum_{\mu\in \N^d}\sum_{j\in
\N}\frac{i^{|\mu|}\partial_{\xi}^{\mu}\partial_y^{\mu}a_j}{\mu!}(x,\xi,x)\left[\chi_j\left(\frac{c\theta(x,\xi,x)}{j}\right)
\chi_{|\mu|}\left(\frac{c'\theta(x,\xi,y)}{|\mu|}\right)-\chi_{j+|\mu|}\left(\frac{c\theta(x,\xi,x)}{j+|\mu|}\right)\right].
\]
We first observe that this is zero whenever $j\geq 2c|\theta|$, and also whenever $j+|\mu|\leq c|\theta|$ and $|\mu|\leq 2c'|\theta|$. Therefore we are only interested in situations where $j\leq 2c|\theta|$ and either $|\mu|\leq 2c'|\theta|$, $j\geq (c-2c')|\theta|$, or $|\mu|\geq 2c'|\theta|$; since $c'\leq \frac c4$, this means that we are only considering cases where either $j$ or $|\mu|$ is larger than $2c'|\theta|$. In this setting,
\[
\left|\frac{\partial_{\xi}^{\mu}\partial_y^{\mu}a_j}{\mu!}(x,\xi,x)\right|\leq C(a)\mu!j!|\xi|^{d-j-|\mu|}\leq Ce^{-c_1|\theta|}.
\]
This concludes the proof of Proposition~\ref{prop:ampl_to_symb_pseudos}.
\end{proof}


\begin{prop}\label{prop:inversion_formal_pseudos}
Define the Moyal product as follows: given two formal analytic symbols $a(x,\xi),b(x,\xi)$, let
\[
(a\sharp b)_k=\sum_{n=0}^k\sum_{l=0}^{k-n}\sum_{|\beta|=n}\frac{1}{\beta!}\partial^\alpha_xa_l\partial_{\xi}^{\beta}b_{k-n-l}.
\]
\begin{enumerate}
\item\label{prop4.16.1} The symbol product of two formal analytic symbols is a formal analytic symbol; in fact the topology of formal analytic symbols is induced by a countable family of Banach spaces of analytic symbols that are Banach algebras for the Moyal product.

\item\label{prop4.16.2} If both $a$ and $b$ are defined on $U\times \R^n\times V$ where $U\Subset V$, then
\[
\Opp(a)\Opp(b)-\Opp(a\sharp b)
\]
is regularising.

\item\label{prop4.16.3} Moreover, for every formal analytic amplitude $a$ such that $a_0$ is bounded away from $0$, there exists a formal analytic amplitude $b$ such that $(a\sharp b)=1$, and there exists a formal analytic amplitude $c$ such that $c\sharp c=a$.

\item\label{prop4.16.4} Define the adjoint $a^*$ of a formal analytic amplitude $a$ as
\[
(a^*)_k:(x,\xi)=\sum_{|\mu|\leq k}\frac{i^{|\mu|}}{\mu!}\partial_x^{\mu}\partial_{\xi}^{\mu}\bar{a}_{k-|\mu|}(x,\xi).
\]
Then, given a formal analytic amplitude $a$ with $a^*=a$ and $a_0>0$, there exists a formal analytic amplitude $b$ such that $b^*\sharp b=a$.
\end{enumerate}
\end{prop}
\begin{proof}
Claims~\eqref{prop4.16.1} and~\eqref{prop4.16.3} only depend on the formal calculus and are well-known; see~\cite{sjostrand_singularites_1982}. A related fact, proved in Appendix~\ref{sec:algebra-symb-pseud}, is the following: with the symbol norms introduced in Definition~\ref{def:formal-amp}, for every $m$ large enough (depending on the dimension) and every $R\geq 2^{d+1}\rho^2$,
\[
\|a\sharp b\|_{S_m^{\rho,R}}\leq 12\|a\|_{S_m^{\rho,R}}\|b\|_{S_m^{\frac{\rho}{2},\frac{R}{4}}}.
\]

Claim~\eqref{prop4.16.2} is a particular case of Proposition~\ref{prop:stat_phase_in_FIO}.

To prove claim~\eqref{prop4.16.4}, fix a Banach space $S^{\rho,R}_m$ of symbols as in Definition~\ref{def:formal-amp}, containing $a,\sqrt{a_0},(\sqrt{a_0})^{-1}$ (where we invert with respect to the Moyal product) and satisfying the requirements of Proposition~\ref{prop:composition-formal}. Letting $b_0=\sqrt{a_0}$, by claim~\eqref{prop4.16.3},
\[
(b_0^*)^{-1}\sharp a\sharp (b_0)^{-1}=1+r
\]
where the inverse is taken with respect to the Moyal product and where $r$ is a formal analytic symbol of degree $-1$.

Up to increasing $R$, we can assume $\|r\|_{S^{\frac{\rho}{2},\frac{R}{4}_m}}<\frac{1}{12}$. Then the power series
\[
\sqrt{1+r}:=1+\frac{r}{2}-\frac{r\sharp r}{8}+\frac{r\sharp r
\sharp r}{16}+\dots
\]
converges to an element $c$ of $S^{\rho,R}_m$ such that $c^*=c$ and $c\sharp c = 1+r$. Setting $b=c\sharp b_0$ then concludes the proof of Proposition~\ref{prop:inversion_formal_pseudos}.
\end{proof}

Our next move is to translate the pseudodifferential algebra into one that is more suited to our microlocal needs, as a workaround to Remark~\ref{rem:pseudo_not_microlocal}.

\begin{defi}\label{def:FBI}
The FBI transform from $L^2(\R^n)$ to $L^2(\R^n\times S^{n-1})$ is the operator with singular kernel
\[
T:\R^n\times S^{n-1}\times \R^n\ni (x,\omega,y)\mapsto \int_0^{+\infty} e^{it\left((x-y)\cdot
\omega+i\frac{|x-y|^2}{2}\right)}t^{\frac{n+1}{4}}\dd t.
\]
\end{defi} The power of $t$ is chosen such that $T$ is continuous from $L^2(\R^n)$ to $L^2(\R^n\times S^{n-1})$ and $T^*T$ is a degree~$0$ pseudodifferential operator.

\begin{prop}\label{prop:FBI_holom_ext}
The FBI transform of a compactly supported distribution extends into an holomorphic function of $x-i\omega$ on $\R^n\times B(0,1)$.
\end{prop}

\begin{proof}
It suffices to remark that, writing the phase as
\[
t\left((x-y)\cdot \omega+i\frac{|x-y|^2}{2}+i\frac{(1-|\omega|^2)}{2}\right),
\]
we obtain a holomorphic function of $x-i\omega$ on $\R^n\times B(0,1)$, whose imaginary part is non-negative.
\end{proof}

Before proceeding further, we introduce a helpful notation: given $z\in \C^n$ we set $z^2=\sum_{j=1}^nz_j^2$; this is the holomorphic extension of $\R^n\ni x\mapsto |x|^2$. We will reserve $|z|^2$ for the Hermitian norm of $z$.

\begin{prop}\leavevmode
\label{prop:T*T}
\begin{enumerate}
\item\label{prop4.19.1} The Lagrangian associated with $T$ is
\[
\Lambda_T=
\left\{
\begin{aligned}
(x,t(\omega+i(x-y)),\omega,t(x-y),y,t(\omega+i(x-y))),\\
t>0,(x-y)\cdot \omega=-i\frac{(x-y)^2}{2}
\end{aligned}
\right\}.
\]
\item\label{prop4.19.2} $\Lambda_{T^*}\circ
\Lambda_T=\diag(T^*\C^n\setminus \{0\})$, so that $T^*T$ is an elliptic pseudodifferential operator.

\item\label{prop4.19.3} The projection of $\Lambda_{T}\circ\Lambda_{T^*}$ onto the base has corank $1$ everywhere. Moreover the intersection of $\Lambda_{T}\circ\Lambda_{T^*}$ and the real locus is a half-line bundle over $\diag(\R^{n}\times S^{n-1})$, so that the integral kernel of any operator of the form $T^*\Opp(a)T$, for a formal analytic amplitude $(a_j)_{j\in \N}$, or more generally of any Fourier integral operator with the same Lagrangian as $T^*T$, is real-analytic away from the diagonal, and near the diagonal takes the form,
\[
(x_1,\omega_1,x_2,\omega_2)\mapsto
\int_0^{+\infty}e^{it\psi(x_1,\omega_1,x_2,\omega_2)}b(x_1,\omega_1,x_2,\omega_2,t)\dd t+r(x_1,\omega_1,x_2,\omega_2).
\]
In this formula,
\begin{itemize}
\item
\begin{equation}\label{eq:phase_Toep}
\psi(x_1,\omega_1,x_2,\omega_2):=i\left(1+\left(\frac{x_1-i\omega_1-x_2-i\omega_2}{2}\right)^2\right),
\end{equation}
and in particular $\psi$ extends into a holomorphic function of $x_1-i\omega_1$ and $x_2+i\omega_2$ on $\R^n\times B(0,1)\times
\R^n\times B(0,1)$, and $\im(\psi)$ is positive away from $\diag(\R^n\times S^{n-1})$.
\item $b$ is the realisation of a formal analytic amplitude.
\item $r$ and $b$ extend into holomorphic functions of $x_1-i\omega_1$ and $x_2+i\omega_2$ on a neighbourhood of $\diag(\R^n\times S^{n-1})$ in $\R^{4n}$.
\item $b$ realises a formal analytic amplitude obtained from $a$ by usual stationary phase. The first term $b_0$ is the only function such that $b_0(x,\omega,x,\omega)=a_0(x,\omega)$ and which extends into holomorphic function of $\omega_1+ix_1$ and $\omega_2-ix_2$. More generally, for any $j,\ell\in \N$, $b_j(x,\omega,x,\omega)$ only depends on $a_l$ and its derivatives up to degree~$2j-\ell$ at $x,\omega$.
\end{itemize}
\end{enumerate}
\end{prop}
\begin{proof}
The description of $\Lambda_T$ follows by an immediate computation. Now
\begin{align*}
\Lambda_{T^*}\circ \Lambda_T=\left\{
\begin{aligned}
(y_1,\eta_1,y_2,\eta_2),\exists (x,t_1,t_2,\omega),\quad &\eta_2=t_2(\omega+i(x-y_2))\\
&\eta_1=t_1(\omega-i(x-y_1))\\
&t_2(\omega+i(x-y_2))=t_1(\omega-i(x-y_1))\\
&t_2(x-y_2)=t_1(x-y_1)\\
&(x-y_2)\cdot \omega=-i\frac{(x-y_2)^2}{2}\\
&(y_1-x)\cdot \omega=-i\frac{(x-y_1)^2}{2}
\end{aligned}
\right\}.
\end{align*}
The three first conditions ensure $\eta_1=\eta_2$. In particular, $\eta_1^2=\eta_2^2$, so that (using the last two conditions) $t_1=t_2$; finally the fourth condition gives $y_1=y_2$. Reciprocally, given $(y,\eta)\in T^*\R^d\setminus \{0\}$, setting $x=y,t_1=t_2=|\eta|,\omega=\eta/|\eta|$ yields $(y,\eta,y,\eta)\in \Lambda_{T^*}\circ \Lambda_T$.

We now consider
\begin{align*}
\Lambda_{T}\circ \Lambda_{T^*}=
\left\{
\begin{aligned}
(x_1,\xi_1,\omega_1,v_1,x_2,\xi_2,\omega_2,v_2),&\exists (y,t_1,t_2),\,\\
&\xi_1=t_1(\omega_1+i(x_1-y))\\
&v_1=t_1(x_1-y)\\
&\xi_2=t_2(\omega_2+i(y-x_2))\\
&v_2=t_2(x_2-y)\\
&t_1(\omega_1+i(x_1-y))=t_2(\omega_2+i(y-x_2))\\
&(x_1-y)\cdot \omega_1=-i\frac{(x_1-y)^2}{2}\\
&(y-x_2)\cdot \omega_2=-i\frac{(y-x_2)^2}{2}
\end{aligned}
\right\}.
\end{align*}
Given $(x_1,\omega_1,x_2,\omega_2)$ fixed, conditions 1, 3, and 5 give $\xi_1=\xi_2$; together with the two last conditions, $\xi_1^2=\xi_2^2$ yields $t_1=t_2$. Using again condition 5, we can write $y=\frac{x_1-i\omega_1+x_2+i\omega_2}{2}$, and then either the last two conditions amounts to
\[
\left(\frac{x_1-i\omega_1-x_2-i\omega_2}{2}\right)^2=-1;
\]
on this codimension~1 manifold, whose intersection with the real locus is exactly the diagonal, the space of solutions is a half-line described by $t_1$.



The rest of the claim follows from there by applying Propositions~\ref{prop:compo_well_def}, \ref{prop:stat_phase_in_FIO}, and~\ref{prop:FIO-rank1-phase}. The property of holomorphic extension of $b$ comes from that of the kernel of $T$ itself; the fact that holomorphic extension of real-analytic functions prescribed on the diagonal is unique comes from the fact that $\R^n\times B(0,1)$ is strongly pseudoconvex.
\end{proof}

The microlocal structure of operators of the form $T\Opp(a)T^*$ is exactly that of the candidate for the Szeg\H{o} projector at the boundary of $\R^n\times B(0,1)$, which is strongly pseudoconvex since its natural defining function $|\omega|^2-1$ is strongly p.s.h. Indeed,
\begin{align*}
\psi(x_1,\omega_1,x_2,\omega_2)&=\frac{i}{4}\left((x_1-x_2)^2+(\omega_1-\omega_2)^2+2i(x_1-x_2)\cdot(\omega_1+\omega_2)\right)
\\ &=-i\left(\frac 14\left[(x_1+i\omega_1)^2+(x_2-i\omega_2)^2-2(x_1+i\omega_1)\cdot (x_2-i\omega_2)\right]-1\right)
\end{align*}
is the holomorphic extension of the function equal to $i(1-|\omega|^2)$ on the diagonal $(x_1,\omega_1)=(x_2,\omega_2)$.

\begin{prop}\label{prop:Szego_is_FIO}
The Szeg\H{o} kernel on $\R^n\times S^{n-1}$ is a Fourier integral operator with phase $\psi$.
\end{prop}

\begin{proof}
Recall from Proposition~\ref{prop:m_n_amplitude} that $a:\xi\mapsto
e^{-|\xi|}\int_{S^{n-1}}e^{\xi\cdot y}\dd y$ is, up to an exponential factor, the realisation of an elliptic formal analytic amplitude. Thus
\[
b:\R^n\ni \xi \mapsto \frac{e^{-|\xi|}}{\int_{S^{n-1}}e^{\xi\cdot y}\dd y}
\]
is a radial function and the realisation of an elliptic formal analytic amplitude. Passing to spherical coordinates, the Szeg\H{o} kernel computed in Proposition~\ref{prop:flat_Grauert} now reads
\[
K(z,w)=\frac{1}{(2\pi)^n}\int_0^{+\infty}b(2r)\left[e^{-2r}\int_{S^{n-1}}e^{ir(z-\bar{w})\cdot
\nu}\dd \nu\right]r^{n-1}\dd r,
\]
which we simplify into
\[
K(z,w)=\frac{1}{(2\pi)^n}\int_0^{+\infty}e^{r\left(-2+\sqrt{-\sum_{j=1}^n\left(z_j-\bar{w_j}\right)^2}\right)}b(2r)a(ir(z-\bar{w}))r^{n-1}\dd r.
\]
If $(z,w)$ belongs to a neighbourhood to $\{z=w\in \R^n\times S^{n-1}\}$, then $i(z-\bar{w})$ lies in a neighbourhood of $2S^{n-1}$ in $\C^n$, so that $ir(z-\bar{w})\in \Omega$.

To conclude, near the diagonal, the Szeg\H{o} kernel takes the form of a Fourier integral operator whose phase is the holomorphic extension of
\[
\psi_2(z,\bar{z})=-2+2\lvert\im(z)|
\]
which is a defining function for $\R^n\times S^{n-1}$. By Proposition~\ref{prop:FIO-rank1-phase}, one can rewrite it as a Fourier integral operator with phase $\psi$. Away from the diagonal we already know by Proposition~\ref{prop:flat_Grauert} that the Szeg\H{o} kernel is real-analytic.
\end{proof}

\begin{prop}\label{prop:invesion_Toeplitz}
Given $U\subset \R^n\times
S^{n-1}$, open and relatively compact, we~let
\[
\clH(U)=\ker_{\clE'(\bar{U})}(\partial_b)/(C^{\omega}(\bar{U})+\clE'(\partial U)),
\]
and $S(U)$ the space of formal analytic amplitudes on $\bar{U}\times \R_+$.
\begin{enumerate}
\item\label{prop4.21.1} Given $V\subset \R^n\times S^{n-1}$ and $a\in S(V)$, there exists a neighbourhood $W$ of $\bar{V\times V}$ and a unique formal analytic amplitude $\widetilde{a}$ on $W\times \R_+$, holomorphic in the second variable and anti-holomorphic in the first variable.
\item\label{prop4.21.2} Given $U\Subset V$ open and $\epsilon>0$ small enough, letting $\Omega=\{(x,t,y)\in U\times \R^+\times V,\linebreak
\dist(x,y)<\epsilon\}$, then $(\phi,\Omega,U)$ satisfies~\eqref{eq:boundary_positivity} and $I_{\psi,\Omega,U}(\widetilde{a})$ maps $\clF(V)$ into $\clH(U)$.
\item\label{prop4.21.3} There exists a sequence $(B_j)_{j\in \N}$ of bi-differential operators on $\R^n\times S^{n-1}$ such that for every $j$, $B_j$ has total degree~$2j$, and for every relatively compact $V\subset \R^n\times S^{n-1}$, $S(V)$ is a unit algebra for the product
\[
(S(V))^2\ni (a,b)\mapsto (a\sharp b)_k=\left(\sum_{j+\ell=k}B_{k-j-\ell}(a_j,b_\ell)\right)_k.
\]
\item\label{prop4.21.4} Invertible elements for the product above are exactly the formal amplitudes $a=(a_j)_{j\in \N}$ such that $a_0$ never vanishes.
\item\label{prop4.21.5} For every $U_1\Subset U_2\Subset V$ and $a,b\in S(V)$, letting $\epsilon$ small enough and
\begin{align*}
\Omega_1&=\left\{(x,t,y)\in U_1\times \R^+\times U_2,\;
\dist(x,y)<\epsilon\right\}\\
\Omega_2&=\left\{(x,t,y)\in U_2\times \R^+\times V,\;
\dist(x,y)<\epsilon\right\},
\end{align*}
one has
\[
I_{\psi,\Omega_1,U_1}(\widetilde{a})I_{\psi,\Omega_2,U_2}(\widetilde{b})=I_{\psi,\Omega_1,U_1}(\widetilde{a\sharp b}).
\]
\item\label{prop4.21.6} The unit for $\sharp$ acts as the identity on $\clH$.
\end{enumerate}
\end{prop}
\begin{proof}
Since $\R^n\times S^{n-1}$ is strongly pseudoconvex, the diagonal of $(\R^n\times S^{n-1})^2$ is a totally real submanifold, where we reverse the CR structure on the left factor. Therefore any real-analytic function on an open set $V\subset \R^n\times S^{n-1}$ can be extended CR-holomorphically to a neighbourhood of $\diag(V)$ in $(\R^n\times S^{n-1})^2$.

Given a formal analytic amplitude $\widetilde{a}$ near the diagonal such that $\partial_{b;1}\widetilde{a}=0$ (by which we mean: it is anti-CR-holomorphic in the first variable), since $\partial_{b;1}\psi=0$ as well, local realisations of $I_{\psi,\Omega,U}(\widetilde{a})$ as in Definition~\ref{def:FIO_fixed_fibre} will indeed map $\clE'(\bar{V})$ into functions $u$ such that $\partial_{b;1}u$ is real-analytic on a neighbourhood of $\bar{U}$. Therefore, by Proposition~\ref{prop:loc-anal-hypoell-Grauert}, if $S$ denotes the (antiholomorphic) Szeg\H{o} projector one has indeed $u=(1-S)u+Su$ where $\partial_{b}Su=0$ and $(1-S)u$ is real-analytic on $\bar{U}$. In particular, $I_{\psi,\Omega,U}(\widetilde{a})$ indeed maps $\clF(V)$ into $\clH(U)$.

Recalling from Proposition~\ref{prop:T*T} that $\Lambda_{\psi}=\Lambda_{T}\circ \Lambda_{T^*}$ and that $\Lambda_{T^*}\circ \Lambda_T$ is the diagonal of $(T^*\C^n)^*$, we obtain that $\Lambda_{\psi}\circ \Lambda_{\psi}=\Lambda_{\psi}$. We also proved in Proposition~\ref{prop:T*T} that any analytic Fourier integral operator with Lagrangian $\Lambda_{\psi}$ can be written as a Fourier integral operator with phase $\psi$. Therefore, given $V,a,b, U_1,U_2,\Omega_1,\Omega_2$ as in the claim, there exists a formal analytic amplitude $c$ such that $I_{\psi,\Omega_1,U_1}(\widetilde{a})I_{\psi,\Omega_2,U_2}(\widetilde{b})=I_{\psi,\Omega_1,U_1}(c)$. The amplitude $c$ is obtained from $a$, $b$ and $\psi$ by stationary phase and change of variables in the middle variables, so that $c$ is CR-holomorphic on $\Omega_1$. By the usual $C^{\infty}$ calculus of these operators (see for instance~\cite{boutet_de_monvel_sur_1975}), as a formal amplitude, the restriction of $c$ to the diagonal is indeed of the form $a\sharp b$ as in the claim.

We know by Proposition~\ref{prop:Szego_is_FIO} that the Szeg\H{o} projector has a kernel real-analytic from the diagonal and of the form $I_{\psi,\Omega,U}(a)$, where $U=\R^n\times S^{n-1}$, $\Omega=(\R^n\times S^{n-1})^2$ and $a$ depends only on the fibre variable. By definition, $S$ acts as identity on $\clH$; its amplitude is a unit for the algebra above.

It remains to prove that elliptic amplitudes can be inverted, which is a claim of formal nature (by the product formula, it suffices to check it at the level of expansion of amplitudes, and not a priori as operators making sense from some function space to another). To this end, we first suppose that amplitudes are globally defined, so that for every $a\in S(\R^n\times S^{n-1})$, by Proposition~\ref{prop:T*T}, the operator $T^*I_{\psi,\Omega,V_1}(a)T$ is a pseudodifferential operator on $\R^n$. Moreover if $a$ is elliptic then so is this pseudodifferential operator, so that by Proposition~\ref{prop:inversion_formal_pseudos} it admits an inverse with amplitude $b$.

Now, $M=I_{\psi}(a)T\Opp(b)T^*$ satisfies, for every $c\in S_{\clH}(\R^n\times S^{n-1})$,
\[
T^*MI_{\psi}(c)T=TI_{\psi}(c)T^*,
\]
and since $T$ and $T^*$ are invertible, we conclude that $M$ acts as identity on $\clH$. Moreover, by the previous considerations, $M$ can be written in the form $I_{\psi}(d)$ for some $d\in S_{\clH}(\R^n\times S^{n-1})$.

As a formal amplitude, $d$ is obtained from $a$ by a formula of the form $d=\sum_{k=0}^{+\infty}C_k(a,a_0^{-1})$, for some sequence of bidifferential operators $C_k$ (this is a consequence of Propositions~\ref{prop:T*T} and~\ref{prop:inversion_formal_pseudos}; one can also directly prove it using the $C^{\infty}$ calculus). Applying this formula to a locally defined amplitude $a$ leads necessarily its inverse for the formal product.
\end{proof}

In the last proof and from now on, given $V_0\subset \R^n\times S^{n-1}$ and given a formal analytic amplitude $a$ on $V_0\times (0,+\infty)$, we will sometimes drop the domains and consider $I_{\psi}(a)$ as the collection of operators of the form $I_{\psi,\Omega,U}(\widetilde{a}):\clF(V)\to \clH(U)$ where $U\Subset V\Subset V_0$ and $\Omega$ is a small neighbourhood of $U\times V\times (0,+\infty)$ on which $a$ extends into $\widetilde{a}$, holomorphic in the first variable and anti-holomorphic in the second variable. These operators will always strictly decrease the definition set of the function on which they act.



The operator $T^*$ sends $\clO(\R^n\times S^{n-1})$ to $\clO(\R^n)$, and therefore $T$ is well-defined from $\clO'(\R^n)$ to $\clO'(\R^n\times S^{n-1})$. Thus we can define the \emph{analytic wave front set} of an analytic functional as follows.

\begin{defi}\label{def:WFa}
The \emph{analytic wave front set} of $u\in \clO'(\R^n)$ is defined as the analytic singular support of $Tu$.
\end{defi}

This is equivalent to the usual definition of analytic singular support by means of the FBI transform, which asks for exponential decay of a related quantity. In particular, it is equivalent to all of the classical notions of analytic wave front set~\cite{bony_equivalence_1977}.

We conclude this section with a few comments. Our algebra of operators defined in Proposition~\ref{prop:invesion_Toeplitz} is inspired from ``covariant Toeplitz operators'', originally defined in the semiclassical setting in~\cite{charles_berezin-toeplitz_2003} as an alternative description of operators of the ``contravariant Toeplitz'' form $SaS$, where $S$ is the Szeg\H{o} projector and $a$ acts by multiplication; these operators also form an algebra in analytic regularity~\cite{deleporte_toeplitz_2018,rouby_analytic_2018}.

It a bit awkward that the natural range of $T$, and therefore the space on which our algebra acts, extends into a space of \emph{anti}holomorphic functions on $\R^n\times B(0,1)$. Alternative conventions, however, are not more satisfactory; one can set the FBI phase as $(y-x)\cdot
\omega+i\frac{|x-y|^2}{2}$ and obtain holomorphic functions, but at the end we obtain a characterisation of the analytic wave front set of $u$ as the analytic singular support of $(x,\omega)\mapsto Tu(x,-\omega)$. One can also simply swap the variables $x$ and $\omega$, but this is at odds with the interpretation of $\R^n\times S^{n-1}$ as the cosphere bundle over $\R^n$ with a natural CR structure obtained by identifying $T^*\R^n$ and $\widetilde{\R^n}=\C^n$; this~``Grauert tube'' vision of the FBI transform extends naturally to manifolds with a real-analytic Riemannian structure~\cite{guillemin_grauert_1991,guillemin_grauert_1992,leichtnam_intrinsic_1996}, see also the semiclassical approach~\cite{sjostrand_density_1996}.

The Grauert tube approach is essentially contained in our toolbox; the FBI transform as introduced in~\cite{leichtnam_intrinsic_1996}, for instance, can be written as an analytic Fourier integral operator whose Lagrangian has similar properties to the flat case, and the definition of the wave front set is equivalent to the one above, via analytic charts. We will say more on this in Section~\ref{sec:appl-grau-tubes}.


\subsection{Advanced properties of Fourier {integral operators}}\label{sec:advanced-properties}

With help of the FBI transform and the structure of Toeplitz operators, we can prove that general Fourier {integral operators} push the analytic wave front set of a distribution as expected, and also an microlocal ellipticity result for pseudodifferential operators.

\begin{prop}\label{prop:WF_Pu}
Let $a$ be a formal analytic symbol on $\R^n$. Then
\[
WF_a(u)\subset \{a_0=0\}\cup WF_a(\Opp(a)u).
\]
\end{prop}
\begin{proof}
Let $(x,\omega)\in \R^n\times S^{n-1}$ and suppose that $T\Opp(a)u$ is real-analytic near $(x,\omega)$ and $a_0$ is bounded away from $0$ near $(x,\omega)$. Modulo an analytic function, one~has
\[
T\Opp(a)u=T\Opp(a)\Opp(r)T^*Tu=I_{\psi}(b)Tu,
\]
where $b$ is a formal analytic amplitude such that $b_0$ is bounded away from $0$ near $(x,\omega)$. By Proposition~\ref{prop:invesion_Toeplitz}, on a neighbourhood $V$ of $(x,\omega)$, $b$ admits a formal inverse for the $I_{\psi}$ calculus. Letting $d$ denote the formal analytic amplitude of its inverse, given $U\Subset V$, one has, modulo an analytic function,
\[
I_{\psi,U\times V,U}(d)(T\Opp(a)u) = I_{\psi,U}(d)I_{\psi}(b)Tu=u;
\]
now, by Proposition~\ref{prop:reg_to_reg}, since $T\Opp(a)u$ is real-analytic on $V$, the left-hand side is real-analytic on $U$. This concludes the proof.
\end{proof}


\begin{prop}\label{prop:FIO-WF}
Let $(\phi,\Omega,U)$ satisfy~\eqref{eq:boundary_positivity} and let $a$ be a formal analytic amplitude on $\Omega$. Then
\[
WF_a(I_{\phi,\Omega,U}(a)u)\subset \left\{(x,\xi)\in T^*\R^{n_x}\setminus \{0\},\;
\exists (y,\xi)\in WF_a(u), \; (x,\xi,y,\eta)\in (\Lambda_{\phi})_{\R}\right\}.
\]
\end{prop}
\begin{proof}
The operator $T\circ I_{\phi,\Omega,U}(a)\circ T^*$ is a Fourier integral operator whose real locus of the Lagrangian projects on the base onto
\[
\left\{(x,\omega,y,v)\in \left(\R^{n_x}\times S^{n_x-1}\times \R^{n_y}\times S^{n_y-1}\right)\cap (\Lambda_{\phi})_{\R}\right\};
\]
thus, we are left with the claim
\[
SS_a(I_{\phi,\Omega,U}(a)u)\subset \left\{x\in \R^{n_x},\: \exists
\xi,\;\eta\neq 0, \:\exists y\in SS_a(u), \;(x,\xi,y,\eta)\in (\Lambda_{\phi})_{\R}\right\},
\]
which was proved in Theorem~\ref{prop:FIOs}.
\end{proof}

Another important application of the FBI transform is the construction of unitary Fourier {Integral} operators associated with arbitrary real symplectic maps, also called ``quantized contact transformations''.

\begin{prop}\label{prop:quantized-contact}
Let $\kappa$ be a one-homogeneous symplectic transformation between respective neighbourhoods of two open cones $U,V$ of $T^*\R^n$. Let $[U]=U\cap (\R^n\times S^{n-1})$ and $[V]=V\cap (\R^n\times S^{n-1})$. Then there exists $\hat{\kappa}:\clH([V])\to \clH([U])$ and $\widehat{\kappa^{-1}}:\clH([U])\to \clH([V])$ such that the following is true.
\begin{itemize}
\item\label{prop4.25.1} For every formal analytic amplitude $a\in S([V])$, there exists $b\in S([U])$, such that, for every $U_0\Subset U$, there exists $V_0\Subset V$, a neighbourhood $\Omega$ of $\diag([U_0])$ in $(\R^n\times S^{n-1})\times [U]\times (0,+\infty)$, a neighbourhood $\Omega'$ of $\diag([V_0])$ in $(\R^n\times S^{n-1})\times [V]\times (0,+\infty)$ with
\[
\hat{\kappa}I_{\psi,\Omega',[V_0]}(b)\widehat{\kappa^{-1}}=I_{\psi,\Omega,[U_0]}(a):\clH([U])\to \clH([U_0]),
\]
where $b_0=a_0\circ \kappa$; if $a$ is the amplitude of the Szeg\H{o} projector then so is $b$.
\item\label{prop4.25.2} For every formal analytic amplitude $b\in S([U])$, there exists $a\in S([V])$ such that, for every $V_0\Subset V$, there exists $U_0\Subset U$, a neighbourhood $\Omega$ of $\diag([U_0])$ in $(\R^n\times S^{n-1})\times [U]\times (0,+\infty)$, a neighbourhood $\Omega'$ of $\diag([V_0])$ in $(\R^n\times S^{n-1})\times [V]\times (0,+\infty)$with
\[
\widehat{\kappa^{-1}}I_{\psi,\Omega,[U_0]}(a)\hat{\kappa}=I_{\psi,\Omega',[V_0]}(b):\clH([V])\to \clH([V_0]),
\]
where $a_0=b_0\circ \kappa^{-1}$; if $b$ is the amplitude of the Szeg\H{o} projector then so is $a$.
\item\label{prop4.25.3} For every $U_0\Subset U$ and $V_0\Subset V$, the distribution
\[
(x,y)\mapsto
\widehat{\kappa^{-1}}(x,y)-\overline{\hat{\kappa}(y,x)}
\]
is real-analytic on $V_0\times U_0$.
\end{itemize}
\end{prop}
\begin{proof}
For convenience we drop the bracket notations and identify open sets of $\R^n\times S^{n-1}$ with open cones in $T^*\R^n$. Thus $\kappa$ is a contact transformation on $\R^n\times S^{n-1}$. We let $\phi_{\kappa}$ be any function from a neighbourhood $W$ of the graph of $\kappa$ in $V_1\times U_1$, where $U_1\Subset U$ and $V_1\Subset V$, which is holomorphic in $x_1-i\omega_1$ and in $x_2+i\omega_2$, and such that $\im(\phi_{\kappa})\asymp
\dist((x_1,\omega_1),\kappa(x_2,\omega_2))$. Such a function exists thanks to Proposition~\ref{prop:completing_Lagrangians}: it is a phase function for the only Lagrangian containing the graph of $\kappa$.

We now set
\[
\widehat{\kappa_0}=I_{\psi_{\kappa},W,V_1}\1_{(x_1,\omega_1,x_2,\omega_2)\in W}\int_0^{+\infty}e^{it\psi_{\kappa}(x_1,\omega_1,x_2,\omega_2)}t^{n-1}\dd t
\]
and we similarly define $\widehat{\kappa_0^{-1}}$ based on $\kappa^{-1}$.

Given $a$ as in the claim, the operator $\widehat{\kappa_0^{-1}}I_{\psi,\Omega,U_0}(a)\widehat{\kappa_0}$ is, by Theorem~\ref{prop:FIOs}, a Fourier integral operator with the same Lagrangian as the Szeg\H{o} projector (indeed, the real locus of the Lagrangian is the diagonal, and this Lagrangian is holomorphic for the skew CR structure on $(\R^n\times S^{n-1})^2$). Therefore it is of the form $I_{\psi,\Omega',[V_0]}(b)$ for another analytic formal amplitude $b$ obtained from $a$ by stationary phase. In particular, letting $S$ denote the Szeg\H{o} projector,
\[
\widehat{\kappa_0^{-1}}\widehat{\kappa_0}=\widehat{\kappa_0^{-1}}S\widehat{\kappa_0}=I_{\psi,\Omega',V_1}(f),
\]
where $f_0\neq 0$ and $V_1$ is any relatively compact open set of $V$. By Proposition~\ref{prop:invesion_Toeplitz}, $f$~admits an inverse for the product of analytic Fourier {integral operators} with phase~$\psi$, and therefore there exists a formal analytic amplitude $g$ on a neighbourhood of~$V$, such that
\[
I_{\psi,\Omega',V_1}(g)\widehat{\kappa_0^{-1}}\widehat{\kappa_0}=S:\clH(V)\to
\clH(V_1).
\]
We now set $\hat{\kappa}=\widehat{\kappa_0}I_{\psi,
\Omega',V_1}(\sqrt{g})$ and $\widehat{\kappa^{-1}}=I_{\psi,\Omega',V_1}(\sqrt{g}^*)\widehat{\kappa_0^{-1}}$. In this algebra, the square root and adjoint of an amplitude are well-defined because, as in the proof of Proposition~\ref{prop:invesion_Toeplitz}, the calculus of these amplitudes is conjugated via $T$ to the symbol calculus of pseudo-differential operators. With this correction, we obtain the\linebreak last claim.

We also obtain that the second part of the claim holds with the quantifiers on $a$ and $b$ exchanged (for every $a$, there exists $b, \dots)$.

Given $b$, the operator
\[
\hat{\kappa}I_{\psi,\Omega',V_0}(b)\widehat{\kappa^{-1}}
\]
is also an analytic Fourier integral operator whose Lagrangian is that of $\psi$, so that it is of the form $I_{\psi,\Omega,U_0}(a)$, and we obtain the first part of the claim with the quantifiers on $a$ and $b$ exchanged (for every $b$, there exists $a, \dots)$.

The particular case $\hat{\kappa}S\widehat{\kappa^{-1}}$ is a projector, because we already made sure that $\widehat{\kappa^{-1}}\hat{\kappa}$ acts as identity; however it is also an invertible operator (because it is a Fourier integral operator with phase $\psi$ and nonvanishing principal amplitude). Therefore it is the identity.

To conclude, given $a\in S(V)$, one recovers $b$ so that the first point of the proposition is true by letting $b$ be the amplitude of $\widehat{\kappa^{-1}}I_{\psi,\Omega,U_0}(a)\hat{\kappa}$, and conversely.
\end{proof}

\begin{rema}\label{rem:replacement}
To conclude this section, we can go back to Remark~\ref{rem:EtoO}. Now that we have characterised the analytic wave front set of the singular kernel of an analytic Fourier integral operator (Proposition~\ref{prop:FIO-WF}), we observe that these operators are well-defined when acting on singularity hyperfunctions, that is, elements of $\clF_{\omega}(V):=\clO'(\bar{V})/(\clO'(\partial V)+ \clO(\bar{V}))$. Indeed, as in the smooth case the wave front set conditions (see~\cite[Theorems~8.5.1 and~8.5.2]{hormander_analysis_2003}) are satisfied. Thus, in all of Section~\ref{sec:four-integr-oper}, we can now replace $\clF$ with $\clF_\omega$, and $\clH$ with $\clH_\omega:U\mapsto
\ker_{\clO'(\bar{U})}(\bar{\partial}_b)/(\clO(\bar{U})+\clO'(\partial U))$. This replacement is crucial in the proof of Theorem~\ref{prop:analytic-Szego}.
\end{rema}

\section{The Szeg\H{o} kernel and Toeplitz operators for general pseudoconvex open sets}\label{sec:szegho-kern-toepl}



\subsection{Normal form for the \texorpdfstring{$\overline{\partial}_b$}{partial\_b} operator}\label{sec:norm-form-overl}

The features of $\bar{\partial}_b$ on the boundary of a strongly pseudoconvex manifold are the following: locally, it is a vector-valued, degree~$1$ analytic pseudodifferential operator $(D_1,\dots,D_d)$ on a $2d-1$-dimensional manifold, such that $[D_j,D_k]=0$ for every $j\neq k$ and the matrix of principal symbols $\sigma(i[D_j,D_k^*])$ is either positive definite or negative definite on $\Sigma=\{\sigma(D_j)=0\forall 1\leq j\leq d\}$.

All of this is also true of the following model operator $D_0$, acting on $\R^{d-1}_x\times \R^{d}_y$:
\[
(D_0)_j=-i\frac{\partial}{\partial x_j}-x_j\frac{\partial}{\partial y_1}.
\]

Our goal is to prove the following result attributed to Sato: microlocally near any point of $\Sigma$, $\bar{\partial}_b$ is conjugated with $D_0$ by an analytic Fourier integral operator (as always, modulo an operator continuously mapping $\clE'$ into real-analytic-functions). Here, as before, ``microlocally'' will mean ``locally after conjugation with the FBI transform''.


The first step (Proposition~\ref{prop:canon_transf_to_local_model}) is to solve the classical problem, and find a symplectic transformation which maps the symbol of $\bar{\partial}_b$ to the symbol of $D_0$, multiplied by an invertible matrix; this was described in Section~\ref{sec:class-norm-form}. W move on to an exact conjugation in the algebra of analytic formal symbols of pseudodifferential operators (Proposition~\ref{prop:correct_D0_error_formal}), and we conclude by realising all requested transformations on the FBI side (Proposition~\ref{prop:Sato_conj}).



Given a unitary Fourier integral operator $\hat{\kappa}$ which quantizes $\kappa$, the principal symbol of $\widehat{\kappa^{-1}}\bar{\partial}_b\hat{\kappa}$ is precisely that of $\Opp(C)D_0$. It remains to correct the subprincipal (degree~$0$) remainder, which we can do by conjugating with an elliptic, diagonal matrix-valued pseudodifferential operator.

\begin{prop}\label{prop:correct_D0_error_formal}
Let $r=(r_1,\dots,r_d)$ be a formal analytic degree~0 vector-valued symbol in the neighbourhood of a point where $\sigma(D_0)=0$. Then there exists a smaller neighbourhood of this point and a vector of degree~0 elliptic symbols $a=(a_1,\dots,a_d)$ such that, on this neighbourhood,
\[
(\sigma(D_0)_j+r_j)\sharp a_j=a_j\sharp \sigma(D_0)_j, \qquad
\forall 1\leq j\leq d.
\]
\end{prop}

\begin{proof}
We first remove the order 0 term in $r$, then all other terms.

At order 0, the following equation links the principal symbol of $a_j$ to the principal symbol of $r_j$
\[
\partial_{x_j}a_j^0+ix_j\partial_{y_1}a_j^0-i\eta_1\partial_{\xi_j}a_j^0=ir_j^0a_j^0.
\]
This transport equation admits a solution, together with the Cauchy data $a_j^0|_{x_j=0}=1$, written as
\[
a_j^0(x,y,\xi,\eta)=\exp\left[i\int_0^{x_j}r_j^0\left(x-(1-t)e_j,y+i\frac{t^2}{2}e_1,\xi+i(1-t)\eta_1e_j,\eta\right)\dd t\right].
\]
Since $r_j^0$ is a real-analytic function, this solution is well-defined for $x_j$ sufficiently close to $0$, and is again real-analytic. Hence, we can solve for $a$ at principal order. Thus
\[
(a^0_j)^{\sharp -1}\sharp \left(\sigma(D_0)_j+r_j\right)\sharp a_j^0=\sigma(D_0)_j+r_j',
\]
where $r_j'$ is a degree $-1$ analytic symbol.

Consider now the family
\[
[0,1]\ni t\mapsto \sigma(D_0)_j+tr_j'.
\]
We wish to find a degree $-1$ elliptic symbol $a_j'(t)$ such that
\begin{equation}\label{eq:rprimeaprime}
\left(\sigma(D_0)_j+tr_j'\right)\sharp\left(1+a_j'(t)\right)=\left(1+a_j'(t)\right)\sharp (D_0)_j.
\end{equation}

We will find $a_j'(t)$ as a solution of the differential equation
\[
\partial_ta_j'(t)=ib_j(t)\sharp\left(1+a_j'(t)\right).
\]
First, differentiating~\eqref{eq:rprimeaprime} with respect to $t$ yields
\[
r_j'=-i\left[\sigma(D_0)_j+tr_j',b_j(t)\right],
\]
and we want to solve this equation for $b_j(t)$.

One has
\begin{equation}\label{eq:exact_conj_D0}
\left[(D_0)_j,\Opp(b_j(t))\right]
=\Opp\left(-i\partial_{x_j}b_j(t)-x_j\partial_{y_1}b_j(t)+\eta_1\partial_{\xi_j}b_j(t)\right).
\end{equation}
From now on we drop the index $j$ and the parameter $t$ on $r'$ and $b$. We write their full symbols as
\begin{align*}
r'=\sum_{k\geq 1} r_k'\qquad b=\sum_{k\geq 1}b_k,
\end{align*}
where $b_k$ and $r_k'$ are homogeneous of degree $-k$.

The equation above yields the following family of equations on $b_k$:
\[
-i\partial_{x_j}b_k +\eta_1\partial_{\xi_j}b_k-x_j\partial_{y_{1}}b_k=ir_k'+t
\left(r'\sharp b-b\sharp r'\right)_k.
\]
Note that the right-hand-side only involves $b_l$ for $1\leq l\leq k-1$. Thus, we can find the $b_k$ by induction, recursively solving a transport equation of the form
\[
i\partial_{x_j}u-\eta_{1}\partial_{\xi_j}u+x_j
\partial_{y_{1}}u=g,
\]
using again the method of characteristics: the solution (with Cauchy data $u=0$ on $x_j=0$) is
\[
u:(x,y,\xi,\eta)\mapsto
\int_0^{x_j}g\left(x-(1-t)e_j,\;y+i\frac{t^2}{2}e_1,\;\xi+i(1-t)\eta_1e_j,\eta\right)\dd t.
\]

We have to prove that the successive solutions $(b_k)_{k\in \N}$ form an analytic symbol. It turns out that a convenient analytic class is given by the infinite jet of the $S^{\rho,R}_m$ class at $x_j=0$; namely we let
\[
\|a\|_{JS^{\rho,R}_m}=\sup_{\alpha,k}\frac{\sup_\Omega(|\nabla^{\alpha}a_k|)(1+|\alpha|+k)^m}{(|\alpha|+k)!\rho^{|\alpha|}R^k},
\]
where
\[
\Omega=\left\{|\eta_{1}|\in \left[\frac 14,\; \frac 12\right],\;|x|+|y|+|\xi|+|(\eta_2,\dots,\eta_{d-1})|<\epsilon,\quad x_j=0
\right\}.
\]
Since the product of symbols is an infinitesimal operation, we have (see Proposition~\ref{prop:composition-formal})
\[
\left\|(b\sharp r'-r'\sharp b)\right\|_{JS^{\rho,R}_m}\leq 24\|r'\|_{JS^{\frac{\rho}2, \frac R4}_m}\|b\|_{JS^{\rho,R}_m}.
\]

Suppose by induction that, for all $l\leq k-1$, one has
\[
\sup_{\Omega}\|\nabla^{\alpha}b_l\|\leq C_b\frac{\rho^{|\alpha|}R^l(|\alpha|+l)!}{(|\alpha|+l+1)^m}.
\]
Then one has readily, for all $t\in [0,1]$, if $m\geq C(d)$,
\[
\sup_{\Omega}\left|\nabla^\alpha\left(r'_k+t\left(r'\sharp b - b \sharp r'\right)_k\right)\right|\leq (24C_b+1)\|r\|_{JS^{\frac{\rho}2,\frac R4}_m}\frac{\rho^{|\alpha|}R^k(|\alpha|+k)!}{(|\alpha|+k+1)^m}.
\]
Denote $g_k=ir'_k+t(r'\sharp b - b \sharp r')_k$, and let
\begin{align*}
C(g_k,k) &=\sup_{\alpha}\frac{\sup_\Omega(|\nabla^{\alpha}g_k|)(1+|\alpha|+k)^m}{(|\alpha|+k)!\rho^{|\alpha|}R^k}
\\
\intertext{and}
C(b_k,k) &=\sup_{\alpha}\frac{\sup_\Omega(|\nabla^{\alpha}b_k|)(1+|\alpha|+k)^m}{(|\alpha|+k)!\rho^{|\alpha|}R^k}.
\end{align*}
We just proved that
\[
C(g_k,k)\leq \|r'\|_{JS^{\frac{\rho}2, \frac R4}_m}\left(1+24\sup_{ \ell
\leq k-1}C(b_\ell,\ell)\right).
\]
Let us now prove that
\[
C(b_k,k)\leq C(g_k,k).
\]
To this end, we write the transport equation
\[
i\partial_{x_j}b_k-\eta_{1}\partial_{\xi_j}b_k+x_j
\partial_{y_{1}}b_k=g_k
\]
in terms of the power series at $x_j=0$
\begin{align*}
b_k=\sum_{n=1}^{+\infty} \frac{b_k^{(n)}}{n!}x_j^n  \qquad g_k=\sum_{n=1}^{+\infty}
\frac{g_k^{(n)}}{n!}x_j^n.
\end{align*}
One has, for all $n\geq 1$,
\begin{equation}\label{eq:transport}
ib_{k}^{(n+1)}-\eta_{1}\partial_{\xi_j}b_{k}^{(n)}+n\partial_{y_{1}}b_{k}^{(n-1)}=g_k^{(n)},
\end{equation}
and for $n=0$,
\[
b_{k}^{(0)}=0.
\]

Let us prove by induction on $n$ that
\[
\left|\nabla^{\alpha}b_k^{(n)}\right|\leq C(g_k,k)\frac{\rho^{n+|\alpha|}(n+|\alpha|+k)!}{(n+|\alpha|+k+1)^m}.
\]
This is true at $n=0$. If it is true at $n-1$ and $n$, then there are three terms in
\begin{equation}\label{eq:ctrl_b_g}
\left|\nabla^{\alpha}b_k^{(n+1)}\right|\leq\left|\nabla^{\alpha}g_k^{(n)}\right|+\left|\nabla^{\alpha}\left(\eta_{1}\partial_{\xi_j}b_k^{(n)}\right)\right|-
\left|n\nabla^{\alpha}\partial_{y_{1}}b_k^{(n-1)}\right|.
\end{equation}
First, one has directly
\begin{align*}
\left|\nabla^{\alpha}g_k^{(n)}\right|
&\leq C(g_k,k)\frac{\rho^{(n+|\alpha|)}R^k(n+|\alpha|+k)!}{(n+|\alpha|+k+1)^m}\\
&\leq C(g_k,k)\frac{\rho^{(n+1+|\alpha|)}R^k(n+1+|\alpha|+k)!}{(n+|\alpha|+k+2)^m}\times
\frac{\left(1+\frac{1}{n+|\alpha|+k+1}\right)^m}{\rho(n+1+|\alpha|+k)}.
\end{align*}
If $\rho\geq 6\left(\frac{3}{2}\right)^m$, then we obtain
\[
\left|\nabla^{\alpha}g_k^{(n)}\right|\leq
\frac{1}{6}C(g_k,k)\frac{r^{(n+1+|\alpha|)}R^k(n+1+|\alpha|+k)!}{(n+|\alpha|+k+2)^m}.
\]
We can also directly control the last term in~\eqref{eq:ctrl_b_g} in the same way, using the induction hypothesis:
\begin{multline*}
\left|n\nabla^{\alpha}\partial_{y_{d+1}}b_k^{(n-1)}\right|\leq nC(g_k,k)\frac{\rho^{n+|\alpha|}R^k(n+|\alpha|+k)!}{(n+|\alpha|+k+1)^m}\\
\leq C(g_k,k)\frac{\rho^{n+1+|\alpha|}R^k(n+1+|\alpha|+k)!}{(n+|\alpha|+k+2)^m}\times
\frac{\left(1+\frac{1}{n+|\alpha|+k+1}\right)^m}{\rho}\frac{n}{n+1+|\alpha|+k}.
\end{multline*}
As previously, if $\rho\geq 6\left(\frac 32\right)^m$, we obtain
\[
\left|n\nabla^{\alpha}\partial_{y_{1}}b_k^{(n-1)}\right|\leq
\frac{1}{6}C(g_k,k)\frac{\rho^{(n+1+|\alpha|)}R^k(n+1+|\alpha|+k)!}{(n+|\alpha|+k+2)^m}.
\]

In the second term of~\eqref{eq:ctrl_b_g}, we have to distinguish between two cases, depending on whether or not one of the differentials has hit $\eta_1$. Thus
\[
\left|\nabla^{\alpha}\left(\eta_{1}\partial_{\xi_j}b_k^{(n)}
\right)\right|\leq |\alpha|\:\left|\nabla^{\alpha-\gamma}\partial_{\xi_j}b_k^{(n)}\right|+|\eta_{1}|\:\left|\nabla^{\alpha}\partial_{\xi_j}b_k^{(n)}\right|,
\]
where $\gamma$ is the length $1$ polyindex corresponding to the direction of $\eta_{1}$.

As before,
\begin{multline*}
|\alpha|\left|\nabla^{\alpha-\gamma}\partial_{\xi_j}b_k^{(n)}\right|\leq |\alpha|C(g_k,k)\frac{\rho^{n+|\alpha|}R^k(n+|\alpha|+k)!}{(n+|\alpha|+k+1)^m}\\
\leq C(g_k,k)\frac{\rho^{n+1+|\alpha|}R^k(n+1+|\alpha|+k)!}{(n+|\alpha|+k+2)^m}\times
\frac{\left(1+\frac{1}{n+|\alpha|+k+1}\right)^m}{\rho}\frac{|\alpha|}{n+1+|\alpha|+k},
\end{multline*}
and if $\rho\geq 6\left(\frac 32\right)^m$, we obtain
\[
|\alpha|\left|\nabla^{\alpha-\gamma}\partial_{\xi_j}b_k^{(n)}\right|\leq
\frac{1}{6}C(g_k,k)\frac{\rho^{(n+1+|\alpha|)}R^k(n+1+|\alpha|+k)!}{(n+|\alpha|+k+2)^m}.
\]

The last term is
\[
|\eta_{1}|\left|\nabla^{\alpha}\partial_{\xi_j}b_k^{(n)}\right|
\leq |\eta_{1}|C(g_k,k)\frac{\rho^{(n+1+|\alpha|)}R^k(n+1+|\alpha|+k)!}{(n+|\alpha|+k+2)^m}.
\]
Recall that we are controlling derivatives on the set $\Omega$ where $|\eta_1|\leq \frac 12$. We obtain
\[
|\eta_{1}|\left|\nabla^{\alpha}\partial_{\xi_j}b_k^{(n)}
\right|\leq
\frac 12 C(g_k,k)\frac{\rho^{(n+1+|\alpha|)}R^k(n+1+|\alpha|+k)!}{(n+|\alpha|+k+2)^m}.
\]
Altogether, the induction is complete, and we obtain that
\[
C(b_k,k)\leq C(g_k,k)\leq \left(24\max_{\ell\leq k-1}C(b_{\ell},\ell)+1\right)\|r'\|_{JS^{\frac{\rho}2, \frac R4}_{m}},
\]
provided
\begin{align*}
m\geq C(d),\qquad\rho\geq 6\left(\frac 34\right)^m,\qquad R\geq 2^{d+1}\rho.
\end{align*}
It remains to find $\rho,R,m$ as above such that $\|r'\|_{S^{\rho,R}_m}<\frac{1}{24}$. Since $r'$ is of degree $-1$, this can be done by fixing $m$ and $\rho$, then choosing $R$ large enough.

Having found $b$, it remains to build $a'$ such that
\begin{align*}
&(1+a')(0)=0,& &\partial_t (1+a')(t) = b(t)\sharp (1+a')(t).
\end{align*}
We know that $b(t)\in JS^{\rho,R}_m$ for all $t\in [0,1]$, with uniformly bounded norm. Thus, by Proposition~\ref{prop:composition-formal}, this Cauchy problem satisfies the hypotheses of the Picard--Lindel\"of theorem on the Banach space $JS^{2\rho,4R}_m$. In particular, we can find a solution in this space, and $a'(1)$ is the requested operator.
\end{proof}

Grouping together Propositions~\ref{prop:correct_D0_error_formal}, \ref{prop:canon_transf_to_local_model} and~\ref{prop:quantized-contact}, Theorem~\ref{prop:FIOs} and Remark~\ref{rem:replacement}, we obtain the following quantum normal form.

\begin{prop}\label{prop:Sato_conj}
Let $(z,\zeta)\in \Sigma^{\pm}$. Identify a neighbourhood of $z$ in $X$ with a neighbourhood of $0$ in $\R^{2d-1}$. There exists an open neighbourhood $U$ of $(z,\zeta)$ in $T^*\R^{2d-1}$ such that, for every pair of open sets $z\in U_2\Subset U_1 \Subset U$, there exists:
\begin{itemize}
\item open neighbourhoods $V_2\Subset V_1$ of $\{x=y=\xi=0,\eta=(\pm 1, 0,\dots, 0)\}$ in $T^*(\R^d_x\times \R^{d-1}_y)$,
\item an analytic Fourier integral operator $\hat{\kappa}:\clH_{\omega}(U_1)\to\clH_{\omega}(V_1)$,
\item an analytic Fourier integral operator $\widehat{\kappa^{-1}}:\clH_{\omega}(V_2)\to \clH_{\omega}(U_2)$,
\item an elliptic matrix-valued formal symbol $C\in S(V_1)$,
\item a neighbourhood $\Omega$ of $\diag(V_1)\times (0,+\infty)$ in $(\R^{2d-1}\times S^{2d-2})^2\times (0,+\infty)$,
\end{itemize}
such that
\[
T\bar{\partial}_bT^*=\widehat{\kappa^{-1}}I_{\psi,\Omega,V_2}(C)TD_0T^*\widehat{\kappa},
\qquad S=\widehat{\kappa^{-1}}\hat{\kappa},
\]
as operators mapping $\clF_{\omega}(U_1)$ into $\clF_{\omega}(U_2)$.

Reciprocally, there exists an open neighbourhood $V$ of $\{x=y=\xi=0,\eta=(\pm 1, 0,\dots, 0)\}$ in $T^*(\R^d_x\times \R^{d-1}_y)$ such that, for every pair of open sets $V_2\Subset V_1\Subset V$ containing this point, there exists:
\begin{itemize}
\item open neighbourhoods $U_2\Subset U_1$ of $(z,\zeta)$ in $T^*R^{2d-1}$,
\item an analytic Fourier integral operator $\hat{\kappa}:\clH_{\omega}(U_2)\to\clH_{\omega}(V_2)$,
\item an analytic Fourier integral operator $\widehat{\kappa^{-1}}:\clH_{\omega}(V_1)\to \clH_{\omega}(U_1)$,
\item an elliptic matrix-valued formal symbol $C\in S(V_1)$,
\item a neighbourhood $\Omega$ of $\diag(V_1)\times (0,+\infty)$ in $(\R^{2d-1}\times S^{2d-2})^2\times (0,+\infty)$,
\end{itemize}
such that
\[
\widehat{\kappa}T\bar{\partial}_bT^*\widehat{\kappa^{-1}}=I_{\psi,\Omega,V_2}(C)TD_0T^*, \qquad
\qquad S=\hat{\kappa}\widehat{\kappa^{-1}},
\]
as operators mapping $\clF_{\omega}(V_1)$ into $\clF_{\omega}(V_2)$.
\end{prop}


\subsection{The Szeg\H{o} kernel parametrix and the algebra of Toeplitz operators}\label{sec:szegho-kern-param}

In this section, we broadly follow the method proposed in~\cite{boutet_de_monvel_sur_1975} to obtain a parametrix for the Szeg\H{o} kernel starting with the normal form of Proposition~\ref{prop:Sato_conj}.

In this section, the index $_0$ refers to the ``model'' case $\R^d\times S^{2d-1}$. In particular $S_0$ is the Szeg\H{o} projector of Proposition~\ref{prop:flat_Grauert}, and $\Sigma^{\pm}_0$ is the positive, resp. negative, characteristic set of $(\bar{\partial}_b)_0$, the boundary Cauchy--Riemann operator on $\R^d\times S^{2d-1}$. By contrast, the notations $\bar{\partial}_b, S,
\Sigma^{\pm},$ refer to the boundary of a general, relatively compact, strongly pseudo-convex open set whose boundary is denoted by $X$.


\begin{theo}\label{prop:analytic-Szego}
The Szeg\H{o} projector
admits a (singular) integral kernel which is real-analytic away from the diagonal. Given a defining function $\rho$ and its polarisation $\psi$, well-defined on a small neighbourhood $\Omega$ of the diagonal in $X\times X$, there exists a realisation of a formal analytic amplitude $a$ on $X\times X\times (0,+\infty)$ (see Definitions~\ref{def:formal-amp}, \ref{def:analytic_amp} and Proposition~\ref{prop:real_formal_symb} for details about the nature of $a$) such that
\[
(x,y)\mapsto S(x,y)-\int_0^{+\infty}e^{it\psi(x,y)}a(x,y;t)\dd t
\]
is real-analytic.
\end{theo}

\begin{proof}
The proof consists first in exhibiting a parametrix for
$S$ in the form of a Fourier integral operator as above, then proving that this parametrix indeed approximates $S$ in the sense that the difference is a continuous operator from $\clO'(X)$ to $\clO(X)$. This requires several steps, as the mere fact that $S$ acts nicely on $\clO'(X)$ requires a proof.
\begin{enumerate}
\item\label{prftheo5.3.pt1} Given $(z,\zeta)\in \Sigma^{\pm}$ and $(z_0,\zeta_0)\in \Sigma^{\pm}_0$, applying Proposition~\ref{prop:Sato_conj} twice, we obtain a pair of Fourier integral operators $\hat{\kappa}$ and $\widehat{\kappa^{-1}}$ between neighbourhoods of $(z,\zeta)$ and $(z_0,\zeta_0)$ (such that $\hat{\kappa}\widehat{\kappa^{-1}}$ and $\widehat{\kappa^{-1}}\hat{\kappa}$ slightly decrease the definition set and are equal to the restriction map) such that, for every pair of small neighbourhoods $V_1\Subset V$ of $(z,\zeta)$ and every pair of small neighbourhoods $U_1\Subset U$ of $(z_0,\zeta_0)$,
\begin{align*}
\hat{\kappa}T\bar{\partial}_bT^*\widehat{\kappa^{-1}}&=I_{\psi}(C)T(\bar{\partial}_b)_0T^*:\clF_{\omega}(U)\to \clH_{\omega}(U_1)\\
T\bar{\partial}_bT^*&=\widehat{\kappa^{-1}}I_{\psi}(C)T(\bar{\partial}_b)_0T^*\hat{\kappa}:\clF(V)\to \clH(V_1).
\end{align*}
We now present the following local candidate for $TST^*$:
\[
\clS=\widehat{\kappa^{-1}}TS_0T^*\hat{\kappa}.
\]
One readily checks that
\[
T\bar{\partial}_bT^*\clS=0:\clF_{\omega}(V)\to
\clH_{\omega}(V_1).
\]
Moreover $\clS$ is self-adjoint, because of the constraints on $\widehat{\kappa}$ and $\widehat{\kappa^{-1}}$ in Proposition~\ref{prop:quantized-contact} (modified by Remark~\ref{rem:replacement}).

\item\label{prftheo5.3.pt2} From the definition of $\clS$, we see that $T^*\clS T$ is an elliptic Fourier integral operator whose Lagrangian $\Lambda$ is a half-line bundle over its projection onto the base; it is also a projector. When writing $T^*\clS T$ in a form given by Proposition~\ref{prop:FIO-rank1-phase}, once the phase is chosen, the fact that it is a self-adjoint projector determines every term of the formal amplitude. Therefore different pieces of $T^*\clS T$ can be glued together into a global Fourier integral operator $\widetilde{S}$ on $X$; in fact, this Fourier integral operator takes the form proposed in the statement, with a formal analytic amplitude.

\item\label{prftheo5.3.pt3} Let us prove that $(\widetilde{S}-1)S$ is continuous from $L^2(X)$ to $\clO(X)$. Given $u\in L^2(X)$, the wave front of $Su$ is included in $\Sigma$ by Proposition~\ref{prop:WF_Pu} because $\bar{\partial}Su=0$. Now, given $(z,\zeta)\in \Sigma$ and small open neighbourhoods $V_1\Subset V$, letting $v$ be the restriction of $TSu$ to $V$, we can apply the previous technique and obtain that $\widehat{\kappa^{-1}}T(\bar{\partial}_b)_0T^*\hat{\kappa}v=0$ as an element of $\clF_{\omega}(V_1).$ By Proposition~\ref{prop:loc-anal-hypoell-Grauert}, we deduce that $\widehat{\kappa^{-1}}T(S_0-1)T^*\hat{\kappa}v=0$ as an element of $\clF_{\omega}(V_1)$; but this is equal to the restriction of $(\widetilde{S}-1)S u$ to $V_1$.

\item\label{prftheo5.3.pt4} Let us similarly prove that $\widetilde{S}(1-S)$ is continuous from $L^2(X)$ to $\clO(X)$. By the results of Kohn~\cite{kohn_harmonic_1963} and Boutet de Monvel~\cite{boutet_de_monvel_hypoelliptic_1974}, there exists $F:L^2(X)\to L^2(X)$, such that $I_{L^2}=S+F\bar{\partial}_b$. Therefore, as operators on $L^2(X)$, one has
\[
\widetilde{S}=\widetilde{S}S+(F\bar{\partial}_b\widetilde{S})^*=\widetilde{S}S+\widetilde{S}(\bar{\partial}_b)^*F^*.
\]
As functions of the second variable, the phase and all the terms of the formal amplitude of $\widetilde{S}$ belong to the kernel of $\partial_b$, and therefore, integrating by parts, for every $u\in L^2(X)$, $\widetilde{S}(\bar{\partial}_b)^*u\in C^{\omega}(X)$.

\item\label{prftheo5.3.pt5} Since both $\widetilde{S}S-S$ and $\widetilde{S}-\widetilde{S}S$ are continuous from $L^2(X)$ to $\clO(X)$, we deduce that $\widetilde{S}-S$ is continuous from $L^2(X)$ to $\clO(X)$. In particular, $S=\widetilde{S}+(S-\widetilde{S})$ is continuous from $\clO(X)$ to $\clO(X)$, and therefore by duality $S$ acts continuously on $\clO'(X)$.

\item\label{prftheo5.3.pt6} Using~\eqref{prftheo5.3.pt5}, we can repeat the proof of~\eqref{prftheo5.3.pt3} starting with $u\in
\clO'(X)$, and we now obtain that $(\widetilde{S}-1)S$ is continuous from $\clO'(X)$ to $\clO(X)$.

\item\label{prftheo5.3.pt7} Using~\eqref{prftheo5.3.pt5} again, we obtain that, by duality, $S-\widetilde{S}$ is continuous from $\clO'(X)$ to $L^2(X)$, and therefore $(S-\widetilde{S})^2=S-\widetilde{S}S-S\widetilde{S}+\widetilde{S}$ is continuous from $\clO'(X)$ to $\clO(X)$. Then, by~\eqref{prftheo5.3.pt6}, both $\widetilde{S}S-S$ and its dual $S\widetilde{S}-S$ are continuous from $\clO'(X)$ to $\clO(X)$. Therefore $S-\widetilde{S}$ is continuous from $\clO'(X)$ to $\clO(X)$, and the proof of Theorem~\ref{prop:analytic-Szego} is complete.\qedhere
\end{enumerate}
\end{proof}

An interesting consequence of Theorem~\ref{prop:analytic-Szego} is the
\emph{analytic hypoellipticity} of $\bar{\partial}_b$: if $u\in
\clO'(X)$ is such that $\bar{\partial}_bu$ is real-analytic on an open set $V$, then $(1-S)u$ is real-analytic on $V$; in fact, we obtain the much stronger property
\[
WF_a((1-S)u)= WF_a(\bar{\partial}_bu).
\]
This seems to be a new result, interesting in its own right. Away from the characteristic $\Sigma$ this property is of course the usual ellipticity result (Proposition~\ref{prop:WF_Pu}), and there are now several proofs of microlocal analytic hypoellipticity on $\Sigma_-$~\cite{tartakoff_elementary_1981,treves_analytic_1978,sjostrand_analytic_1983}.

We are now able to generalise the construction of Section~\ref{sec:pseud-oper-1} to define the algebra of Toeplitz operators on $X$.

\begin{prop}\label{prop:Toeplitz_general}
Let $S(X)$ denote the space of formal analytic amplitudes on $X\times (0,+\infty)$. After a holomorphic extension, these amplitudes lead to Fourier integral operators acting on $\clF_{\omega}(X)$, with singular kernel of the form
\[
I_{\psi}(a):(x,y)\mapsto
\int_0^{+\infty}e^{-t\psi(x,y)}t^{n-1}\widetilde{a}(x,y;t)\dd t, \qquad a\in S(X).
\]
\begin{enumerate}
\item\label{prop5.4.1} These operators form an algebra under composition; the product law takes the form
\[
(a\sharp b)_{k}=\sum_{n+\ell=0}^kB_n(a_{\ell},b_{k-n-\ell}),
\]
where $(B_n)_{n\in \N}$ is a sequence of bilinear differential operator such that, for all $n\in \N$, $B_n$ is of total degree~$2n$.

\item\label{prop5.4.2} The principal symbol of $[I_{\psi}(a),I_{\psi}(b)]$ is $\{a,b\}$ (using the natural symplectic structure on $X\times (0,+\infty)$).

\item\label{prop5.4.3} If the degree of $a$ is $d$, then $I_{\psi}(a)$ continuously sends $H^s(X)$ into $H^{s-d}(X)$ for every $s\in \R$.

\item\label{prop5.4.4} The Szeg\H{o} projector $S=I_{\psi}(s)$ is the unit for this algebra.
\item\label{prop5.4.5} Invertible elements are exactly those for which the principal symbol $a_0$ is bounded away from $0$.
\item\label{prop5.4.6} The space $S(X)$ is the union of Banach spaces on which $a,b\mapsto a\sharp b$ is continuous.
\end{enumerate}
\end{prop}

\begin{proof}
Except for the last two items, all claims follow from Theorem~\ref{prop:FIOs}, Theorem~\ref{prop:analytic-Szego}, the usual properties of Fourier integral operators in smooth regularity, and the results of~\cite{boutet_de_monvel_spectral_1981}. In particular, if $P$ is a (smooth or analytic) pseudodifferential operator on $X$, then $\mathit{SPS}$ is a (smooth or analytic) Fourier integral operator with phase $\psi$ and its principal symbol is the restriction of the principal symbol of $P$ to $\Sigma_+$ (identified with $X\times (0,+\infty)$ by the section~$x\mapsto \partial \psi(x,x)$), so we can apply~\cite[Chapter~2, Corollary~2]{boutet_de_monvel_spectral_1981} to compute commutators at main order.


To prove invertibility of elliptic amplitudes and continuity of the product in suitable spaces of formal analytic amplitudes, we use again Proposition~\ref{prop:Sato_conj}: the wave front set of these operators is $\{(z,\zeta,z,\bar{\zeta}),(z,\zeta)\in \Sigma^+\}$ and microlocally near any point of $\Sigma^+$, a Fourier integral operators conjugates $I_{\psi}(a)$, where $a$ is elliptic, with an operator of the form of Proposition~\ref{prop:invesion_Toeplitz} with elliptic symbol. The latter can be inverted, and conjugating back we obtain, locally, an amplitude for the inverse. Since the composition rule is local,
\end{proof}

One can quantize homogeneous contact transformations as in the flat case; an important subclass of these transformations are fundamental solutions to the Schr\"od\-inger equation for degree~$1$ Toeplitz operators.

\begin{prop}
Given a one-homogeneous symplectic transformation $\kappa:X\times (0,+\infty)\to X\times (0,+\infty)$, there exists a neighbourhood $\Omega$ of the graph of $[\kappa]:X\to X$, a function $\psi:\Omega\to \C$, CR-holomorphic in the first variable, anti-CR-holomorphic in the second variable, such that $\im(\psi)\equiv \dist(\cdot,\Graph([\kappa]))^2$. In particular, the function $\Omega\times (0,+\infty)\ni (x,y,t)\mapsto t\psi(x,y)$ is a positive phase function, whose Lagrangian $\Lambda$ is CR-holomorphic in the first variable, anti-CR-holomorphic in the second variable, and such that $\pi(\Lambda_{\R})=\Graph([\kappa])$.
\end{prop}
\begin{proof}
We recover $\psi$ from its Lagrangian, which is unique by Proposition~\ref{prop:completing_Lagrangians}. By design $\Lambda_{\R}$ is a half-line bundle over its projection, and this condition is open, so that $\Lambda$ satisfies the hypotheses of Proposition~\ref{prop:FIO-rank1-phase}. This concludes the proof.
\end{proof}

\pagebreak
\section{A few applications}\label{sec:few-applications}



\subsection{Grauert tubes, FBI transforms, and quantum propagators}\label{sec:appl-grau-tubes}

Let $(M,g)$ be a compact, real-analytic Riemannian manifold. The holomorphic extension of the exponential map gives a natural isomorphism between small neighbourhoods of $M$ in $\widetilde{M}$ and small neighbourhoods of the zero section in $T^*M$.

There exists~\cite{guillemin_grauert_1991} a real-analytic K\"ahler structure on a neighbourhood of the zero section in $T^*M$, compatible with the complex structure (via the aforementioned identification) and whose restriction to $M$ is $g$. The \emph{Grauert tube}
\[
B_r=\{(x,\xi)\in T^*M, \;\|\xi\|_{g(x)}<r\}
\]
is then a strongly pseudoconvex open set for $r$ small, and its boundary $X_r$ is a real-analytic submanifold of $T^*M$.

The Toeplitz operators on $X_r$ acquire a particular importance because of the FBI transform (from $\clO'(M)$ to $\clH_{\omega}(X_r)$), defined as follows: let $\Omega=\{(x,\xi,y)\in \bar{B_r}\times M, \dist(x,y)<r\}$; let $\psi:\Omega\to \C$ be holomorphic in its first factor and such that $\psi(x,\xi,x)=\frac{i}{2}(r-\|\xi\|_{g(x)}^2)$. This function, when restricted to $(X_r\times M)\cap \Omega$, satisfies the following properties:
\begin{itemize}
\item $-C\dist(x,y)^2\leq \im(\psi)(x,\xi,y)\leq -c\dist(x,y)^2$ for some $0<c<C$.
\item $(\nabla_x,\nabla_\xi,\nabla_y)\psi(x,\xi,x)=(-\xi,-i\xi,\xi)$.
\end{itemize}
In particular, $X_r\times M\times (0,+\infty)\ni (x,\xi,y)\mapsto t\psi(x,\xi,y)$ is a positive phase function whose Lagrangian $\Lambda$ is the only Lagrangian which is CR-holomorphic in the first component and such that
\[
\Lambda_\R=\left\{(x,-t\omega,\omega,0,y,-t\omega),\;t\in \R,\;(x,\omega)\in S^*M\right\}.
\]
We can introduce our first guess for the FBI transform:
\[
T_0(x,\xi,y)=\int e^{t\psi(x,\xi,y)}t^{\frac{n+1}{4}}\dd t.
\]
As in the flat case, the power of $t$ is chosen such that $T_0$ is a bounded operator, with bounded inverse, between $L^2(M)$ and $\ker_{L^2(X_r)}(\bar{\partial}_b)$; in other terms, $T_0^*T_0$ is an elliptic pseudodifferential operator of degree~$0$.

$T_0$ may not exactly be unitary, but we can remedy to this situation in the following way: $T_0T_0^*$ has the same Lagrangian as $S$, and therefore it is of the form $I_{\psi}(t_0)$ for some self-adjoint amplitude $t_0$; there exists an inverse square root $r$ for $t_0$ (because the calculus of these Toeplitz operators is, locally, conjugated to that of pseudodifferential operators on flat space), and therefore, letting
\[
T=I_{\psi}(r)T_0,
\]
then $T$ is a unitary transformation between $L^2(M)$ and $\ker_{L^2(X_r)}(\bar{\partial}_b)$. Useful properties of $T$ are as follows:
\begin{itemize}
\item If $A$ is an analytic pseudodifferential operator on $M$, then $TAT^*$ is a Toeplitz operator with principal symbol $(x,\omega,t)\mapsto
\sigma(A)(x,-t\omega)$.

\item More generally, if $A$ is an analytic Fourier integral operator on $M$, with Lagrangian $\Lambda_A$, then $TAT^*$ is a Fourier integral operator whose Lagrangian $\Lambda_{TAT^*}$ is a complex Lagrangian, CR-holomorphic in the first factor, anti-CR-holomorphic in the second factor, and such that
\begin{align*}
(\Lambda_{TAT^*})_{\R}=\left\{
\begin{aligned}
&(x_1,-t_1\omega_1,\omega_1,0,x_2,-t_2\omega_2,\omega_2,0),\\
&(x_1,\omega_1,x_2,\omega_2)\in (\Lambda_A)_\R, \quad t_1>0, t_2>0, \\ &(x_1,\omega_1)\in S^*M, \quad \left(x_2,\frac{t_2}{t_1}\omega_2\right)\in S^*M
\end{aligned}
\right\}.
\end{align*}
Existence and uniqueness of this Lagrangian are guaranteed by Proposition~\ref{prop:completing_Lagrangians}. In particular, in this case one has $\Sigma_+\simeq T^*M\setminus \{0\}$.
\end{itemize}

We can use this transformation to describe quantum propagators as analytic Fourier integral operators.

\begin{prop}\label{prop:quantum_propag}
Let $P$ be an analytic self-adjoint degree~1 pseudodifferential operator on $M$. For every $s\in \R$, the fundamental solution $e^{-isP}$ of the Schr\"odinger equation for $P$ is an analytic Fourier integral Operator whose Lagrangian is the graph of the bicharacteristic flow $\phi$ of $P$ at time $s$. In particular, $Te^{-isP}T^*$ is a Fourier integral operator with a singular kernel of the form
\[
(x,y)\mapsto \int_0^{+\infty}e^{it\psi(s,x,y)}a(s,x,y;t)\dd t.
\]
Here, $\psi(s,x,y)$ is CR-holomorphic with respect to $x$, anti-CR-holomorphic with respect to $y$, and satisfies $\im(\psi(x,y))\asymp -\dist(y,\phi_s(x))$; moreover $a$ is the realisation of a formal analytic amplitude.
\end{prop}

\begin{proof}
We will prove the claimed description of $Te^{-isP}T^*$, from which the more abstract fact that $e^{-isP}=T^*(Te^{-isP}T^*)T$ is a Fourier integral operator will follow.

By Proposition~\ref{prop:repr_BG_smooth}, a pseudodifferential operator $Q$ on $X$ is such that $[Q,S]$ and $\mathit{SQS}-TPT^*$ are continuous from $L^2$ to $L^2$, if and only if the principal symbol of $Q$, near $\Sigma_+$ is the CR-holomorphic extension of the principal symbol of $P$ (seen as a real-analytic function on the totally real manifold $\Sigma_+$).

We know from the $C^{\infty}$ theory that $e^{-isQ}$ is, up to a smooth remainder, a Fourier integral operator whose Lagrangian is the graph of the bicharacteristic flow of the principal symbol of $Q$. Since the principal symbol Poisson-commutes with that of $(\bar{\partial}_b,\partial_b)$, the (smooth) Fourier integral operator $\Pi e^{-isQ}\Pi$ has exactly the requested Lagrangian. Letting $U_0(s)$ be any elliptic degree~0 analytic Fourier integral operator with the same Lagrangian (obtained by Proposition~\ref{prop:quantized-contact}, for instance), the principal (order 1) symbols of $\frac{\dd}{\dd s}U_0(s)$ and $iU_0(s)TPT^*$ coincide (alternatively, one can compute principal symbols to show this). The phase of this Fourier integral operator is the same as that of $U_0(s)$.

The operator $R(s)=U_0(s)e^{-isTPT^*}$ then solves the equation
\[
\frac{\dd}{\dd s}R(s)=\left[\frac{\dd}{\dd s}U_0(s)-iU_0(s)TPT^*\right]e^{-isTPT^*}.
\]
Introducing $V_0(s)$ an inverse for $U_0(s)$ (we already have an inverse if we applied Proposition~\ref{prop:quantized-contact}), modulo a real-analytic remainder,
\begin{equation}\label{eq:Rs}
\frac{\dd}{\dd s}R(s)=\left[\frac{\dd}{\dd s}U_0(s)-iU_0(s)TPT^*\right]V_0(s).
\end{equation}
The right-hand side is a degree~$0$ analytic Fourier integral operator with the same phase as $\Pi$; it is a degree~$0$ Toeplitz operator in the sense of Proposition~\ref{prop:Toeplitz_general}. Choosing a Banach norm of formal analytic amplitudes which contains the total symbol of $\left[\frac{\dd}{\dd s}U_0(s)-iU_0(s)TPT^*\right]V_0(s)$ and on which the formal product is continuous, one can give an approximate solution for~\eqref{eq:Rs} by the Picard--Lindel\"of theorem; thus there exists an analytic Toeplitz operator $\widetilde{R(s)}$ such that, modulo an analytic remainder,
\[
\frac{\dd}{\dd s}\widetilde{R(s)}=\left[\frac{\dd}{\dd s}U_0(s)-iU_0(s)TPT^*\right]V_0(s).
\]
To conclude, we can apply the Duhamel formula (the true propagator $e^{isTPT^*}$ preserves real-analytic functions) and we obtain that, modulo an analytic remainder,
\[
e^{isTPT^*}=\widetilde{R(s)}^{-1}V_0(s).
\]
This proves the second part of the claim, and conjugating back, we obtain that $T^*\widetilde{R(s)}^{-1}V_0(s)T$ is a Fourier integral operator which coincides with $e^{isP}$ up to an analytic remainder.
\end{proof}

\subsection{From microlocal to semiclassical analysis}\label{sec:from-micr-semicl}
In this section we broadly describe how one can obtain semiclassical results from the description above.

In what follows, we let $X$ be the boundary of a compact, strongly pseudoconvex open set and suppose that $X$ admits an $S^1$ action $(r_{\theta})_{\theta\in S^1}$ which preserves all the structure. Two important examples are:
\begin{enumerate}
\item\label{sect6.2case1} The case where $X$ is the Grauert tube around a compact Riemannian manifold of the form $S^1\times M$. The action of $S^1$ on $S^1\times M$ (by rotation of the first factor) is an isometry and lifts to a symplectic transformation on $T^*(S^1\times M)$, which is also an isometry for the K\"ahler structure near the zero section.
\item\label{sect6.2case2} The case where $X$ is the circle bundle of the dual of a positive line bundle over a K\"ahler manifold.
\end{enumerate}

Among Toeplitz operators, those whose formal amplitude is itself invariant under the $S^1$ action commute with the $S^1$ action, and thus one can decompose their action over the eigenmodes $\clH_k,k\in \Z$ of the $S^1$ action. Changing variables, the singular kernel of a Toeplitz operator $I_{\psi}(a)$ restricted to $\clH_k$ is then
\[
(x,y)\mapsto k\int_{S^1}\int_{0}^{+\infty}e^{ik(s\psi(r_{\theta}x,y)-\theta)}a(x,y,ks)\dd
\theta \dd s.
\]
To perform a stationary phase in the variables $(s,\theta)$, we need to satisfy the condition
\[
\alpha(V)\neq 0
\]
where $\alpha$ is the contact 1-form on $X$ of Proposition~\ref{prop:CR-structure} and $V$ is the tangent vector field of the $S^1$ action. In case~\eqref{sect6.2case2} above, this is always true because the Reeb flow coincides with the $S^1$ action (in fact one can choose $\psi$ with simple and explicit dependence on $\theta$, see for instance the unlabelled equation following (24) in~\cite{zelditch_pointwise_2018}). In case~\eqref{sect6.2case1}, the Reeb flow is the geodesic flow on $S^1\times M$, so we restrict our attention to the open set
\[
X_+=\left\{(y,\eta,x,\xi)\in S^*(S^1\times M),\; \eta>0\right\}.
\]
Thus, in both cases, we are in position to apply the stationary phase lemma in the variables $(s,\theta)$ (see for instance~\cite{zelditch_szego_2000}), and obtain a description modulo errors of size $e^{-ck}$ for some $c>0$ (indeed, elements of analytic function spaces have exponentially fast decaying Fourier modes). In case~\eqref{sect6.2case2} above, we obtain semiclassical Berezin--Toeplitz operators; in case 1, after a stereographic change of variables $X_+\to S^1\times T^*M$, we obtain semiclassical pseudodifferential operators.

With this point of view, one can for instance recover the expressions of the semiclassical Szeg\H{o} kernel, and many properties of semiclassical analytic pseudodifferential operators after conjugating back by a semiclassical Bargmann transform, since we know that the Toeplitz and pseudodifferential algebras are equivalent in the analytic semiclassical case~\cite{rouby_analytic_2018}.

This procedure can be generalised to other ``small parameter'' techniques, notably to study the action of (microlocal) analytic Fourier integral operators on WKB states. We hope to describe this in future work.


\appendix

\section{The algebra of symbols of pseudodifferential operators}\label{sec:algebra-symb-pseud}

\begin{prop}\label{prop:composition-formal}
Define the following formal analytic symbol norm on $\R^{d}_x\times \R^d_{\xi}$, for $m\in \R, \rho>0,R>0$:
\[
\|a\|_{S^{\rho,R}_m}=\sup_{k\in \N, \alpha\in \N^{2d}, (x,\xi)\in
\R^{2d}}\frac{|\nabla^{\alpha}a_k(x,\xi)|(1+k+|\alpha|)^m}{\rho^{|\alpha|}R^k(|\alpha|+k)!}.
\]
For the star-product of left-quantization on $T^*\R^d$, there exists $m_0(d)$ such that
\[
\|a\sharp b\|_{S^{\rho,R}_m}\leq 12\|a\|_{S^{\rho,R}_m}\|b\|_{S^{\frac{\rho}2,\frac R4}_m}
\]
if
\[
R\geq 2^{d+2}\rho^2, \qquad m\geq m_0(d).
\]
\end{prop}
\begin{proof}
We start with the formula
\[
(a\sharp b)_k=\sum_{n=0}^k\sum_{l=0}^{k-n}\sum_{|\beta|=n}\frac{1}{\beta!}\partial^\beta_xa_l\partial_{\xi}^{\beta}b_{k-n-l}.
\]
Since
\[
\sum_{|\beta|=n}\frac{n!}{\beta!}=n^d\leq 2^{nd}
\]
we obtain that, for every $\alpha$ with $|\alpha|=j$,
\begin{multline*}
\left|\nabla^{\alpha}(a\sharp b)_k\right|\leq
\|a\|_{S^{\rho,R}_m}\|b\|_{S^{\frac{\rho}2,\frac R4}_m}\frac{\rho^jR^k(j+k)!}{(j+k+1)^m} \\
\times \sum_{n=0}^k\left(\frac{2^{d+1}\rho}{R}\right)^n\,\sum_{l=0}^{n-k}\sum_{j_1=0}^j\sum_{\substack{|\gamma|=j_1\\
\gamma\leq
\alpha}}
\frac{(n+j_1+l)!(k-l+j-j_1)!j!}{4^{k-l}2^{j-j_1-n}n!(k+j)!j_1!(j-j_1)!}
\\ \times \left(\frac{1+j+k}{(1+n+j_1+l)(1+k-l+j-j_1)}\right)^m.
\end{multline*}
Let us prove that, if $0\leq j_1\leq j$ and $0\leq l+n\leq k$, then
\[
\frac{(n+j_1+l)!(k-l+j-j_1)!j!}{2^{k-l+j-j_1}n!(k+j)!j_1!(j-j_1)!}\leq 1.
\]
If the other parameters are fixed, then
\[
l\mapsto
\frac{(n+j_1+l)!(k-l+j-j_1)!j!}{2^{k-l+j-j_1}n!(k+j)!j_1!(j-j_1)!}
\]
is log-convex; there are two extremal points.



At \underline{$l=0$}, we obtain
\[
\frac{(n+j_1)!(k+j-j_1)!j!}{2^{k+j-j_1}n!(k+j)!j_1!(j-j_1)!}\leq
\frac{(n+j_1)!j!k!}{n!(k+j)!j_1!}.
\]
This increasing function of $j_1$ reaches a maximum at $j_1=j$, where we obtain
\[
\frac{(n+j)!k!}{n!(k+j)!}\leq 1
\]
since $n\leq k$.

In the other case \underline{$l=k-n$}, we obtain
\[
\frac{(k+j_1)!(n+j-j_1)!j!}{2^{n+j-j_1}n!(k+j)!j_1!(j-j_1)!}\leq
\frac{(k+j_1)!j!}{(k+j)!j_1!}\leq 1
\]
since $j_1\leq j$. It remains to bound, for fixed $l,n$ such that $l+n\leq k$,
\[
\sum_{j_1=0}^j\sum_{\substack{|\gamma|=j_1\\ \gamma\leq
\alpha}}\left(\frac{1+j+k}{(1+n+j_1+l)(1+k-l+j-j_1)}\right)^m
\]
To this end we use the fact that
\[
\#(|\gamma|=j_1,\gamma\leq \alpha)\leq \min(j_1,(j-j_1))^{d-1};
\]
thus, this sum is bounded by
\[
\sum_{j_1=0}^j\frac{(1+j+n+k)^m\min(j_1,j-j_1)^{d-1}}{(1+n+j_1+l)^m(1+k-l+j-j_1)^m}.
\]
Now we use~\cite[Lemma~2.13]{deleporte_toeplitz_2018}; for $m$ large enough (depending only on $d$), this sum is smaller than $3$. To conclude,
\[
\sum_{l=0}^{n-k}\frac{1}{2^{k-n-l}}\leq 2
\]
and, if $R\geq 2^{d+2}\rho^2$,
\[
\sum_{n=0}^k\left(\frac{2^{d+1}\rho^2}{R}\right)^n\leq 2.\qedhere
\]
\end{proof}

\section{Functional spaces in analytic regularity}\label{sec:funct-spac-analyt}

This section collects the basic properties of the spaces of analytic functions and their duals. Most technical facts claimed here are proved in~\cite{hormander_introduction_1973}; see also~\cite[Chapter~6]{treves_analytic_2022}.

\begin{defi}\label{def:space_hol_func}
Let $\Omega\subset \C^n$ (or any complex paracompact manifold) be an open set. We define $\clO(\Omega)$ as the space of holomorphic functions from $\Omega$ to $\C$, endowed with the topology of local uniform convergence.
\end{defi} The topology on $\clO'(\Omega)$ coincides with that of compactly supported Radon measures, of which it is a closed subspace. Elements of $\clO'(\Omega)$ will be called \emph{analytic functionals}.

One can generalise this definition into ``germs'' of holomorphic functions and their duals.

\begin{defi}
Let $E\subset \C^n$ (or any complex paracompact manifold) be a general set (not necessarily open). Define $\clO(E)$ as the union of the spaces $\clO(\Omega)$ for all $\Omega$ open and containing $E$ (endowed with the colimit topology).
\end{defi} If $E$ is an open subset or $\R^n$, or more generally an open subset of a real-analytic submanifold of $\R^n$, then $\clO(E)$ is the space of real-analytic functions on $E$. If $E$ is a compact subspace of $\R^n$ with non-empty interior, then $\clO(E)$ is the space of functions on $E$ which are real-analytic up to the boundary (and, therefore, which extend into real-analytic functions on some open neighbourhood of $E$).

Given two sets $E\subset F\subset \C^n$ one can naturally define a restriction map from $\clO(F)$ to $\clO(E)$. By duality, this defines a natural map $\clO'(E)\to \clO'(F)$. This map can only be injective when the restriction $\clO(F)\to \clO(E)$ has dense image; this is called the \emph{Runge property}. We will only be interested in the case where $E$ is a compact subset of $\R^n$, in which case the Runge property is always satisfied, by elementary approximation theory.

\begin{prop}
[{\cite[Theorem~6.2.14]{treves_analytic_2022}}]
\label{prop:real_cpts_are_Runge}
Let $K$ be a compact subset of~$\R^n$. Then entire functions are dense in $\clO(K)$; that is to say, $K$ is Runge.
\end{prop}

The proof roughly consists in showing that, with $\rho_N:x\mapsto=c_nN^{-n}\exp(-N|x|^2)$, given $f\in \clO(K)$ and $g$ any smooth extension of $f$ to $\R^n$, the sequence of entire functions $(\rho_N*g)_{N\in \N}$ converges towards $f$ in the topology of $\clO(K)$.

One of the main useful properties of Runge sets is the fact that one can solve the $\bar{\partial}$ problem on them.

\begin{prop}[{\cite[Theorem~2.7.8]{hormander_introduction_1973}}]
\label{prop:solve_dbar_on_Runge}
Let
$K\Subset \C^n$ be compact and Runge. Then there exists a basis of neighbourhoods $(\Omega_j)_{j\in \N}$ of $K$ in $\C^n$ such that the following is true. For every $j\in \N$ and every $f=(f_1,\dots,f_n)\in C^1(\Omega_j,\C^n)$ which is $\bar{\partial}$-closed, in the sense that
\[
\bar{\partial}_{k}f_\ell=\bar{\partial}_{\ell}f_k,\qquad
\forall 1\leq k,\ell\leq n,
\]
there exists $u\in C^1(\Omega_j,\C)$ such that $\bar{\partial}u=f$.
\end{prop}

The ability to prove the $\bar{\partial}$-problem on some neighbourhoods of real compact sets allow us to describe how the spaces $\clO$ and $\clO'$ behave under natural operations on their domains.

\begin{prop}\label{prop:inclusion_union_O_spaces}
Let $K_1,K_2$ be compact subsets of $\R^n$.
\begin{enumerate}
\item\label{propB.5.1} $\clO(K_1)\cap \clO(K_2)=\clO(K_1\cup K_2)$.
\item\label{propB.5.2} $\clO(K_1)+ \clO(K_2)=\clO(K_1\cap K_2)$.
\item\label{propB.5.3} $\clO(K_1)'\cap \clO'(K_2)=\clO'(K_1\cap K_2)$.
\item\label{propB.5.4} $\clO(K_1)'+ \clO'(K_2)=\clO'(K_1\cup K_2)$.
\end{enumerate}
\end{prop}

\begin{proofc}
\begin{itemize}
\item[\eqref{propB.5.1}] Let $U_1,U_2$ be respective neighbourhoods of $K_1$ and $K_2$ in $\C^n$. Then holomorphic functions on $U_1\cup U_2$ are exactly functions that are both holomorphic on $U_1$ and on $U_2$, because holomorphy is a local property.
\item[\eqref{propB.5.2}] The inclusion from left to right is obvious and it remains to prove, given $f\in \clO(K_1\cap K_2)$, that it is of the form $f_1+f_2$, where $f_1\in \clO(K_1)$ and $f_2\in
\clO(K_2)$.

Let $\Omega_1$ and $\Omega_2$ be small neighbourhoods in $\C^n$ of respectively $K_1$ and $K_2$, so that $f$ extends to $\widetilde{f}\in \clO(\Omega_1\cap \Omega_2)$ (the function $\widetilde{f}$ may or may not correspond to the usual notion of holomorphic extension, depending on whether $K$ is the closure of its interior in $\R^n$).

The compact sets $K_1\setminus \Omega_2$ and $K_2\setminus \Omega_1$ do not intersect and therefore lie at positive distance from each other. We let $\chi\in C^{\infty}(\Omega_1\cup \Omega_2,\R)$ be such that $\chi=1$ on a neighbourhood $U_1$ of $K_1\setminus
\Omega_2$ and $\chi=0$ on a neighbourhood $U_2$ of $K_2\setminus
\Omega_1$. The one-form $\alpha=\bar{\partial}(\widetilde{f}\chi)$, well-defined on $\Omega_1\cap \Omega_2$, is equal to $0$ on $(U_1\cup U_2)\cap (\Omega_1\cap \Omega_2)$, and therefore can be smoothly extended by $0$ on $\Omega=U_1\cup U_2\cup (\Omega_1\cap \Omega_2)$.

$\Omega$ is a neighbourhood of $K_1\cup K_2$, which is a Runge set by Proposition~\ref{prop:real_cpts_are_Runge}; by Proposition~\ref{prop:solve_dbar_on_Runge}, this means that there exists a smaller neighbourhood $\Omega'\subset \Omega$ of $K_1\cup K_2$ and $u\in C^{\infty}(\Omega')$ such that $\bar{\partial}u=\alpha$.

To conclude, the function $h_1=(1-\chi)\widetilde{f}+u$ is well-defined on $(U_1\cup (\Omega_1\cap \Omega_2))\cap \Omega'$, which is an open neighbourhood of $K_1$; it satisfies $\bar{\partial}h_1=\alpha-\bar{\partial}u=0$; therefore $h_1\in \clO(K_1)$. In the same way, $h_2=\chi\widetilde{f}-u$ is well-defined and holomorphic on $(U_1\cup (\Omega_1\cap \Omega_2))\cap \Omega'$ which is an open neighbourhood of $K_2$; moreover, on $K_1\cap K_2$, one has $h_1+h_2=f$; this concludes this part of the proof.
\item[\eqref{propB.5.3}] This statement is the dual of~\eqref{propB.5.2}.
\item[\eqref{propB.5.4}] This statement is the dual of~\eqref{propB.5.1}.\qedhere
\end{itemize}
\end{proofc}

Given two compacts $K_1\subset K_2$ of $\R^n$, there is no natural restriction map from $\clO'(K_2)$ to $\clO'(K_1)$ (in fact, we have a natural injective \emph{extension} map from $\clO'(K_1)$ to $\clO'(K_2)$). Nevertheless, there is a well-defined notion of \emph{support} of an analytic functional, thanks to the injectivity of this extension map and Proposition~\ref{prop:inclusion_union_O_spaces}$\MK$\eqref{propB.5.3}.

Let us remark that if we define $C^{\infty}(K)=\clE(K)$, for $K$ any compact of $\R^n$, following Whitney~\cite{whitney_analytic_1934}, then the equivalents of all items of Proposition~\ref{prop:inclusion_union_O_spaces} are also true, with the small caveat that the dense maps $\clE(\R^n)\to
\clE(K)$ fail to be injective. We will use the following restriction theorem on locally compact distributions: they can be restricted to a smaller subset modulo a boundary indeterminacy.

\begin{prop}\label{prop:presheaf}
Let $K_1\subset K$ be two compact sets. The inclusion map $\clE'(K_1)\to \clE'(K)$ has a left inverse modulo $\clE'(\partial K_1)$. This right inverse does not depend on adding an element of $\clE'(\partial K)$, and therefore defines a well-defined map
\[
\rho_{K,K_1}:\clE'(K)/\clE'(\partial K)\to
\clE'(K_1)/\clE'(\partial K_1).
\]
\end{prop}
\begin{proof}
Let $K_2=\bar{K\setminus K_1}$. Recall that $\clE'(K)=\clE'(K_1)+\clE'(K_2)$, and given $u\in \clE'(K)$, choose $u_1\in \clE'(K_1)$ such that $u-u_1\in \clE'(K_2)$. The distribution $u_1$ is not unique, but given any other choice $v_1$, the difference $u_1-v_1\in
\clE'(K_1)$ is such that $(u-u_1)-(u-v_1)=v_1-u_1\in
\clE'(K_2)$. Therefore $u_1-v_1\in \clE'(K_2\cap K_1)=\clE'(\partial K_1)$, and the class of $u_1$ modulo $\clE'(\partial K_1)$ is uniquely defined. If one had $u \in \clE'(K_1)$ to begin with, then one can choose $u_1=u$, so that this map is indeed a left inverse to the extension map.

If $u\in \clE'(\partial K)$, then, since $\partial K\subset (\partial K_1)\cup K_2$, one can choose $u_1\in
\clE'(\partial K_1)$ in the lines above, and therefore $u$ is mapped to $0$; this concludes the proof.
\end{proof}

The same property holds when $\clE'$ is replaced with $\clO'$. The two families of spaces $\clE'(K)/\clE'(\partial K)$ and $\clO'(K)/\clO'(\partial K)$ are quite practical because of this well-defined restriction map, which mimics the ability to restrict analytic or smooth functions to a smaller set. Elements of the spaces $\clO'(K)/\clO'(\partial K)$ are called hyperfunctions. Proposition~\ref{prop:presheaf} and its equivalent for $\clO'$ means that these families of spaces (indexed by $\mathring{K}$) form \emph{pre-sheaves}.

\begin{defi}\label{def:support_analytic_functional}
Let $K\Subset \R^n$ and let $f\in \clO'(K)$. The support of $f$ is the smallest compact $K_1\subset K$ such that $f\in \clO'(K_1)$.
\end{defi}

This notion is well-defined: first of all ``$f\in \clO'(K_1)$'' makes sense because the extension $\clO'(K_1)\to
\clO'(K)$ is injective; second, if two compacts $K_1$ and $K_2$ are such that $f\in \clO'(K_1)\cap \clO'(K_2)$, then $f\in \clO'(K_1\cap K_2)$ by Proposition~\ref{prop:inclusion_union_O_spaces}$\MK$\eqref{propB.5.3}.

If $U$ is a relatively compact open set in $\R^n$ then $\clO(\bar{U})\subset C^{\infty}(\bar{U})$ and therefore $\clE'(\bar{U})\subset
\clO'(\bar{U})$. In this case, the support of a distribution $T\in \clE'(\bar{U})$ coincides with its support as an element of $\clO'(\bar{U})$.

We now reach the main result of this appendix: one can patch together analytic functionals defined on different compact sets which agree on the intersection (meaning that the support of their difference lies away from the intersection), and if they only agree on the intersection modulo a real-analytic function then we can patch them together modulo a real-analytic function.

\begin{prop}\label{prop:patching_analytic_func_or_distros}
Let $K_1,K_2$ be compact subsets of $\R^n$. Let $u_1\in
\clO'(K_1)$ and $u_2\in \clO'(K_2)$ be such that $\supp(u_1-u_2)\cap K_1\cap K_2=\emptyset$. Then there exists $u\in
\clO'(K_1\cup K_2)$ such that $\supp(u-u_1)\cap K_1=\emptyset$ and $\supp(u-u_2)\cap K_2=\emptyset$. If $u_1,u_2$ belong to $\clE'$, then so does $u$.

More generally, if there exists $f\in \clO(K_1\cap K_2)$ such that $\supp(u_1-u_2-f)\cap K_1\cap K_2=\emptyset$, then there exists $u\in \clO'(K_1\cup K_2)$, $g_1\in \clO(K_1)$ and $g_2\in \clO(K_2)$ such that $\supp(u-u_1-g_1)\cap K_1=\emptyset$ and $\supp(u-u_2-g_2)\cap K_2=\emptyset$. Again if $u_1,u_2$ belong to $\clE'$, then so does $u$.
\end{prop}
\begin{proof}
By hypothesis, there exists two compacts $L_1\in K_1\setminus K_2$ and $L_2\in K_2\setminus K_1$ such that $u_1-u_2\in \clO'(L_1\cup L_2)$. Using Proposition~\ref{prop:inclusion_union_O_spaces}, there exists $f_1\in \clO'(L_1)$ and $f_2\in \clO'(L_2)$ such that $u_1-u_2=f_1+f_2$.

Then $u=u_1-f_2=u_2+f_1$ is an element of $\clO'(K_1\cup K_2)$ such that $u-u_1=f_2\in \clO'(L_2)$ and $u-u_2=f_1\in
\clO'(L_1)$. This concludes this part of the proof, and we can seamlessly replace $\clO'$ with $\clE'$ in the lines above.

If now $u_1-u_2\in \clO'(L_1\cup L_2)+\clO(K_1\cup K_2)$, then by writing $\clO(K_1\cup K_2)=\clO(K_1)+\clO(K_2)$ we can correct $u_1$ and $u_2$ by respective elements of $\clO(K_1)$ to reduce ourselves to the previous case.
\end{proof}



To conclude, we mention the generalisation of these results to general paracompact real-analytic manifolds. Any paracompact real-analytic manifold is an analytic submanifold of $\R^n$ for $n$ large enough~\cite{grauert_levis_1958}, so that we only need to restrict our attention to compact sets which belong to this submanifold; thus there is no additional difficulty. Of course, when considering the restriction map as in Proposition~\ref{prop:presheaf}, the boundary is then taken with respect to the topology of the submanifold, not with respect to the ambient $\R^n$ topology.



\bibliography{20260219_deleporte}
\end{document}
